\documentclass[11pt,a4paper]{article}
\usepackage{amsmath,amsxtra,amsthm,amssymb,mathrsfs,bbm}
\usepackage{hyperref}
\usepackage[margin=2cm]{geometry}
\usepackage{graphicx}
\usepackage[pdftex]{color}
\usepackage[normalem]{ulem}
\usepackage[shortlabels]{enumitem}
%\usepackage[small,nohug,noPostScript]{diagrams}
%\usepackage{stmaryrd}
\usepackage{makeidx}
\usepackage[all]{xy}
\usepackage{float}
\usepackage{tocloft}
\usepackage{tikz}
%\usetikzlibrary{positioning}
\usetikzlibrary{shapes,arrows}
\usetikzlibrary{patterns}
\usetikzlibrary{decorations.shapes}
\usetikzlibrary{decorations.pathreplacing}
\usetikzlibrary{decorations.pathmorphing}
\numberwithin{equation}{section}

\setlength{\parindent}{0mm}
\setlength{\parskip}{11pt}

%set default enumerate environment
\setlist[enumerate]{label={\rm(\arabic*)},itemsep=3pt,topsep=-6pt}
\setlist[itemize]{itemsep=3pt,topsep=-6pt}

\newenvironment{claim}[1][]
{\begin{list}{}{\setlength{\itemsep}{0pt}
	\setlength{\listparindent}{0pt}
	\setlength{\labelsep}{0.2cm}
	\setlength{\topsep}{2pt}
	\setlength{\leftmargin}{0.5cm}}

	\item[] \textbf{Claim:$\mbox{}\quad$}#1
	\item[] \textbf{Proof of Claim: }\\
}
{\end{list}}

%%% set up proof environment
\makeatletter
\renewenvironment{proof}[1][\proofname]{\par
    \pushQED{\qed}%
    \normalfont \topsep0\p@\@plus6\p@ \labelsep1em\relax
    \trivlist
    \item[\hskip\labelsep\indent
        \bfseries #1]\ignorespaces
    \mbox{}
}{
    \popQED\endtrivlist\@endpefalse
}
\makeatother

%%% Remove section numbering but retain section counting
%\makeatletter
%\def\@seccntformat#1{%
%  \expandafter\ifx\csname c@#1\endcsname\c@section\else
%  \csname the#1\endcsname\quad
%  \fi}
%\makeatother

%add dots to sections in TOC
\renewcommand{\cftsecleader}{\cftdotfill{\cftdotsep}}
%remove section number in TOC
\makeatletter
\renewcommand{\cftsecpresnum}{\begin{lrbox}{\@tempboxa}}
\renewcommand{\cftsecaftersnum}{\end{lrbox}}
\makeatother

%abovespace,belowspace,bodyfont,indent,headfont,headpunct(e.g.\newline),headspace,custom-head-spec
\newtheoremstyle{all}{10pt}{0pt}{\itshape}{0pt}
{\bfseries}{.}{5pt plus 1pt minus 1pt}{\thmname{#1}\thmnumber{ #2}\thmnote{ (#3)}}
\theoremstyle{all}

\newtheorem{theorem}{Theorem}[section]
\newtheorem{proposition}[theorem]{Proposition}
\newtheorem{lemma}[theorem]{Lemma}
\newtheorem{corollary}[theorem]{Corollary}
\newtheorem{example}[theorem]{Example}
\newtheorem*{example*}{Example}
\newtheorem{definition}[theorem]{Definition}
\newtheorem{exercise}[theorem]{Exercise}
\newtheorem*{exercise*}{Exercise}
\newtheorem{assumption}[theorem]{Assumption}
\newtheorem*{convention}{Convention}
\newtheorem*{notation}{Notation}

\theoremstyle{remark}
\newtheorem{remark}[theorem]{Remark}
\newtheorem*{remark*}{Remark}

\DeclareMathOperator{\isom}{\cong}
\DeclareMathOperator{\ifonlyif}{\Leftrightarrow}
\DeclareMathOperator{\imply}{\Rightarrow}
\DeclareMathOperator{\onto}{\twoheadrightarrow}
\DeclareMathOperator{\into}{\hookrightarrow}
\newcommand{\ses}[3]{0\rightarrow #1 \rightarrow #2 \rightarrow #3 \rightarrow 0}

\renewcommand{\P}{\mathbb{P}}
\newcommand{\Z}{\mathbb{Z}}
\newcommand{\R}{\mathbb{R}}
\newcommand{\C}{\mathbb{C}}
\newcommand{\bbA}{\mathbb{A}}
\newcommand{\CB}{\mathcal{B}}
\newcommand{\CC}{\mathcal{C}}
\newcommand{\CD}{\mathcal{D}}
\newcommand{\rmC}{\mathrm{C}}
\newcommand{\calX}{\mathcal{X}}

\newcommand{\id}{\operatorname{id}}
\newcommand{\soc}{\operatorname{soc}\nolimits}
\newcommand{\rad}{\operatorname{rad}\nolimits}
\renewcommand{\top}{\operatorname{top}\nolimits}
\newcommand{\Ext}{\operatorname{Ext}\nolimits}
\newcommand{\Tor}{\operatorname{Tor}\nolimits}
\newcommand{\Hom}{\operatorname{Hom}\nolimits}
\newcommand{\End}{\operatorname{End}\nolimits}
\newcommand{\Aut}{\operatorname{Aut}\nolimits}
\newcommand{\op}{\mathrm{op}}
\renewcommand{\Im}{\operatorname{Im}\nolimits}
\newcommand{\Ker}{\operatorname{Ker}\nolimits}
\newcommand{\Cok}{\operatorname{Cok}\nolimits}

\DeclareMathOperator{\moduleCategory}{\mathsf{mod}} \renewcommand{\mod}{\moduleCategory}
\DeclareMathOperator{\Mod}{\mathsf{Mod}}
\DeclareMathOperator{\proj}{\mathsf{proj}}
\DeclareMathOperator{\inj}{\mathsf{inj}}
\DeclareMathOperator{\ind}{\mathsf{ind}}

\DeclareMathOperator{\Char}{char}
\newcommand{\triv}{\operatorname{triv}\nolimits}
\newcommand{\sgn}{\operatorname{sgn}\nolimits}
\newcommand{\Mat}{\operatorname{Mat}\nolimits}
\newcommand{\GL}{\operatorname{GL}\nolimits}
\DeclareMathOperator{\Sn}{\mathfrak{S}}
\newcommand{\Res}{\operatorname{Res}\nolimits}
\newcommand{\Ind}{\operatorname{Ind}\nolimits}
\newcommand{\Symm}{\operatorname{Sym}\nolimits}
\newcommand{\Stab}{\operatorname{Stab}\nolimits}
\newcommand{\Tr}{\operatorname{Tr}\nolimits}
\newcommand{\bfm}{\mathbf{m}}
\newcommand{\bfT}{\mathbf{T}}

\newcommand{\Ann}{\operatorname{Ann}\nolimits}
\newcommand{\linA}{\vec{\mathbb{A}}}
\newcommand{\pdim}{\operatorname{pdim}}
\newcommand{\idim}{\operatorname{idim}}
\newcommand{\gldim}{\operatorname{gldim}\nolimits}
\newcommand{\Rej}{\operatorname{Rej}\nolimits}
\newcommand{\std}{\Delta}
\newcommand{\costd}{\nabla}
\newcommand{\filt}{\mathcal{F}}
\newcommand{\add}{\operatorname{add}}
\newcommand{\dimvec}{\mathbf{dim}}
\newcommand{\nunlhd}{\ntrianglelefteq}
\newcommand{\nlhd}{\ntriangleleft}
\newcommand{\nunrhd}{\ntrianglerighteq}
\newcommand{\nrhd}{\ntriangleright}
\newcommand{\LL}{\mathrm{LL}}

\newcommand{\Ab}{\operatorname{\mathcal{A}b}}

\newcommand{\comm}[1]{{\color[rgb]{0.1,0.9,0.1}#1}}
\newcommand{\old}[1]{{\color{red}#1}}
\newcommand{\new}[1]{{\color{blue}#1}}
\definecolor{darkblue}{rgb}{0,0,0.7}
\newcommand{\defn}[1]{\textsl{\color{darkblue} #1}}

\makeindex

\begin{document}
\begin{center}
\textsc{{\LARGE Topics in Mathematical Science VIII}\\
Autumn 2023}\\
\textsc{\Large Introduction to quiver representations and homological algebras}

\textsc{Aaron Chan}
\end{center}

Last update: \today

%\tableofcontents

%{\bf Prerequisites.}  One should be familiar with basic group theory, ring theory, and module theory; the recommended book by Erdmann-Holm will provide some help on the module theory in case some reminder is needed.  Linear algebra is absolutely essential.

\section*{Convention}
 
Throughout the course, $\Bbbk$ will always be a field.
All rings are unital and associative.
We only really work with artinian rings (but sometimes noetherian is also OK).
We always compose maps from right to left.

\section{Reminder on some basics of rings and modules}

\begin{definition}
Let $R$ be a ring.
A \defn{right $R$-module}\index{module} $M$ is an abelian group $(M,+)$ equipped with a (linear) \defn{$R$-action on the right of $M$} $\cdot:M\times R\to M$, meaning that for all $r,s\in R$ and $m,n\in M$, we have
\begin{itemize}
\item $m\cdot 1 = m$,
\item $(m+n)\cdot r=m\cdot r+n\cdot r$,
\item $m\cdot (r+s)=m\cdot r+m\cdot s$,
\item $m(sr)=(ms)r$.
\end{itemize}
Dually, a \defn{left $R$-module} is one where $R$ acts on the left of $M$ (details of definition left as exercise).
Sometimes, for clarity, we write $M_A$ for right $A$-module and ${}_AM$ for left $A$-module.
\end{definition}
Note that, for a commutative ring, the class of left modules coincides with that of right modules.

\begin{example}
$R$ is naturally a left, and a right, $R$-module.  Both are \defn{free}\index{free} $R$-module of rank 1.  Sometimes this is also called regular modules but it clashes with terminology used in quiver representation and so we will avoid it.

In general, a free $R$-module $F$ is one where there is a basis $\{x_i\}_{i\in I}$ such that for all $x\in F$, $x=\sum_{i\in I} x_ir_i$ with $r_i\in R$.  We only really work with free modules of finite rank, i.e. when the indexing set $I$ is finite.  In such a case, we write $R^n$.
\end{example}

\begin{convention}
All modules are right modules unless otherwise specified.
\end{convention}

\begin{definition}
Suppose $R$ is a \emph{commutative} ring.
A ring $A$ is called an \defn{$R$-algebra}\index{algebra} if there is a (unital) ring homomorphism $\theta:R\to A$ with image $\theta(R)$ being in the \defn{center}\index{center} $Z(A):=\{z\in A\mid za=az\;\forall a\in A\}$ of $A$.
In such a case, $A$ is an $R$-module and so we simply write $ar$ for $a\in A,r\in R$ instead of $a\theta(r)$.

An (unital) $R$-\defn{algebra homomorphism}\index{homomorphism!algebra} $f:A\to A'$ is a (unital) ring homomorphism $f$ that \defn{intertwines} $R$-action, i.e. $f(ar)=f(a)r$.

The \defn{dimension} of a $\Bbbk$-algebra $A$ is the dimension of $A$ as a $\Bbbk$-vector space; we say that $A$ is \defn{finite-dimensional} if $\dim_{\Bbbk} A<\infty$.
\end{definition}

Note that commutative ring theorists usually use dimension to mean Krull dimension, which has a completely different meaning.

\begin{example}
Every ring is a $\Z$-algebra.

The matrix ring $M_n(R)$ given by $n$-by-$n$ matrices with entries in $R$ is an $R$-algebra.
\end{example}

We will only really work with $\Bbbk$-algebras, where $\Bbbk$ is a field.
Most of the time, we will also assume $\Bbbk$ is algebraically closed for simplicity.
But it worth reminding there are many interesting $R$-algebras for different $R$, such as group algebra.
Recall that the \defn{characteristic}\index{characteristic} of $R$, denoted by $\Char R$, is $0$ if the additive order of the identity $1$ is infinite, or else the additive order itself.

\begin{example}
Let $G$ be a finite (semi)group and $R$ a commutative ring.
Let $A:=R[G]$ be the free $R$-module with basis $G$, i.e. every $a\in A$ can be written as the formal $R$-linear combination $\sum_{g\in G} \lambda_g g$ with $\lambda_g\in R$.
Then group multiplication extends ($R$-linearly) to a ring multiplication on $R[G]$, making $A$ an $R$-algebra.
\end{example}

\begin{example}
Recall that the \defn{direct product} of two rings $A,B$ is the ring $A\times B=\{(a,b)\mid a\in A, b\in B\}$ with unit $1_{A\times B}=(1_A,1_B)$.
It is straightforward to check that if $A,B$ are $R$-algebras, then $A\times B$ is also an $R$-algebra.
\end{example}

\begin{example}
Suppose that $A$ is a $\Bbbk$-algebra and $B$ is a $\Bbbk$-subspace of $A$ containing $1_A$ and closed under multiplication.  Then $B$ is also a $\Bbbk$-algebra.  We call such a $B$ a \defn{subalgebra} of $A$.  For a concrete example, the space of diagonal matrices forms a subalgebra of $M_n(\Bbbk)$.
\end{example}

\begin{definition}
A map $f:M\to N$ between right $R$-modules $M,N$ is a \defn{homomorphism}\index{homomorphism} if it is a homomorphism of abelian groups (i.e. $f(m+n)=f(m)+f(n)$ for all $m,n\in M$) that intertwines $R$-action (i.e. $f(mr)=f(m)r$ for all $m\in  M$ and $r\in R$).
Denote by $\Hom_R(M,N)$ the set of all $R$-module homomorphisms from $M$ to $N$.
We also write $\End_R(M):=\Hom_R(M,M)$.
\end{definition}

\begin{lemma}
$\Hom_R(M,N)$ is an abelian group with $(f+g)(m)=f(m)+g(m)$ for all $f,g\in\Hom_R(M,N)$ and all $m\in M$.
If $R$ is commutative, then $\Hom_R(M,N)$ is an $R$-module, namely, for a homomorphism $f:M\to N$ and $r\in R$, the homomorphism $fr$ is given by $m\mapsto f(mr)$.
\end{lemma}
%\begin{proof}
%Abelian group: exercise.
%
%Identity action: $(f1)(m)=f(m1)=f(m)$
%
%Action distribute over addition: $((f+g)r)(m)=(f+g)(mr)=f(mr)+g(mr)=f(m)r+g(m)r=(fr)(m)+(gr)(m)=(fr+gr)(m)$.
%
%Addition distributive over action:
%$(f(r+s))(m)=f(m(r+s))=f(mr+ms)=f(mr)+f(ms)=(fr)(m)+(fs)(m)=(fr+fs)(m)$.
%
%Associativity: $(f(sr))(m)=(f(rs))(m)=f(m(rs))=f((mr)s)=(fs)(mr)=((fs)r)(m)$.
%\end{proof}

\begin{definition}
$\End_R(M)$ is an associative ring where multiplication is given by composition and identity element being $\id_M$.  We call this the \defn{endomorphism ring}\index{endomorphism ring} of $M$.
\end{definition}

\begin{lemma}
If $A$ is an $R$-algebra over a commutative ring $R$, then any right $A$-module is also an $R$-module, and $\Hom_A(M,N)$ is also an $R$-module (hence, $\End_R(M)$ is an $R$-algebra).
\end{lemma}
%\begin{proof}
%The $A$-action map extends to an $R$-action map by precomposing with the natural map $R\times M\to Z(A)\times M\to  A\times M$ and so yields an $R$-module structure on any $A$-module $M$.
%By definition, since any $A$-module homomorphism is also $R$-module homomorphism, $\Hom_A(M,N)\subset\Hom_R(M,N)$ is an $R$-submodule.
%\end{proof}

\begin{example}
$A\cong \End_A(A)$ given by $a\mapsto (1_A\mapsto a)$ is an isomorphism of rings (or of $R$-algebras if $A$ is an $R$-algebra).  Note that if we work with left modules, then $A\cong \End_{A}({}_AA)^\op$, where $(-)^\op$ denotes the \defn{opposite ring} given by the same underlying set with reverse direction of multiplication, i.e. $a\cdot_{\op}b:=b\cdot a$.
\end{example}

Recall that an $R$-module $M$ is \defn{finitely generated}\index{finitely generated} if there exists as surjective homomorphism $R^n\onto M$, or equivalently, there is a finite set $X\subset M$ such that for any $m\in M$, we have $m=\sum_{x\in X}xr_x$ for some $r_x\in R$.

\begin{notation}
We write $\mod A$ for the collection of all finitely generated right $A$-modules.
\end{notation}

\pagebreak 
%%%%%%%%%%%%%%%%%%%%%%%%%%%%%%%%%%%%%%%%%%%%%%%%%%%%%%%%%%%%%%%%%%%%%%%%%%%%%%%
%%%%%%%%%%%%%%%%%%%%%%%%%%%%%%%%%%%%%%%%%%%%%%%%%%%%%%%%%%%%%%%%%%%%%%%%%%%%%%%
%%%%%%%%%%%%%%%%%%%%%%%%%%%%%%%%%%%%%%%%%%%%%%%%%%%%%%%%%%%%%%%%%%%%%%%%%%%%%%%
%%%%%%%%%%%%%%%%%%%%%%%%%%%%%%%%%%%%%%%%%%%%%%%%%%%%%%%%%%%%%%%%%%%%%%%%%%%%%%%

\section{Indecomposable modules and Krull-Schmidt property}


We recall two types of building blocks of modules.  The first one is indecomposability.

\begin{definition}
Let $M$ be a $R$-module and $N_1,\ldots, N_r$ be submodules.
We say that $M$ is the \defn{direct sum}\index{direct sum} $N_1\oplus \cdots \oplus N_r$ of the $N_i$'s if $M=N_1+\cdots +N_r$ and $N_j\cap (N_1+\cdots +N_{\hat{j}}+\cdots N_r)=0$.  Equivalently, every $m\in M$ can be written \uline{uniquely} as $n_1+n_2+\cdots +n_r$ with $n_i\in N_i$ for all $i$.
In such a case, we write $M\cong N_1\oplus \cdots \oplus N_r$. 
Each $N_i$ is called a \defn{direct summand}\index{direct summand} of $M$.

$M$ is called \defn{indecomposable}\index{indecomposable} if $M\cong N_1\oplus N_2$ implies $N_1=0$ or $N_2=0$.

We say that $M=\bigoplus_{i=1}^m M_i$ is an \defn{indecomposable decomposition}\index{decomposition} (or just decomposition for short if context is clear) of $M$ if each $M_i$ is indecomposable.  
%Such a decomposition is said to be unique if for any other decomposition $M=\bigoplus_{j=1}^n N_j$, we have $n=m$ and the $N_j$'s are permutation of the $M_i$'s.
\end{definition}

\begin{convention}
We write $(n_1,\ldots, n_r)$ instead of $n_1+\cdots+n_r$ with $n_i \in N_i$ for a direct sum $N_1\oplus \cdots \oplus N_r$.
\end{convention}

We will only work with direct sum with finitely many indecomposable direct summands.

%\begin{definition}
%Let $M$ be an $R$-module.
%\begin{enumerate}
%\item We say that $M=\bigoplus_{i=1}^m M_i$ is an \defn{indecomposable decomposition}\index{decomposition} (or just decomposition for short if context is clear) of $M$ if each $M_i$ is indecomposable.  Such a decomposition is said to be unique if for any other decomposition $M=\bigoplus_{j=1}^n N_j$, we have $n=m$ and the $N_j$'s are permutation of the $M_i$'s.
%\item We say that $A$ satisfies the \defn{Krull-Schmidt property}\index{Krull-Schmidt} if every finitely generated $A$-module $M$ admits a unique indecomposable decomposition.
%\end{enumerate}
%\end{definition}

\begin{example}
Suppose that $R_R$ is indecomposable as an $R$-module.
If $F$ is a free $R$-module of rank $n$, then $R^{\oplus n}:=R\oplus R\oplus \cdots \oplus R$ (with $n$ copies of $R$) is a decomposition of $F$.
\end{example}

\begin{example}
Consider the matrix ring $A:=\Mat_n(\Bbbk)$ over a field $\Bbbk$.  Let $V$ be the `row space', i.e. $V=\{(v_j)_{1\leq j\leq n}\mid v_j\in \Bbbk\}$ where $X\in \Mat_n(\Bbbk)$ acts on $v\in V$ by $v\mapsto vX$ (matrix multiplication from the right).  
Since for any pair $u,v\in V$, there always exist $X$ so that $v=uX$, we see that there is no other $A$-submodule of $V$ other than $0$ or $V$ itself.
Hence, $V$ is an indecomposable $A$-module.
In particular, the $n$ different ways of embedding a row into an $n$-by-$n$-matrix yields an $A$-module isomorphism between $V^{\oplus n}\cong A_A$, which is the decomposition of the free $A$-module $A_A$.
\end{example}

The above example shows indecomposability by showing that $V$ is a \emph{simple} $A$-module, which is a stronger condition that we will come back later.
Let us give an example of a different type of indecomposable (but non-simple) modules.

\begin{example}
Let $A=\Bbbk[x]/(x^k)$ the \defn{truncated polynomial ring} for some $k\geq 2$.
This is an algebra generated by ($1_A$ and) $x$, and an $A$-module is just a $\Bbbk$-vector space $V$ equipped with a linear transformation $\rho_x\in \End_{\Bbbk}(V)$ (representing the action of $x$) such that $\rho_x^k=0$.

Consider a 2-dimensional space $V=\Bbbk \{v_1,v_2\}$ and a linear transformation
\[\rho_x = \begin{pmatrix}0 & 0 \\ 1 & 0\end{pmatrix}.\]
By definition $(av_1+bv_2)x = (a+b)v_2$, and so any submodules must contains $\Bbbk v_2$, i.e. $v_2$ spans a unique non-zero submodules.   
If, on the contrary, $V$ is not indecomposable, then we have $V=U_1\oplus U_2$ for (at least) two non-zero submodules $U_1,U_2$.  But $v_2$ must be contained in any submodule of $V$, hence, we have $v_2 \in U_1\cap U_2$, i.e. $U_1\cap U_2\neq 0$ -- a contradiction not decomposability.
\end{example}


\begin{proposition}\label{prop:hom commute direct sum}
There is a canonical $R$-module isomorphism
\[
\xymatrix@C=40pt@R=0pt{\Hom_A(\bigoplus_{j=1}^m M_j, \bigoplus_{i=1}^n N_i) \ar[r]^(.55){\cong}& \bigoplus_{i,j}\Hom_A(M_j,N_i)\\
f \ar@{|->}[r]& (\pi_i f\iota_j)_{i,j}
}
\]
where $\iota_j:N_j\to \bigoplus_j N_j$ is the canonical inclusion for all $j$ and $\pi_i:\bigoplus_iM_i \to M_i$ is the canonical projection for all $i$.
\end{proposition}

One can think of the right-hand space above as the space of $m$-by-$n$ matrix with entries in each corresponding Hom-space.

Recall that an \defn{idempotent} $e\in R$ is an element with $e^2=e$.
For example, the identity map $\id_M\in \End_A(M)$ (the unit element of the endomorphism ring) is an idempotent. 
From the previous proposition, we see that for a decomposition $M=N_1\oplus N_2$, we have idempotents
\[
e_i : M\xrightarrow{\pi_i} N_i \xrightarrow{\iota_i} M
\]
for both $i=1,2$.
Hence, being decomposable implies existence of multiple idempotents; this turns out characterise indecomposability completely.

\begin{proposition}\label{prop:indec characterisations}
Let $A$ be a finite-dimensional algebra and $M$ be a finite-dimensional non-zero $A$-module.
Then the following hold.
\begin{enumerate}
\item (Fitting's lemma) For any $f\in \End_A(M)$, there exists $n\ge 1$ such that $M\cong \Ker(f^n)\oplus \Im(f^n)$.

\item The following are equivalent.
\begin{itemize}
\item $M$ is indecomposable.
\item The endomorphism algebra $\End_A(M)$ does not contain any idempotents except $0$ and $\id_M$.
\item Every homomorphism $f\in \End_A(M)$ is either an isomorphism or is nilpotent.
\item $\End_A(M)$ is \defn{local} (see below).
\end{itemize}
\end{enumerate}
\end{proposition}
\begin{remark}
It is known that if $M$ is only artinian or only noetherian, then Fitting's lemma (and hence part (2)) fails.
Nevertheless, in general, the proposition still hold for $M$ that is both artinian and noetherian.  
\end{remark}

Let us briefly recall various characterisation of local rings.

\begin{definition}
A ring $R$ is \defn{local}\index{local ring} if it has a unique maximal right (equivalently, left; equivalently, two-sided) ideal.
\end{definition}
\begin{remark}
When $R$ is non-commutative, the `non-invertible elements' are the ones that do not admit (right) inverses.
\end{remark}

\begin{lemma}
The following are equivalent for a finite-dimensional algebra $A$.
\begin{itemize}
\item $A$ is local (i.e. has a unique maximal right ideal).
\item Non-invertible elements of $A$ form a two-sided ideal.
\item For any $a\in A$, one of $a$ or $1-a$ is invertible.
\item $0$ and $1_A$ are the only idempotents of $A$. 
\item $A/J(A)\cong \Bbbk$ as rings, where $J(A)$ is the two-sided ideal of $A$ given by the intersection of all maximal right (equivalently, left) ideals.
\end{itemize} 
\end{lemma}

\begin{example}\label{eg:A2 first}
Consider the upper triangular $2$-by-$2$ matrix ring
\[
A=\begin{pmatrix}
\Bbbk & \Bbbk  \\
0 & \Bbbk  
\end{pmatrix} = { \left\{(a_{i,j})_{1\leq i\leq j\leq 2}\middle| \begin{array}{l}a_{i,j}\in \Bbbk \;\forall i\leq j\\
a_{i,j}=0\;\forall i>j \end{array}\right\} }.
\]

Let $M =\{(x,y)\in \Bbbk^2\}$ be the $2$-dimensional space where $A$ acts as matrix multiplication (on the right).
Suppose $f\in\End_A(M)$, say, $f(x,y)=(ax+by,cx+dy)$ for some $a,b,c,d\in\Bbbk$.
Then being an $A$-module homomorphisms means that
\[
(ax+by,cx+dy)\begin{pmatrix}u & v \\ 0& w\end{pmatrix} = f\left((x,y)\begin{pmatrix}u & v \\ 0& w\end{pmatrix}\right) = (aux+bvx+wy,cux+dvx+dwy)
\]
for all $u,v,w,x,y\in\Bbbk$.
This means that 
\[
\begin{cases}
buy=bvx+bwy \\
avx+bvy+cxw=cux+dvx 
\end{cases}.
\]
The first line yields $b=0$, and the second line yields $c=0=b$ and $a=d$.
In other words, $\End_A(M)\cong \Bbbk$ which is clearly a local algebra.  Hence, $M$ is indecomposable.
\end{example}

A natural question is to ask when is a decomposition of modules, if it exists, unique up to permuting the direct summands.

\begin{definition}
We say that an indecomposable decomposition $M=\bigoplus_{i=1}^m M_i$ is \defn{unique} if any other indecomposable decomposition $M=\bigoplus_{j=1}^n N_j$ implies that $m=n$ and there is a permutation $\sigma$ such that $M_i\cong N_{\sigma(i)}$ for all $1\leq i \leq m$.
$\mod A$ is said to be \defn{Krull-Schmidt}\index{Krull-Schmidt} if every (finitely generated) $A$-module $M$ admits a unique indecomposable decomposition.
\end{definition}

\begin{theorem}
For a finite-dimensional algebra $A$, $\mod A$ is Krull-Schmidt.
\end{theorem}
\begin{remark}
This is a special case of the Krull-Schmidt theorem - whose proof we will omit to save time.
\end{remark}


\begin{theorem}[Krull-Schmidt]
Suppose $M=\bigoplus_{i=1}^m M_i$ is an indecomposable decomposition of $M$.  If $\End_A(M_i)$ is local for all $1\leq i\leq m$, then the decomposition of $M$ is unique.
\end{theorem}
\begin{remark}
Some people refer to this result as Krull-Remak-Schmidt theorem.
\end{remark}



\pagebreak 
%%%%%%%%%%%%%%%%%%%%%%%%%%%%%%%%%%%%%%%%%%%%%%%%%%%%%%%%%%%%%%%%%%%%%%%%%%%%%%%
%%%%%%%%%%%%%%%%%%%%%%%%%%%%%%%%%%%%%%%%%%%%%%%%%%%%%%%%%%%%%%%%%%%%%%%%%%%%%%%
%%%%%%%%%%%%%%%%%%%%%%%%%%%%%%%%%%%%%%%%%%%%%%%%%%%%%%%%%%%%%%%%%%%%%%%%%%%%%%%
%%%%%%%%%%%%%%%%%%%%%%%%%%%%%%%%%%%%%%%%%%%%%%%%%%%%%%%%%%%%%%%%%%%%%%%%%%%%%%%


\section{Simple modules, Schur's lemma}
\begin{definition}
Let $M$ be an $R$-module.
\begin{enumerate}
\item M is \defn{simple}\index{module!simple} if $M\neq 0$, and for any submodule $L\subset M$, we have $L=0$ or $L=M$.

\item $M$ is \defn{semisimple}\index{module!semisimple} if it is a direct sum of simples.
\end{enumerate}
\end{definition}
\begin{remark}
In the language of representations, simple modules are called \defn{irreducible}\index{representation!irreducible} representations, and semisimple modules are called \defn{completely reducible}\index{representation!completely reducible} representations.
\end{remark}
\begin{remark}
Note that a module is semisimple if and only if every submodule is a direct summand.
\end{remark}

\begin{example}
Consider the matrix ring $A:=\Mat_n(\Bbbk)$ over a field $\Bbbk$.  Then the row-space representation $V$ is an $n$-dimensional simple module.  Since $A_A\cong V^{\oplus n}$, we have that $A_A$ is a semisimple module.
\end{example}

\begin{example}
The \defn{ring of dual numbers}\index{dual numbers} is $A:=\Bbbk[x]/(x^2)$.  The module $(x)$ is simple.  The regular representation $A$ is non-simple (as $(x)=AxA$ is a non-trivial submodule).  It is also not semisimple.  Indeed, $(x)$ is a submodule of $A$, and the quotient module can be described by $\Bbbk v$ where $v=1+(x)$.  If $A$ is semisimple, then the 1-dimensional space $\Bbbk v$ is isomorphic to a submodule of $A$.  Such a submodule must be generated by $a+bx$ (over $A$) for some $a,b\in \Bbbk$.  If $a\neq 0$, then $(a+bx)A=A$.  So $a=0$, and $\Bbbk v\cong (x)$, a contradiction.
\end{example}

\begin{lemma}
$S$ is a simple $A$-module if and only if for any non-zero $m\in S$, we have $mA:=\{ma\mid a\in A\}=S$.  In particular, simple modules are cyclic (i.e. generated by one element).
\end{lemma}
%\begin{proof}
%$\Rightarrow$:  $mA\subset S$ is a submodule and contains a non-zero element $m$, so by simplicity of $S$ we must have $mA=S$.
%
%$\Leftarrow$: Suppose that there is a non-zero submodule $L\subset S$.  For a non-zero element $m\in L$, the assumption says that we have $mA \subset L\subset S = mA$, and so $L=S$.
%\end{proof}

Let us see how one can find a simple module.

\begin{definition}
Let $M$ be an $A$-module and take any $m\in M$.
The \defn{annihilator}\index{annihilator} of $m$ (in $A$) is the set $\Ann_A(m):=\{a\in A\mid ma=0\}$.
\end{definition}

Note that $\Ann_A(m)$ is a right ideal of $A$ - hence, a right $A$-module.

\begin{lemma}\label{lem:simple as A/Ann}
For a simple $A$-module $S$ and any non-zero $m\in S$, we have $S\cong A/\Ann_A(m)$ as $A$-module.  In particular, if $A$ is finite-dimensional, then every simple $A$-module is also finite-dimensional.
\end{lemma}
%\begin{proof}
%Since $S=mA$, the element $m$ defines a surjective $A$-module homomorphism $f:A_A\to S$ given by $a\mapsto ma$.
%On the other hand, we have $\Ker(f) = \Ann_A(m)$, and so $A/\Ann_A(m)\cong S$.
%\end{proof}

Suppose $I$ is a two-sided ideal of $A$.
Then we have a quotient algebra $B:=A/I$.
For any $B$-module $M$, we have a canonical $A$-module structure on $M$ given by $ma:= m(a+I)$.
This is (somewhat confusingly) the \index{restriction}\defn{restriction of $M$ along the algebra homomorphism $A\twoheadrightarrow A/I$}.

\begin{lemma}
Suppose $B:=A/I$ is a quotient algebra of $A$ by a strict two-sided ideal $I\neq A$.
If $S\in \mod B$ is simple, then $S$ is also simple as $A$-module
\end{lemma}
\begin{proof}
This follows from the easy observation that any a $B$-submodule of $S_B$ is also a $A$-submodule of $S_A$ under restriction.
\end{proof}

The following easy, yet fundamental, lemma describes the relation between simple modules.
Recall that a division ring is one where every non-zero element admits an inverse (but the ring is not necessarily commutative).

\begin{lemma}[Schur's lemma]\index{Schur's lemma}\label{lem:Schur}
Suppose $S,T$ are simple $A$-modules, then
\[
\Hom_A(S,T) = \begin{cases}\text{a division ring},& \text{if }S\cong T;\\
0, & \text{otherwise.}
\end{cases}
\]
\end{lemma}
\begin{remark}
Note that if $A$ is an $R$-algebra, then the division ring appearing is also an $R$-algebra (since it is the endomorphism ring of an $A$-module).
In particular, if $R$ is an algebraically closed field $\Bbbk=\overline{\Bbbk}$, then any division $\Bbbk$-algebra is just $\Bbbk$ itself.
\end{remark}


\begin{proof}
The claim is equivalent to saying that any $f\in\Hom_A(S,T)$ is either zero or an isomorphism.
Since $\Im(f)$ is a submodule of $T$, simplicity of $T$ says that $\Im(f)=0$, i.e. $f=0$, or $\Im(f)\cong T$.
In the latter case, we can consider $\Ker(f)$, which is a submodule of $S$, so by simplicity of $S$ it is either $0$ or $S$ itself.
But this cannot be $S$ as this means $f=0$, hence, $\Im(f)\cong T$ implies that $\Ker(f)=0$, i.e. $f$ is an isomorphism.
\end{proof}

\begin{example}
In Example \ref{eg:A2 first}, we showed that the upper triangular 2-by-2 matrix ring $A$ has a 2-dimensional indecomposable module $P_1=\{(x,y)\mid x,y\in\Bbbk^2 \}$ given by `row vectors'.
It is straightforward to check that there is a 1-dimensional (hence, simple) submodule given by $S_2:=\{(0,y)\mid y\in\Bbbk^2\}$.

Consider the module $S_1:=P_1/S_2$.  This is a 1-dimensional (simple) module spanned by, say, $w$ with $A$-action given by 
\[
w\begin{pmatrix}
a&b\\ 0&c 
\end{pmatrix} := wa.
\]

Consider a homomorphism $f\in \Hom_A(S_1,S_2)$.
This will be of the form $w\mapsto (0,y)$ for some $y \in \Bbbk$ and has to satisfy
\[
(0,ya)=(0,y)a=f(wa)=f( w \begin{pmatrix}
a&b\\ 0&c 
\end{pmatrix} ) = f(w)\begin{pmatrix}
a&b\\ 0&c 
\end{pmatrix} = (0,y)c = (0,yc)
\]
for any $a,b,c\in \Bbbk$.
Hence, we must have $y=0$, which means that $f=0$.
In particular, by Schur's lemma $S_1\ncong S_2$.
\end{example}

\begin{lemma}\label{lem:endo ss}
Suppose that $S$ is a simple $A$-module.
Consider a semisimple $A$-module $M=S_1\oplus \cdots \oplus S_n$ with $S_i\cong S$ for all $i$.
Then $\End_A(M)\cong \Mat_{n}(D)$, where $D:=\End_A(S)$.
\end{lemma}
\begin{proof}
We have canonical inclusion $\iota_j: S_j\into M$ and projection $\pi_i:M\onto S_i$.  So for $f\in \End_A(M)$, we have a homomorphism $\pi_i f \iota_j: S_j\to S_i$, and by Schur's lemma, this is an element of $D$.  Now we have a ring homomorphism
\[
\End_A(M) \to \Mat_r(D), \quad f\mapsto (\pi_i f \iota_j)_{1\leq i,j\leq r},
\]
which is clearly injective.
Conversely, for $(a_{i,j})_{1\leq i,j\leq r}\in\Mat_{r}(D)$, we have an endomorphism $M\stackrel{\pi_j}{\onto} S_j \stackrel{a_{i,j}}{\to} S_i \stackrel{\iota_i}{\into} M$, which yields the required surjection.
\end{proof}

\begin{example}
For a tautological example, take $A=\Bbbk$ to be just a field.  Then we have a 1-dimensional simple $A$-module $S=\Bbbk$ with
$\End_A(S^{\oplus n})=\Mat_n(\End_A(\Bbbk))=\Mat_n(\Bbbk)$.  Note that now we have an $n$-dimensional simple $\Mat_n(\Bbbk)$-module (given by the row vectors).
\end{example}



\pagebreak
%%%%%%%%%%%%%%%%%%%%%%%%%%%%%%%%%%%%%%%%%%%%%%%%%%%%%%%%%%%%%%%%%%%
%%%%%%%%%%%%%%%%%%%%%%%%%%%%%%%%%%%%%%%%%%%%%%%%%%%%%%%%%%%%%%%%%%%
%%%%%%%%%%%%%%%%%%%%%%%%%%%%%%%%%%%%%%%%%%%%%%%%%%%%%%%%%%%%%%%%%%%


\section{Quiver and path algebra}

\begin{definition}
A (finite) \defn{quiver} is a datum $Q=(Q_0, Q_1, s,t:Q_1\to Q_0)$ for finite sets $Q_0, Q_1$.  The elements of $Q_0$ are called \defn{vertices} and those of $Q_1$ are called \defn{arrows}.  The \defn{source} (resp. \defn{target})of an arrow $\alpha\in Q_1$ is the vertex $s(\alpha)$ (resp. $t(\alpha)$).
\end{definition}
This is equivalent to specifying an oriented graph (possibly with multi-edges and loops); Gabriel coined the term quiver as a way to emphasise the context is not really about the graph itself.

\begin{definition}
Let $Q$ be a quiver.
\begin{itemize}
\item A \defn{trivial path}\index{path!trivial} on $Q$ is a ``stationary walk at $i$", denoted by $e_i$ for some $i\in Q_0$.

\item A \defn{path} of $Q$ is either a trivial path or a word $\alpha_1\alpha_{2}\cdots \alpha_{\ell}$ of arrows with $s(\alpha_i)=t(\alpha_{i+1})$.
\end{itemize}
\end{definition}

The source and target functions extend naturally to paths, with $s(e_i)=i=t(e_i)$.
Two paths $p,q$ can be concatenated to a new one $pq$ if $t(p)=s(q)$; note that our convention is to read \underline{\emph{from left to right}}.

\begin{definition}
The \defn{path algebra} $\Bbbk Q$ of a quiver $Q$ is the $\Bbbk$-algebra whose underlying vector space is given by $\bigoplus_{p:\text{paths of }Q} \Bbbk p$, with multiplication given by path concatenation.  That is $x\in \Bbbk Q$ is a formal linear combinations of paths on $Q$.
\end{definition}
%Note that the trivial paths $e_i$ are primitive idempotents of $KQ$, and the radical $\rad KQ$ of $KQ$ is the ideal generated by all arrows.

Note that $e_i e_j = \delta_{i,j}e_i$, where $\delta_{i,j}=1$ if $i=j$ else $0$.
In other words, $e_i$ is an \defn{idempotent}\index{idempotent} of the path algebra $\Bbbk Q$.
Moreover, we have an idempotent decomposition
\[
1_{\Bbbk Q} = \sum_{i\in Q_0} e_i \]
of the unit element of $\Bbbk Q$.


\begin{example}
Consider the \defn{one-looped quiver}, a.k.a. \defn{Jordan quiver},
\[ Q= \Big(\quad \vcenter{\xymatrix{ \bullet \ar@(ul,ur)^{\alpha}}}\quad\Big) \]
Then $\Bbbk Q$ has basis $\{\alpha^k\mid k\geq 0\}$ (note that the trivial path at the unique vertex is the identity element).
Then $\Bbbk Q\cong \Bbbk[x]$.
\end{example}

An \defn{oriented cycle}\index{oriented cycle} is a path of the form $v_1\to v_2\to \cdots v_r\to v_1$, i.e. starts and ends at the same vertex.
If $Q$ does not contain any oriented cycle, we say that it is \defn{acyclic}\index{acyclic}.

\begin{proposition}
$\Bbbk Q$ is finite-dimensional if, and only if, $Q$ is finite acyclic.
\end{proposition}
\begin{proof}
If there is an oriented cycle $c$, then $c^k\in \Bbbk Q$ for all $k\geq 0$, and so $\Bbbk Q$ is infinite-dimensional.
Otherwise, there are only finitely many paths on $Q$.
\end{proof}


\begin{example}
Consider the linearly oriented $\linA_n$-quiver
\[
Q = \linA_n = 1 \xrightarrow{\alpha_1} 2 \xrightarrow{\alpha_2}\cdots \xrightarrow{\alpha_{n-1}} n .
\]
Then the path algebra $\Bbbk Q$ has basis $\{ e_i, \alpha_{j,k}\mid 1\leq i\leq n, 1\leq j\leq k\leq n \}$, where $\alpha_{j,k}:=\alpha_j\alpha_{j+1}\cdots \alpha_k$.

Consider the upper triangular $n$-by-$n$ matrix ring
\[
\begin{pmatrix}
\Bbbk & \Bbbk & \cdots & \Bbbk \\
0 & \Bbbk & \cdots & \Bbbk \\
0 & 0 & \ddots & \vdots \\
0 & 0 & 0 & \Bbbk
\end{pmatrix} = { \left\{(a_{i,j})_{1\leq i\leq j\leq n}\middle| \begin{array}{l}a_{i,j}\in \Bbbk \;\forall i\leq j\\
a_{i,j}=0\;\forall i>j \end{array}\right\} }.
\]
Denote by $E_{i,j}$ the elementary matrix whose entries are all zero except at $(i,j)$ where it is one.
This ring is isomorphic to $\Bbbk Q$ via $E_{i,i}\mapsto e_i$ and $E_{i,j}\mapsto \alpha_{i,j-1}$ for $1\leq j<k\leq n$.
\end{example}


From now on, we will focus in the following setting.

\begin{assumption}
\begin{enumerate}
\item Quivers are finite (i.e. finitely many vertices and arrows). 
\item Representations (equivalently, modules) are finite-dimensional.
\end{enumerate}
%\begin{enumerate}
%\item Quivers are always finite.

%\item Algebras are finite-dimensional over an algebraically closed field $\Bbbk$.

%\item Representations are finite-dimensional.
%\end{enumerate}
\end{assumption}


\section{Duality}

For a quiver $Q$, the \defn{opposite quiver} $Q^\op$ has the same set of vertices with the reverse direction of arrows, i.e. $Q_0^\op=Q_0, Q_1^\op=Q_1, s_{Q^\op}=t_Q$, and $t_{Q^\op}=s_Q$.

\begin{exercise}
Show that there is a canonical isomorphism $(\Bbbk Q)^\op \cong \Bbbk (Q^\op)$.
\end{exercise}

Let $M$ be a finite-dimensional $A$-module.  Then we have a dual space
\[
D(M):= M^* := \Hom_{\Bbbk}(M,\Bbbk),
\]
which has a natural $A^\op$-module structure, namely, $(a\cdot f)(m):=f(ma)$
for any $a\in A, f\in M^*, m\in M$.
Moreover, for an $A$-module homomorphism $\theta:M\to N$, we have also an $A^\op$-module homomorphism $\theta^*:N^*\to M^*$ with $\theta^*(f)(m) = f(\theta(m))$.

\begin{lemma}
There is a $\Bbbk$-vector space isomorphism $\Hom_A(M,N)\cong \Hom_{A^\op}(DN,DM)$.
\end{lemma}
\begin{proof}
Just a straightforward check that $(\theta^*)^*=\theta$.
\end{proof}

We note as a fact that $D$ preserves indecomposability of (finite-dimensional) modules.  This can be seen using the fact that $\Hom_A(M,N)\cong \Hom_{A^\op}(DN,DM)$ and can be upgraded to an algebra isomorphism for the case when $N=M$; then uses characterisation of indecomposable module by local endomorphism ring.
%We note the following facts without detail proof (since cleanest proof uses the language in category theory).
%\begin{itemize}
%\item $\Hom_A(M,N)\cong \Hom_{A^\op}(DN,DM)$ for all $M,N\in\mod A$.
%\item $D$ preserves indecomposability.
%\end{itemize}
%\begin{lemma}
%For $A,M,N$ finite-dimensional, we have $\Hom_A(M,N)\cong \Hom_{A^\op}(DN,DM)$.
%In particular, $D$ preserves indecomposablity of finite-dimensional modules.
%\end{lemma}
%\begin{proof}
%We have a map $\Hom_A(M,N)\to \Hom_{A^\op}(DN,DM) \to \Hom_{{A^\op}^\op}(D(DM),D(DN))$.
%Check that the evaluation map
%\[
%M\to D(DM) \text{ which sends } m\mapsto (DM \ni f\mapsto f(m)) 
%\]
%is an isomorphism of $A$-modules.
%\end{proof}


\begin{example}
The \emph{left} $A$-module ${}_AA$ yields a right $A$-module structure on $D(A)$.
More generally, suppose we have a left ideal $Ae$ of $A$ for some element $e\in A$, then $D(Ae)$ is a right ideal of $A$.
\end{example}

\begin{remark}
There is another natural duality, which we will not used, between $\mod A$ and $\mod A^\op$ given by sending $M$ to $\Hom_A(M,A)$.
In general, this duality is different from the $\Bbbk$-linear dual unless $A$ is a so-called \emph{symmetric algebra}, meaning that $A\cong DA$ as bimodule; in which case, $\Hom_A(-,A)$ dual is naturally isomorphic to $D$ (as functors).
\end{remark}


\section{Representations of quiver}

\begin{definition}
A $\Bbbk$-linear \defn{representation}\index{representation!of quiver} of $Q$ is a datum $(\{M_i\}_{i\in Q_0}, \{M_{\alpha}\}_{\alpha\in Q_1})$ where $M_i$ is a $\Bbbk$-vector space for each $i\in Q_0$ and $M_\alpha: M_{s(\alpha)}\to M_{t(\alpha)}$ is $\Bbbk$-linear map for each $\alpha\in Q_1$.

Such a representation is \defn{finite-dimensional} if $\dim_{\Bbbk} M_i<\infty$ for all $i\in Q_0$.
\end{definition}

\begin{notation}
For a representation $M$ of $Q$, we take $M_p:=M_{\alpha_1}\cdots M_{\alpha_\ell}$ for a path $p=\alpha_1\cdots \alpha_\ell$.
\end{notation}

It is easy to notice that every representation of $Q$ is equivalent to a $\Bbbk Q$-module, namely, 
\[
\text{representation }(\{M_i\}_{i\in Q_0}, \{M_{\alpha}\}_{\alpha\in Q_1}) \leftrightarrow \begin{array}{c}\text{$\Bbbk Q$-module }\prod_{i\in Q_0} M_i \\ \text{s.t. }\sum_{p:\text{path}} \lambda_p p\text{ acts as }\sum_p \lambda_pM_p.\end{array}
\]


\begin{example}[Simple]
For $x\in Q_0$, denote by $S_x$ (or $S(x)$) the representation given by putting a 1-dimensional space on $x$, zero on all other vertices, and zero on all arrows.
This corresponds to a 1-dimensional $\Bbbk Q$-module and so we call it the \defn{simple at $x$}\index{simple}.
\end{example}

Note: at this stage, it is not clear if these are all the simple $\Bbbk Q$-modules (up to isomorphism) yet.

\begin{example}[Projective]
For $x\in Q_0$, denote by $P_x$ (or $P(x)$) the representation given by
$( \{M_y\}_{y\in Q_0}, \{M_\alpha\}_{\alpha\in Q_1} )$, where
\begin{align*}
M_y := \bigoplus_{\substack{p:\text{path with }\\s(p)=x,\\t(p)=y}} \Bbbk p, \qquad \text{and}\qquad 
(M_\alpha : M_y \to M_z) := \sum_{ p\alpha=q } (M_y\twoheadrightarrow \Bbbk p \xrightarrow{\id} \Bbbk q \hookrightarrow M_z).
\end{align*}
This is called the \defn{projective at $x$}\index{projective}.
This corresponds to the right ideal $e_x\Bbbk Q$ of $\Bbbk Q$.
\end{example}

\begin{example}[Injective]
Dual to the projective module construction, for $x\in Q_0$, denote by $I_x$ (or $I(x)$) the representation given by
$( \{M_y\}_{y\in Q_0}, \{M_\alpha\}_{\alpha\in Q_1} )$, where
\begin{align*}
M_y := \bigoplus_{\substack{p:\text{path with }\\s(p)=y,\\t(p)=x}} \Bbbk p, \qquad \text{and}\qquad 
(M_\alpha : M_y \to M_z) := \sum_{ p=\alpha q } (M_y\twoheadrightarrow \Bbbk p \xrightarrow{\id} \Bbbk q \hookrightarrow M_z).
\end{align*}
This is called the \defn{injective at $x$}\index{injective}.
This corresponds to the dual of the left ideal generated by $e_x$, i.e. $D(\Bbbk Q e_x)$.
\end{example}

\begin{example}
The representation of $Q=\vec{\mathbb{A}}_n$ given by 
\[
U_{i,j}:=0\to \cdots 0\to \Bbbk \xrightarrow{\id} \to \cdots \xrightarrow{\id} \Bbbk \to 0\to\cdots \to 0
\]
with a copy of $\Bbbk$ on vertices $i,i+1,\ldots, j$ is the uniserial $\Bbbk Q$-module corresponding to the column space (under the isomorphism of $\Bbbk Q$ with the lower triangular matrix ring) with non-zero entries in the $k$-th row for $i\leq k\leq j$.
\end{example}

\begin{example}
Let $Q$ be the Jordan quiver with unique arrow $\alpha$.
Then a representation of $Q$ is nothing but an $n$-dimensional vector space equipped with a linear endomorphism, equivalently, an $n$-by-$n$ matrix.
\end{example}

\begin{definition}
A \index{homomorphism!of quiver representation}\defn{homomorphism} $f:M\to N$ of ($\Bbbk$-linear) quiver representations $M=(M_i,M_\alpha)_{i,\alpha}$ and $N=(N_i,N_\alpha)_{i,\alpha}$ is a collection of linear maps $f_i:M_i\to N_i$ that intertwines arrows' actions, i.e. we have a commutative diagram
\[
\xymatrix{ M_i \ar[r]^{f_i} \ar[d]_{M_\alpha} & N_i\ar[d]^{N_\alpha}\\
M_j \ar[r]_{f_j} & N_j}
\]
for all arrows $\alpha:i\to j$ in $Q$.

A homomorphism $f=(f_i)_{i\in Q_0}:M\to N$ of quiver representations is \defn{injective}, resp. \defn{surjective}, resp. an \defn{isomorphism}, if every $f_i$ is injective, resp. surjective, resp. an isomorphism, for all $i\in Q_0$.
\end{definition}

\begin{example}
Let $Q$ be the Jordan quiver.
Recall that a representation of $Q$ is equivalent to a choice of $n$-by-$n$ matrix $M_\alpha$.
By definition, the isomorphism class of such a representation is given by the conjugacy classes of $M_\alpha$.
If we assume $\Bbbk$ is algebraically closed, then a representative of the isomorphism class of $M_\alpha$ is given by the Jordan normal form of $M_\alpha$.  That is, $M_\alpha$ can be block-diagonalise into Jordan blocks $J_{m_1}(\lambda_1),\ldots, J_{m_l}(\lambda_l)$, where $J_m(\lambda)$ is the $m$-by-$m$ Jordan block with eigenvalue $\lambda\in \Bbbk$.
\end{example}

\begin{proposition}
There is an isomorphism between the category of representations of $Q$ and $\mod\Bbbk Q$, where $(M_i,M_\alpha)_{i,\alpha}$ corresponds to $M=\prod_{i\in Q_0} M_i$ with $\Bbbk Q$-action given by (linear combinations of compositions of) $M_\alpha$'s, and isomorphism classes of $Q$-representations correspond to isomorphism classes of $\Bbbk Q$-modules.
\end{proposition}



\pagebreak

\section{Idempotents}

Recall that an \index{idempotent}\defn{idempotent} of an algebra $A$ is an element $x$ with $x^2=x$.

The right $A$-modules of the form $eA$ and $D(Ae)$ for an idempotent $e\in A$ are of central importance in representation theory and in homological algebra.

\begin{lemma}\label{lem:yoneda and idem trunc}
The the following hold for any idempotent $e\in A$.
\begin{enumerate}
\item {\rm (Yoneda's lemma)}\index{Yoneda's lemma} $\Hom_A(eA,M)\cong Me$ as a $\Bbbk$-vector space for all $M\in \mod A$.

\item There is an isomorphism of rings $\End_A(eA)\cong eAe$.
\end{enumerate}
\end{lemma}
\begin{proof}
For (1), check that $\Hom_A(eA,M)\ni f\mapsto f(e)=f(1)e\in Me$ defines a $\Bbbk$-linear map with inverse $me\mapsto (ea\mapsto mea)$.  (2) follows from (1) by putting $M=eA$ with straightforward check of correspondence of multiplication on both sides.
\end{proof}
\begin{remark}
Under the isomorphism $A\cong \End_A(A)$, an idempotent $e$ of $A$ corresponds to the `project to direct summand $P=eA$ endomorphism', i.e. $A \twoheadrightarrow P \hookrightarrow A$.  This is compatible with Yoneda lemma (think about this!) which says that there is a vector space isomorphism $fAe\cong \Hom_A(eA, fA)$ for any idempotents $e,f$.
\end{remark}

\begin{lemma}\label{lem:idem conj}
For idempotents $e,f\in A$, we have $eA\cong fA$ as right $A$-module if and only if $f=ueu^{-1}$ for some unit $u\in A^\times$.
\end{lemma}
%\begin{proof}
%$\Leftarrow$: By Yoneda lemma, an isomorphism $\phi \in \Hom_A(fA,eA)$ corresponds to an element in $x\in eAf\subset A$; likewise an isomorphism $\psi \in \Hom_A((1-f)A,(1-e)A)$ corresponds to $y\in (1-e)A(1-f)\subset A$.  
%Let $x'\in fAe$ and $y'\in (1-f)A(1-e)$ be the elements corresponding to $\phi^{-1}$ and $\psi^{-1}$ respectively.
%Since $\phi^{-1}\phi = \id_{eA}$ corresponds to $e\in eAe$, we have
%\[
%x'x = f, xx'=e, y'y=1-f, yy'=1-e.
%\]
%Take $u:=x+y$ and $v:=x'+y'$.  Then we have $vu = f+(1-f) = 1$ and $uv=e+(1-e)=1$.  Therefore, $u,v$ are units such that $uf = x = eu$, i.e. $e=ufu^{-1}$ as required.
%
%$\Rightarrow$: The required isomorphism $fA\to eA$ is given by $fa\mapsto eua$.
%\end{proof}

Given an idempotent $e=e^2\in A$ in an algebra $A$, then $eA$ and $(1-e)A$ are both right ideal of $A$.  Since $e(1-e)=0=(1-e)e$, we have $eA\cap (1-e)A=0$, which means that $A \cong eA\oplus (1-e)A$ as right $A$-module.  
In particular, in the setting of the above lemma, we have that $eA\cong fA$ and $(1-e)A\cong (1-f)A$ by Krull-Schmidt property.
%By Lemma \ref{lem:projective equiv cond} both $Ae$ and $A(1-e)$ are then projective $A$-modules.
%This leads to the following characterisation of idempotent that yields \emph{indecomposable} projective modules.

\begin{definition}
Two idempotents $e,f$ are \defn{orthogonal}\index{idempotent!orthogonal} if $ef=0=fe$.
An idempotent $e$ is \defn{primitive}\index{idempotent!primitive} if $e\neq f+f'$ for some orthogonal (pair of) idempotents $f,f'$.
\end{definition}

It follows from the definition of primitivity that
\[
\text{$eA$ and $D(Ae)$ are indecomposable $A$-modules for a primitive idempotent $e$}.
\]

\begin{example}
The trivial paths $e_x$ for $x\in Q_0$ is (by design) a primitive idempotent of the path algebra $\Bbbk Q$, and $1=\sum_{x\in Q_0}e_x$ is an orthgonal decomposition of primitive idempotents.
Hence, we have a decomposition \[
\Bbbk Q\cong \bigoplus_{x\in Q_0} e_x \Bbbk Q = \bigoplus_{x\in Q_0} P_x\text{ and }D(\Bbbk Q)\cong \bigoplus_{x\in Q_0} D(\Bbbk Q e_x) \cong \bigoplus_{x\in Q_0} I_x.
\]
\end{example}

%
%\begin{example}
%Consider $A=\Bbbk\linA_n$.  Then we have the paths starting from $i\in\{1,\ldots, n\}=Q_0$ are $\alpha_{i,j}$ for $j\geq i$, and $P_i=U_{i,n}$.
%By considering the correspondence with upper triangular matrices, we see that $P_i$ is just the `$i$-th row space' of upper triangular matrices.
%\end{example}

%Indecomposable projective modules - as they are direct summands of $A$ - can be regarded as the `largest unbreakable building block' (not in the sense of dimension, but from the Jordan-H\"{o}lder filtration perspective) of $A$-modules, whereas a simple $A$-modules are the smallest unbreakable building block.
%The following part details their relation.





\pagebreak
%%%%%%%%%%%%%%%%%%%%%%%%%%%%%%%%%%%%%%%%%%%%%%%%%%%%%%%%%%%%%%%%%%%
%%%%%%%%%%%%%%%%%%%%%%%%%%%%%%%%%%%%%%%%%%%%%%%%%%%%%%%%%%%%%%%%%%%
%%%%%%%%%%%%%%%%%%%%%%%%%%%%%%%%%%%%%%%%%%%%%%%%%%%%%%%%%%%%%%%%%%%


\section{Composition series, Jordan-H\"{o}lder Theorem}

\begin{definition}
Let $A$ be a $\Bbbk$-algebra and $M\in A\mod$.
A \defn{composition series}\index{composition!series} of $M$ is a \underline{finite} chain of submodules
\[
0=M_{0} \subset M_{1} \subset \cdots \subset M_\ell = M
\]
such that $M_i/M_{i-1}$ is simple for all $1\leq i \leq \ell$.   The number $\ell$ here is the \defn{length}\index{length} of the composition series.  The module $M_i/M_{i-1}$ for each $1\leq i\leq \ell$ are called the \defn{composition factors}\index{composition!factor} of the series.
\end{definition}

%Composition series allows us to understand the structure of a module by simple modules.  It is desirable to have a rigidity result, namely, (1) composition factors do not depends on the choice of composition series, and (2) the notion of length is well-defined.

\begin{theorem}[Jordan-H\"{o}lder Theorem]
Any two composition series have the same length and the multi-sets of their composition factors (up to isomorphisms) are the same.
\end{theorem}
We omit the proof.  The strategy is basically by induction on the length of series.

%\begin{proof}
%Suppose we have two composition series
%\[
%\begin{array}{rcccccl}
%0=M_0 & \subset & M_1 & \subset & \cdots & \subset & M_\ell = M, \\
%0=N_0 & \subset & N_1 & \subset & \cdots & \subset & N_n = M. \\
%\end{array}
%\]
%Without loss of generality, we can assume $n>\ell$.  We claim that $N_\ell = M$. Indeed, we can do this by induction on $\ell$.
%If $\ell=0$, then clearly $M_0=0=N_0$ and we are done; likewise, when $\ell=1$, then $M$ is simple and we have $M_1=M=N_1$.  For $\ell>1$, we have
%\[
%0 = M_1\cap N_0 \subset M_1\cap N_1 \subset \cdots \subset  M_1\cap N_n = M_1\cap M = M_1.
%\]
%So as $M_1$ simple, there is some $n_0$ such that $N_{n_0}\cap M_1= M_1$ and $N_{i}\cap M_1=0$ for all $i<n_0$.
%
%We now consider two new chains
%\[
%\begin{array}{rcccccccccccccl}
%0 & \subset & \frac{M_2}{M_1} & \subset & \cdots & \cdots& \cdots& \cdots& \cdots& \cdots& \cdots & \cdots &\cdots & \subset & \frac{M_\ell}{M_1} = \frac{M}{M_1}, \\
%0 & \subset & N_1 & \subset & \cdots & \subset & N_{n_0-1} & \subset & \frac{N_{n_0+1}}{M_1} &\subset & \frac{N_{n_0+2}}{M_1} & \subset & \cdots & \subset & \frac{N_n}{M_1} = \frac{M}{M_1}, \\
%\end{array}
%\]
%which are both composition series of $M/M_1$.  By induction hypothesis, we thus have $n-1=\ell-1$ and the composition factors of these two series coincide up to permutation.
%\end{proof}

\begin{remark}
Jordan-H\"{o}lder theorem holds as long as a module, regardless of what kind of algebra, has a (finite) composition series; this condition is actually equivalent to saying that it is noetherian and artinian.
\end{remark}

\begin{remark}
The Jordan-H\"{o}lder theorem may not hold if one relaxes the form of composition factors from simple modules to something else.  There are a few active research themes, including one related to quasi-hereditary algebras, that are stemmed from this.
\end{remark}

\begin{lemma}\label{lem:comp series exists}
Let $M$ be a finite-dimensional right $A$-module.
Then $M$ has a composition series.
\end{lemma}
\begin{proof} Induction on $\dim_{\Bbbk} M$, at each step choose a maximal submodule (i.e. a submodule whose quotient is simple).
%This is by induction on $\dim_K M$.  For $\dim_K M=0$ this is trivial.  For $\dim_K M>0$, if $M$ is simple, then we are done.  Otherwise, $M$ proper non-zero submodule, and we pick $N$ such a submodule so that $M/N$ is simple.  Clearly $\dim_K N<\dim_K M$ and so we can apply induction hypothesis.
\end{proof}

\begin{example}
Let $A=\Bbbk \linA_n$.  Then the module $U_{i,j}$ has a composition series
\[
0 \subset U_{j,j} \subset U_{j-1,j} \subset \cdots \subset U_{i+1,j} \subset U_{i,j}
\]
with composition factors $S_{k} = U_{k,j}/U_{k+1,j}$ for $i\leq k \leq j$.
Note that this composition series is unique - such kind of modules are called \defn{uniserial}\index{uniserial}.
\end{example}

\begin{lemma}\label{lem:submod comp ser}
If $M\in \mod A$ and $N\subset M$ is a submodule, then there is a composition series $(M_i)_{0\leq i\leq \ell}$ so that $N=M_k$ for some $0\leq k\leq \ell$.
\end{lemma}
\begin{proof}
$N$ has a composition series, say, of length $k$, so we take that as the first $k$ terms of the required composition series of $M$.
On the other hand, $M/N$ also has a composition series, and since every submodule of $M/N$ is of the form $L/N$ (for a submodule $U$ of $M/N$, take $L:=\{m\in M\mid m+N\in U\}$; it is routine to check that this is an inverse operation as quotienting $N$ on the submodules of $M$ that contains $N$), a composition series of $M/N$ is of the form $(L_i/N)_{0\leq i\leq r}$.  Now take $M_{k+i} = L_i$.
\end{proof}

\begin{proposition}
Suppose $A$ is a $\Bbbk$-algebra such that $A_A$ has a composition series.
Then there are only finitely many simple $A$-modules up to isomorphisms, and they all appear in the form $A/I$ for some $A$-submodule $I$ of $A$.
\end{proposition}
Note that while this does not require $A$ to be finite-dimensional, it requires $A_A$ to be of finite length (equivalently, noetherian and artinian).
\begin{proof}
The final clause of the claim is just restating Lemma \ref{lem:simple as A/Ann}:  any simple $S$ is given by $A/\Ann_A(m)$ for any non-zero $m\in S$.
Now fix such an $S$ and $I:=\Ann_A(m)$.
Since $A$ has a composition series, $I$ also have one by Lemma \ref{lem:submod comp ser} so that the series ends with $I\subset A$.
Since this is possible for any simple $S$, it follows from Jordan-H\"{o}lder theorem that all simple modules other than $S$ must appear as composition factors of $I$.

Since composition series is a finite chain, there must be finitely many composition factors - hence, the simple modules of $A$ must be finite.
\end{proof}

%To obtain all simple $A$-modules (up to isomorphism), we will look at maximal submodules of $A_A$ - or better - maximal submodules of the indecomposable module $eA$ for a primitive idempotent $e$.
%
%Let us do this for path algebra of finite acyclic quiver first.
%\begin{proposition}
%Consider $x\in Q_0$ of a finite acyclic quiver $Q$.
%Then $P_x=e_x\Bbbk Q$ has a unique submodule $\rad(P_x)$ spanned by the set $e_xQ_{\geq 1}$ of all paths of length at least one starting from $x$, and the cokernel of the canonical inclusion is the simple module $S_x$.
%In particular, $\{S_x\mid x\in Q_0\}$ is the complete set of simple $\Bbbk Q$-modules up to isomorphisms.
%\end{proposition}
%\begin{proof}
%$P_x=e_x\Bbbk Q$ has basis given by all paths starting from $x$.  
%Since all the action of any arrow in $Q$ increases the length of the path, $e_x\Bbbk Q_{\geq 1}$ is a submodule of $P_x$.
%The cokernel of this inclusion is 1-dimensional and spanned by the coset $e_x+(e_x\Bbbk Q_{\geq 1})$; it is easy to check that this cokernel is isomorphic to $S_x$.
%
%Note that as $Q$ is acyclic, there is no path starting from $x$ that ends also at $x$ apart from $e_x$, so $\rad(P_x)e_x=0$.
%Let $M\subsetneqq P_x$ be a strict submodule.  
%We need to shot that $M\subset \rad(P_x)$.  
%Suppose not, then there is some $m=\lambda e_x+a \in M$ with a non-zero $\lambda\in \Bbbk$ and $a\in \rad(P_x)$.
%Then we have 
%\[
%me_x = \lambda e_x^2 + ae_x = \lambda e_x 
%\]
%and so $M\supset e_xA$, which means that $M=e_x A$ - a contradiction.
%
%Now it follows that the maximal $A$-submodule of $A_A$ are of the form $\rad(P_x)\oplus \bigoplus_{y\neq x}P_y$ for each $x\in Q_0$.
%Hence, $\{S_x\mid x\in Q_0\}$ are all the (isomorphism-class representatives of) simple $\Bbbk Q$-modules.
%\end{proof}

\pagebreak
%%%%%%%%%%%%%%%%%%%%%%%%%%%%%%%%%%%%%%%%%%%%%%%%%%%%%%%%%%%%%%%%%%%
%%%%%%%%%%%%%%%%%%%%%%%%%%%%%%%%%%%%%%%%%%%%%%%%%%%%%%%%%%%%%%%%%%%
%%%%%%%%%%%%%%%%%%%%%%%%%%%%%%%%%%%%%%%%%%%%%%%%%%%%%%%%%%%%%%%%%%%

\section{Semisimplicity and Artin-Wedderburn theorem}

In order to obtain all (isomorphism classes of) simple $A$-modules - or equivalently maximal right $A$ ideal (i.e. maximal submodules of $A_A$) - for a finite-dimensional $\Bbbk$-algebra $A$, we will use the following.

\begin{definition}
Let $A$ be a $\Bbbk$-algebra and $M\in \mod A$.
\begin{enumerate}
\item The \defn{(Jacobson) radical}\index{radical}\index{Jacobson radical} $\rad(A)$ (sometimes also written as $J(A)$) of $A$ is the intersection of all maximal right ideals (i.e. maximal $A$-submodules) of $A$.

\item $A$ is \defn{semisimple}\index{algebra!semisimple} if $\rad(A)=0$.
\end{enumerate}
\end{definition}

\begin{example}
For $A=\Bbbk Q$ of a finite quiver $Q$ and $x\in Q_0$.
The projective $P_x$ at $x$ contains a submodule spanned by all paths starting from $x$ with length at least 1.
This is a maximal submodule of $P_x$ since the cokernel of the natural embedding to $P_x$ is a one-dimensional module spanned by the coset of $e_x$ -- in particular, this simple module is isomorphic to $S_x$.
Thus, we have $\rad(A)=\Bbbk Q_{\geq 1}$ the submodule of $A_A$ spanned by all paths of length at least 1.
\end{example}

%\begin{example}
%This example shows that we really need composition series on $A_A$ for things to be well-behaved.
%Let $A=\Bbbk[x]$.  Each irreducible polynomial $f$ generates a maximal ideal $(f)\subset \Bbbk[x]$ and so $\rad(A)\subset \bigcap_{f:\text{ irred.}}(f)$.
%Note that there are infinitely many irreducible polynomials in $\Bbbk[x]$.
%
%We claim that $\rad(A)=0$.
%If, on the contrary, there is some non-zero $g$ in this intersection of ideals, then all irreducible polynomials are factors of $g$; this is a contradiction as $g$ can only has finite degree, i.e. finitely many irreducible factors.
%\end{example}

\begin{proposition}\label{prop:radical}
Suppose $A_A$ has a composition series.
Then the following holds for the Jacobson radical $\rad(A)$.
\begin{itemize}
\item $\rad(A)$ is the intersection of finitely many maximal right ideals.
\item $\rad(A)$ is the intersection of all two-sided ideals $\Ann_A(S):=\{a\in A\mid ma=0\forall m\in S\}$, in other words
\[\rad(A)=\{a\in A\mid Sa=0\text{ for all simple }S\}.\]
\item $\rad(A)$ is a two-sided ideal of $A$.
\item $\rad(A)^\ell=0$ for $\ell$ at most the length of $A_A$. 
\item $(A/\rad(A))_{A/\rad(A)}$ is a semisimple (as a module).
%\item If $I\subsetneqq A$ is a strict two-sided ideal of $A$ and $(A/I)_{A/I}$ is semisimple, then $\rad(A)\subset I$.
\item $A_A$ is a semisimple (as a module) if, and only if, $\rad(A)=0$ (i.e. $A$ semisimple as an algebra).
\end{itemize} 
\end{proposition}

Proof omitted.  We note that all of these claims do make use of the Jordan-H\"{o}lder theorem.


\begin{example}
\begin{enumerate}
\item Direct product of two semisimple algebras is semisimple.

\item $A=\Mat_n(D)$ with $D$ a division $\Bbbk$-algebra is a semisimple $\Bbbk$-algebra.  We have decomposition $A_A\cong V^{\oplus n}$ into $n$ copies of $n$-dimensional simple module \[
V=\{(v_i)_{1\leq i\leq n}\mid v_i\in D\;\forall i\}.\]
%\vspace*{-1cm}


\item $A:=\Bbbk[x]/(x^n)$ is not semisimple for any $n\geq 2$ as it has a non-trivial (unique) maximal ideal $\rad(A)=(x)$.

\end{enumerate}
\end{example}


%\begin{exercise*}
%Recall (or check any reference book) the notion of \defn{free module}\index{module!free} and the \defn{rank}\index{rank} of it.
%Check that for an idempotent $e\in A$, $Ae$ is a direct summand of $A$.  In ring/module theory terms, (by definition) $Ae$ is thus a \defn{projective module}\index{projective} since it is a direct summand of a free module.
%\end{exercise*}


\begin{theorem}[Artin-Wedderburn theorem]\index{Artin-Wedderburn theorem}
Let $A$ be a finite-dimensional $\Bbbk$-algebra and let $r$ be the number of isoclasses of simple $A$-modules, say, with representatives $S_1,\ldots, S_r$.
Let $D_i:=\End_A(S_i)$ be the division $\Bbbk$-algebra given by endomorphism of the simple module $S_i$.
Then there is an isomorphism of $\Bbbk$-algebras
\[
A/\rad(A) \cong \Mat_{n_1}(D_1) \times \cdots \times \Mat_{n_r}(D_r).
\]
\end{theorem}
As before, if we work over algebraically closed field $\Bbbk=\overline{\Bbbk}$, then all the $D_i$'s are just $\Bbbk$.
\begin{proof}
Let $B:=A/\rad(A)$.
By definition of $\rad(A)$, the $A$-module $A/\rad(A)$ is semisimple, and any $A$-submodule $M$ of $A/\rad(A)$ satisfies $M\rad(A)=0$.  Hence, $M=M/M\rad(A)$ is naturally a $B$-module and $\End_B(M)\cong \End_A(M)$ (even as algebras!).

By Lemma \ref{lem:yoneda and idem trunc}, we have $B\cong \End_B(B)$.  Since $B$ is semisimple, the $B_B$ is a semisimple $B$-module, say, $B\cong S_1^{\oplus n_1}\oplus \cdots \oplus S_r^{\oplus n_r}$ where $S_i$ are the (representatives of the) isomorphism classes of simple $B$-modules.  Hence, it follows from Schur's lemma and its consequence (Lemma \ref{lem:Schur} and Lemma \ref{lem:endo ss}) that 
\[
B\cong \End_B(B) \cong \Mat_{n_1}(D_1)\times\cdots \times \Mat_{n_r}(D_r),
\]
where $D_i:=\End_B(S_i)$ for all $1\leq i\leq r$.  This completes the proof.
\end{proof}


\begin{corollary}\label{cor:simple corresp}
For any finite-dimensional $\Bbbk$-algebra $A$, let $\mathrm{Sim}(A)$ be the set of isomorphism-class representatives of simple $A$-modules.
Then there is a one-to-one correspondence 
\[
\xymatrix@R=0pt@C=5pt{
\mathrm{Sim}(A) \ar@{<->}[r]^(.4){1:1} & \mathrm{Sim}(A/\rad(A))\\
**[l]S\ar@{|->}[r] & **[r]{\overline{S}:=S/S\rad(A)}\\
& **[r]{\;\; (=S\text{ as underlying vector space})}\\
**[l]\mathrm{res}T &  T \ar@{|->}[l] }
\]
where $\mathrm{res}T$ is the restriction of $T$ along $A\twoheadrightarrow A/\rad(A)$.
\end{corollary}

\begin{example}
Suppose that $Q$ is finite acyclic, i.e. $\Bbbk Q$ is finite-dimensional.  Since $\rad (\Bbbk Q)$ is spanned by all non-trivial paths, $\Bbbk Q/\rad(\Bbbk Q)$ is just the semisimple $\Bbbk Q$-module $\bigoplus_{i\in Q_0} S_i$.  
In particular, the Artin-Wedderburn decomposition reads
\[
\Bbbk Q \cong \Bbbk \times \cdots \times \Bbbk
\]
with one copy of $\Bbbk$ for each $i\in Q_0$ on the right-hand side.
Moreover, every simple $\Bbbk Q$-module is isomorphic to one of $S_i$ for $i\in Q_0$.
\end{example}

\begin{exercise}
Show that when $Q$ is the Jordan quiver, then $\Bbbk Q$ has infinitely many simple modules and that $\rad(\Bbbk Q)=0$.
\end{exercise}


\pagebreak
%%%%%%%%%%%%%%%%%%%%%%%%%%%%%%%%%%%%%%%%%%%%%%%%%%%%%%%%%%%%%%%%%%%
%%%%%%%%%%%%%%%%%%%%%%%%%%%%%%%%%%%%%%%%%%%%%%%%%%%%%%%%%%%%%%%%%%%
%%%%%%%%%%%%%%%%%%%%%%%%%%%%%%%%%%%%%%%%%%%%%%%%%%%%%%%%%%%%%%%%%%%


\section{Radical and socle}

\begin{definition}
The \defn{radical}\index{radical!of module} of an $A$-module $M$ is $\rad(M):=M\rad(A)$.  In general, take $\rad^0(M):=M$ and denote by $\rad^{k+1}(M):=\rad(\rad^k(M))=\rad^k(M)\rad(A)$ for all $k\geq 0$.

Successively taking the radical yields a series:
\[
0\subset \rad^\ell(M) \subset \cdots \subset \rad(M) \subset M
\]
This is called the \defn{radical series}.
The quotient $M/\rad(M)$ is called the \defn{top}\index{top} of $M$, and is denoted by $\mathrm{top}(M)$.
\end{definition}

\begin{proposition}\label{prop:rad(M)}
The following hold for $M\in\mod A$.
\begin{enumerate}
\item $\rad(M)$ is the intersection of all maximal submodules of $M$.
\item $\mathrm{top}(M):=M/\rad(M)$ is the maximal semisimple quotient of $M$.
\item $\rad(M\oplus N) = \rad(M)\oplus\rad(N)$.

\item If $f:M\to N$ is a surjective $A$-module homomorphism, then $f(\rad M) = \rad N$.
%\item (Nakayama's Lemma, special case 1) $\rad(M)=M  \Rightarrow M=0$.
\item (Nakayama's Lemma, special case) For a submodule $N\subset M$, $(N+\rad(M)=M) \Rightarrow N=M$.
\end{enumerate}
\end{proposition}

Proof omitted; this follows the same kind of arguments as in the case for $\rad(A)$.

\begin{example}\label{eg:P/radP simple}
Let $A$ be a finite-dimensional algebra.
Suppose that $e$ is a primitive idempotent, i.e. $P:=eA$ is an indecomposable $A$-module.
Since $A=P\oplus Q$ (by taking $Q:=(1-e) A$), we have
\[ \rad(P)\oplus \rad(Q) = \rad(P\oplus Q) = \rad(A). \]
Since $P$ and $Q$ has no common (non-trivial) submodule, we get that
\[ A/\rad(A) = \frac{P\oplus Q}{\rad(P\oplus Q)} = P/\rad (P) \oplus \frac{Q}{\rad(Q)}. \]
Thus, it follows from Corollary \ref{cor:simple corresp} that $P/\rad (P)$ is a simple module and that every simple $A$-module arises this way.
In other words, let $\mathrm{PIM}(A)$ be the set of isoclass (=isomorphism class) representatives of indecomposable direct summands of $A$, then we have a correspondence
\begin{align}\label{eq:simple-proj corresp.}
\xymatrix@R=0pt{
\mathrm{PIM}(A)\ar@{<->}[r]^{1:1} & \mathrm{Sim}(A)\\
P\phantom{abc} \ar@{|->}[r] & {P/\rad(P)}
}
\end{align}
\end{example}


For a simple $A$-module $S$, denote by $P_S$ the corresponding direct summand $P$ of $A$ under the correspondence \eqref{eq:simple-proj corresp.}.
%
%\begin{definition}
%A \defn{projective} $A$-module is a direct summand of free $A$-module.  
%\end{definition}
%
%Let $\proj(A)$ be the collection of modules $P$ with decomposition
%\[
%P \cong \bigoplus_{S\in \mathrm{Sim}(A)} P_S^{\oplus m_S}
%\]
%for $m_S\ge 0$, i.e. this is the collection of all finitely generated projective $A$-modules.
%
%\begin{definition}
%For an $A$-module $M\in \mod A$, the \defn{projective cover} of $M$ is a surjective $A$-module homomorphism $p_M:P_M\twoheadrightarrow M$ such that for any surjective homomorphism $p':P'\to M$ with $P'\in \proj(A)$, there exists a surjective homomorphism $q:P'\to P_M$ such that $p'=p_M q$, i.e. there is a commutative diagram:
%\[
%\xymatrix{ & P' \ar@{->>}[d]^{\forall p'} \ar@{-->>}[ld]_{\exists q}\\
%P_M \ar@{->>}[r]^{p_M} & M. }
%\]
%The module $P_M$ is also called the projective cover of $M$.
%\end{definition}
%
%\begin{example}
%The projective cover of $S$ is the module $P_S$ in the correspondence \eqref{eq:simple-proj corresp.}.
%More generally, for a semisimple module $M$, the projective cover of $M$ is just the direct sum of $P_S$, one for each simple direct summand $S$ of $M$.
%\end{example}
%
%\begin{lemma}\label{lem:proj cover characterise}
%Projective cover exists for any finitely generated $A$-module $M$.  Moreover, a surjective homomorphism $p:P\to M$ with $P\in\proj(A)$ is a projective cover if it induces an isomorphism $P/\rad P \cong M/\rad M$.
%\end{lemma}
%\begin{proof}
%By definition, every finitely generated module $M\in \mod A$ admits a surjective $A$-module homomorphism $f:A^{\oplus m} \to M$.
%Thus, projective cover is just the minimal direct summand $P$ of $A^{\oplus m}$ such that $f|_P$ is surjective.
%
%Let $q:M\to M/\rad M$ be the canonical projection map.
%The latter claim follows from Proposition \ref{prop:rad(M)} that $f(\rad P)=\rad M$ by Proposition \ref{prop:rad(M)} and the composition $qf|_P:P\to M/\rad M$ being surjective.
%\end{proof}

There is a construction dual to $\rad(M)$.

\begin{definition}
The \defn{socle}\index{socle} of an $A$-module $M$ is $\soc(M)$, which is defined as the maximal semisimple submodule of $M$.
More generally, take $\soc^0(M)=0$ and for $k\geq 0$, let $\soc^{k+1}(M)$ to be the submodule of $M$ generated by the lift of $\soc(M/\soc^k(M)) \subset M/\soc^k(M)$.
This yields a series
\[
0\subset \soc(M) \subset \soc^2(M) \subset \cdots \subset \soc^\ell(M)=M
\]
called the \defn{socle series} of $M$.
\end{definition}

\begin{example}
Consider a path algebra $\Bbbk Q$ of a finite acyclic (for simplicity) quiver $Q$, and $x\in Q_0$.  The indecomposable injective $I_x=D(\Bbbk Q e_x)$ has a simple socle isomorphic to $S_x$.  Essentially this can be seen by a dual argument in showing $\top(P_x)\cong S_x$.
More generally, analogous to Example \ref{eg:P/radP simple}, for a finite-dimensional algebra $A$, every simple $A$-module appears as $\soc(I)$ for an indecomposable direct summand of $D(A)$.
\end{example}

%There is also a notion that is dual to projective modules and projective cover.  However, we need to assume $A$ to be a finite-dimensional algebra in order for this notion to be compatible with the use of the term in homological algebra.
%
%\begin{definition}
%Let $A$ be a finite-dimensional algebra.  
%An $A$-module $I$ is \defn{indecomposable injective} if it is an indecomposable direct summand of $DA$, i.e. $I\cong D(eA)$ for some primitive idempotent $e\in A$.
%
%An $A$-module $I$ is (finitely generated) \defn{injective} if it is isomorphic to a finite direct sum of indecomposable injective $A$-modules, i.e. $I\cong \bigoplus_{S\in \mathrm{Sim}(A)} I_S^{\oplus m_S}$ with $m_S\ge 0$, where $I_S= D(e_SA)$ with $e_S$ the primitive idempotent corresponding to simple $A$-module $S$.
%
%For an $A$-module $M\in \mod A$, the \defn{injective envelop} of $M$ is an injective $A$-module homomorphism $i_M:M\hookrightarrow I_M$ such that for any injective homomorphism $i':M\to I'$ with $I'\in \inj(A)$, there exists an injective homomorphism $j:I_M\to I'$ such that $i'=ji_M$.
%\end{definition}
%
%Analogously, we have the following.
%
%\begin{lemma}
%Injective envelop exists for any $M\in \mod A$.  Moreover, an injective homomorphism $i:M\to I$ with $I\in \inj A$ is an injective envelop if it induces $\soc(M)\cong \soc(I)$.
%\end{lemma}
%
%We mention the following fact just-in-case of further use.

%In general , the socle series and radical series of a module can be very different.  However, it turns out that they have the same length.

\begin{lemma}
For $M\in\mod A$, the socle series and radical series has the same length, and this length is called the \defn{Loewy length} of $M$, and is denoted by $\LL(M)$.
\end{lemma}
\begin{proof}
Let $r_M$ (resp. $s_M$) denotes the length of the radical (resp. socle) series of $M$.
First, we show that $s_M\le r_M$ by induction on $s_M$.
This is clearly fine if $s_M=0$.
%For $s_M>0$, consider $M$ and take a submodule $N$ whose socle series of length $s_M-1$.  
%Then it follows from the induction hypothesis that the radical series of $N$ is $r'\ge s-1$, and so we have $s-1 \le r' \le r$.

Suppose that $s_M>0$.  By definition we have $\rad^{r-1}(M)$ a semisimple submodule of $M$, and so $\rad^{r-1}(M)\subset \soc(M)$.
This means that there is a surjective homomorphism $M/\rad^{r-1}(M)\twoheadrightarrow M/\soc(M)$, and so  $r_{M/\rad^{r-1}(M)}\ge r_{M/\soc M}$ (EXERCISE!).  In particular, we have 
\[
r_M = r_{M/\rad^{r-1}(M)} + 1 \ge r_{M/\soc M} + 1.
\]
Since $s_{M/\soc M} = s_M-1$, it follows from the induction hypothesis that $s_{M/\soc M} \le r_{M/\soc M}$, and hence
\[
s_M = s_{M/\soc M} + 1 \le r_{M/\soc M} + 1 \le r_M,
\]
as required.

One can show that $r_M\le s_M$ dually.
\end{proof}

Note that the semisimple subquotients in (between the layers of) the socle series and the radical series of a module may not coincide.

\begin{example}
Let $Q$ be the quiver $1\xleftarrow{\alpha}2\xrightarrow{\beta}3\xrightarrow{\gamma}4$ and consider the projective $P_2$ which has the form
\[
\Bbbk \xleftarrow{1} \Bbbk \xrightarrow{1} \Bbbk \xrightarrow{1} \Bbbk
\]
Then we have radical series
\[
0\subset S_4=\Bbbk \beta\gamma \stackrel{S_1\oplus S_3}{\subset} \rad(P_2) = \Bbbk\alpha+\Bbbk\beta+\Bbbk\beta\gamma \stackrel{S_2}{\subset} P_2
\]
and socle series
\[
0\subset S_2\oplus S_4 = \Bbbk\alpha+\Bbbk\beta\gamma \stackrel{S_3}{\subset} \rad(P_2) \subset P_2.
\]
\end{example}


\pagebreak
%%%%%%%%%%%%%%%%%%%%%%%%%%%%%%%%%%%%%%%%%%%%%%%%%%%%%%%%%%%%%%%%%%%
%%%%%%%%%%%%%%%%%%%%%%%%%%%%%%%%%%%%%%%%%%%%%%%%%%%%%%%%%%%%%%%%%%%
%%%%%%%%%%%%%%%%%%%%%%%%%%%%%%%%%%%%%%%%%%%%%%%%%%%%%%%%%%%%%%%%%%%
%


\pagebreak
%%%%%%%%%%%%%%%%%%%%%%%%%%%%%%%%%%%%%%%%%%%%%%%%%%%%%%%%%%%%%%%%%%%
%%%%%%%%%%%%%%%%%%%%%%%%%%%%%%%%%%%%%%%%%%%%%%%%%%%%%%%%%%%%%%%%%%%
%%%%%%%%%%%%%%%%%%%%%%%%%%%%%%%%%%%%%%%%%%%%%%%%%%%%%%%%%%%%%%%%%%%

\section{Example: Topological data analysis}

Topological data analysis concerns the ``rough shape of data".
Here, we regard data as just a finite discrete set $X$ in $\R^d$ (with usual Euclidean metric if you like).  $X$ itself is not particular interesting space (in terms of geometry or topology) for further analysis; yet, we can often see ``pattern" -- whether they look more or less randomly distributed, whether they are distributed in the space in a way that avoid certain areas, etc.

A more well-known mathematical approach to addressing this issue is statistics, where we try to see if the pattern tells us correlation between different parameters.
For topological data analysis (TDA) we want to just tell if the data form some `shapes with holes' (this is where `topology' comes in).
The idea is to replace each data point $x\in X$ by a ball $B_t(x)$ of very small radius $t$, slowly increase the radius and observe how the topology (e.g. by looking at topological invariant such as the `genus') of the space $X_t := \bigcup_{x\in X} B_t(x)$ changes.

Note that if $s\le t$, then we have a subspace $X_s \subset X_t$.
Moreover, in practice, it makes sense to sample $t$ to a finite sequence $t_1 < t_2 < \cdots < t_n$ and take $X_i:=X_{t_i}$.  Since we only concern topology of $X_t$, we can replace $X_t$ by a simplicial complex  $\triangle_t$ where $0$-cells (points) are $x$, and $\{x_1,\ldots, x_r\}$ form an $r$-cells if $B_t(x_1)\cap \cdots \cap B_t(x_r)\neq \emptyset$.  Having a simplicial complex means that we can take (e.g. the $p$-th) homology group $H_p(X_t) = H_p(\triangle_t)$.  The fact that we have $X_{t_i}\subset X_{t_{i+1}}$ means that we have a chain
\[
H_p(X_1) \to H_p(X_2) \to \cdots \to H_p(X_n).
\]
If we linearise these abelian groups to $\Bbbk$-vector spaces, then we get a chain of vector spaces and linear transformations -- this is nothing but a representation of the $\vec{\bbA}_n$-quiver 
\[
1\to 2\to \cdots \to n.
\]
In TDA, such a chain is called \defn{persistence module} (of 1 parameter / finite linear poset).
Understanding the indecomposable decomposition of a persistence module is an important aspect in TDA, this can even be used to characterise the nature of the data set (e.g. one may record some data from various metals, and the topological information can be used to distinguish each metal just from the data set).

An \defn{interval module} $M_{[a,b]}$ for $1\le a \le b\le n$ is the $\vec{\bbA}_n$-quiver representation given by 
\[
0 \to \cdots \to 0\to \Bbbk \xrightarrow{\id} \Bbbk \xrightarrow{\id}  \cdots  \xrightarrow{\id}\Bbbk \to 0 \to\cdots 0 \]
where the non-zero space starts at $a$ and ends at $b$.
This is clearly an indecomposable representation.  In fact, forms \emph{all} indecomposable representation -- known by Gabriel in the 70s (this is one special case of the Gabriel's theorem).
The following is then just a consequence of Krull-Schmidt theorem, but turns out to be fundamental in TDA.

\begin{proposition}
Every persistence module can be decomposed uniquely to a direct sum of interval module.
\end{proposition}

\medskip
The above is what people call `single parameter', or (finite) `linear poset', case.
There are other possible forms:
\begin{enumerate}
\item Multi-parameter case:  the quiver $\vec{\bbA}_n$ is replaced by the `commutative cube', i.e. the bound quiver $\vec{\bbA}_{n_1} \times \cdots \times \vec{\bbA}_{n_r}$ with relation $\alpha\beta -\beta\alpha$ for arrows $\alpha,\beta$ going in different directions.  In other words, a persistence module in this case is the same as an $A$-module, where $A = \bigotimes_{i=1}^r \Bbbk \vec{\bbA}$.

\item Poset case: the quiver $\vec{\bbA}_n$ is replace by the bound quiver $(Q,I)$ whose underlying quiver $Q$ is the Hasse quiver of the poset $P$, and $I$ includes all commutation 
\[
x \xrightarrow{\alpha}  y \xrightarrow{\beta}z - x\xrightarrow{\alpha'} w \xrightarrow{\beta'} z
\]
whenever $x\ge y\ge z$ and $x\ge w\ge z$.  In other words, a persistence module in this case isthe same as a module over the \defn{incidence algebra} of the poset $P$.
\end{enumerate}
In these general cases, one can still define interval modules, but it is no longer true that every $A$-module can be decomposed into interval modules.  Much of the recent algebraic and homological aspect of TDA concerns how to overcome such a problem.


For other aspects and more in-depth study of applying quiver representation to TDA, see, for exmample, book of Steve Oudot.



\pagebreak
%%%%%%%%%%%%%%%%%%%%%%%%%%%%%%%%%%%%%%%%%%%%%%%%%%%%%%%%%%%%%%%%%%%
%%%%%%%%%%%%%%%%%%%%%%%%%%%%%%%%%%%%%%%%%%%%%%%%%%%%%%%%%%%%%%%%%%%
%%%%%%%%%%%%%%%%%%%%%%%%%%%%%%%%%%%%%%%%%%%%%%%%%%%%%%%%%%%%%%%%%%%

\section{Example: Linear matrix pencil}

A \defn{linear matrix pencil} is a matrix $A+\lambda B$ with $A,B\in M_{m\times n}(\Bbbk)$ and $\lambda$ being an indeterminant, i.e. $A+\lambda B \in M_{m\times n}(\Bbbk[k])$.  For simplicity, we just say `matrix pencil' and drop the adjective `linear'.
Matrix pencil is used in the study of the so-called \emph{generalised eigenvalue problem}, and has applications to various applied mathematics like control theory, differential algebraic equations, numerical linear algebra, etc.

% https://www.opt.mist.i.u-tokyo.ac.jp/~iwata/papers/KCF.pdf
% looks fairly readable.

Two matrix pencils $A+\lambda B$ and $A'+\lambda B'$ are \defn{strictly equivalent} if there are invertible matrices $P\in M_m(\Bbbk),Q\in M_n(\Bbbk)$ such that $A'+\lambda B' = P(A+\lambda B) Q$.  This is equivalent to $A'=PAQ$ and $B'=PBQ$.

For simplicity, let us specialise $\Bbbk$ to an algebraically closed field; this means that we can use Jordan canonical form $J_m(\alpha) \in M_m(\Bbbk)$.

%For $m\ge 1$, we define
%\[N_m:= I + \lambda J_m(0)\]
% H_m:=J_m(0)
% K_m:=H+\lambda I, \quad \text{and} \quad 

Let $\overline{H}_{m}$ be the $m\times (m+1)$-matrix given by removing the last row of $J_m(0)$, and $(I_m|0)$ be the $m\times (m+1)$-matrix given by adding a column of zero to the identity matrix $I_m$.
Define
\[
L_m := \lambda (I_m|0) + \overline{H}_m = \begin{pmatrix}
\lambda & 1& 0 & \cdots & 0 \\
0 &\lambda & 1 & \ddots & \vdots \\
\vdots & \ddots & \ddots & \ddots & 0 \\
0 & \cdots & 0 & \lambda & 1
\end{pmatrix} \in M_{m\times (m+1)}(\Bbbk [\lambda]).
\]

Similar to the Smith/Jordan canonical form, each matrix pencil of  is equivalent to one in \defn{Kronecker canonical form}.  

\begin{theorem}
Every matrix pencil is strictly equivalent to a block-diagonal matrix, where each block is of one of the following form:
\begin{enumerate}
\item $L_m$ or $L_m^{\mathrm{tr}}$ for some $m\ge 1$.

\item $I_m+\lambda J_m(0)$ or $J_m(\alpha)+\lambda I_m$, for some $m\ge 1$ and $\alpha\in \Bbbk$.

%\item `strictly regular block', which is strictly equivalent to $J_m(\alpha)+\lambda I$ for some $\alpha\in\Bbbk$ when $\Bbbk = \overline{\Bbbk}$.
\end{enumerate}
\end{theorem}

One way to prove this theorem is to observe the following.

\begin{proposition}
For any $m,n\in \Z_{\ge 0}$, there is a one-to-one correspondence between the strictly equivalent classes of $m\times n$-matrix pencil and the isomorphism classes of representations over the Kronecker quiver $K_2:=(\xymatrix@1{1\ar@<.5ex>[r]\ar@<-.5ex>[r]&2})$ with dimension vector $(n,m)$.
\end{proposition}
\begin{proof}
Exercise.
\end{proof}
%\begin{proof}
%Given a matrix pencil $A+\lambda B \in M_{m\times n}(\Bbbk[\lambda])$, we have a representation $M(A,B)$ over $K_2$ given by
%\[
%\xymatrix{ \Bbbk^n  \ar@<.5ex>[r]^{A} \ar@<-.5ex>[r]_{B} & \Bbbk^m}.
%\]
%Suppose that $A'+\lambda B' = P(A+\lambda B) Q$ is a strictly equivalence between $A+\lambda B$ and $A'+\lambda B'$.  Then we have
%\[
%\xymatrix{ \Bbbk^n  \ar@<.5ex>[r]^{A} \ar@<-.5ex>[r]_{B} \ar[d]_{P} & \Bbbk^m \ar[d]^{Q} \\
%\Bbbk^n  \ar@<.5ex>[r]^{A'} \ar@<-.5ex>[r]_{B'} & \Bbbk^m}
%\]
%where all squares commute and $P,Q$ are isomorphisms -- this is precisely the definition of an isomorphism $(P,Q) : M({A,B})\xrightarrow{\sim} M({A',B'})$.
%\end{proof}

Under this correspondence, the Kronecker canonical form (the blocks appearing in the block-diagonal form) corresponds to indecomposable representations of the Kronecker quiver.  
In quiver representation theory, such classification can be done using a Kac's theorem.

\begin{corollary}
Every indecomposable $\Bbbk K_2$-modules is isomorphic to one of the following.
\begin{enumerate}
\item Preinjective modules, which correspond to $L_m$ for $m\ge 1$.
\item Preprojective modules, which correspond to $L_m^{\mathrm{tr}}$ for $m\ge 1$.
\item Regular modules $R_m(x)$ for $x\in \mathbb{P}_{\Bbbk}^1$ and $m\ge 1$, where
\[
\begin{cases}
R_m(\alpha) \text{ corresponds to } J_m(\alpha)+\lambda I_m & \text{if $\alpha=[x:1]$;}\\
R_m(\infty) \text{ corresponds to } I_m+\lambda J_m(0)& \text{if $\alpha=[1:0]=\infty$.}
\end{cases}
\]
\end{enumerate}
\end{corollary}
%\begin{proof}
%The preinjective modules correspond to $L_m$, the preprojective modules corerspond to $L_m^{\mathrm{tr}}$.  The regular modules correspond to $J_m(\alpha)+\lambda I_m$ for the $\alpha\in\Bbbk$ case and $I_m+\lambda J_m(0)$ for the $\infty$ case.
%\end{proof}

With this, various problems about linear matrix pencil can be transformed to problems about representations of the Kronecker quiver.  There are also `higher variation' of matrix pencils that correspond to the $n$-Kronecker quiver where there are $n\ge 2$ arrows between the 2 vertices (instead of just $n=2$).
Examples problem includes the ``matrix subpencil" problem, which are studied by Claus Ringel, Han Yang, \c{S}tefan \c{S}uteu-Sz\"{o}ll\H{o}si.

\begin{exercise}
Write down the indecomposable $\Bbbk K_2$-modules as representations.
\end{exercise}



\pagebreak
%%%%%%%%%%%%%%%%%%%%%%%%%%%%%%%%%%%%%%%%%%%%%%%%%%%%%%%%%%%%%%%%%%%
%%%%%%%%%%%%%%%%%%%%%%%%%%%%%%%%%%%%%%%%%%%%%%%%%%%%%%%%%%%%%%%%%%%
%%%%%%%%%%%%%%%%%%%%%%%%%%%%%%%%%%%%%%%%%%%%%%%%%%%%%%%%%%%%%%%%%%%

\section{Bounded path algebra}


For general quiver, we loses finite-dimensionality, and so many nice things we explained do not hold any more.  To retain finite-dimensionality, we need to consider nice quotients of path algebras.


\begin{definition}
An ideal $I\lhd \Bbbk Q$ is \defn{admissible}\index{ideal!admissible} if $(\Bbbk Q_1)^k \subset I\subset (\Bbbk Q_1)^2$ for some $k\geq 2$, i.e. $I$ is generated by linear combinations of paths of finite length at least 2.
The pair $(Q,I)$ is sometimes called \index{quiver!bounded}\defn{bounded quiver}.
A \defn{bounded path algebra} or \defn{quiver algebra} (with relations) is an algebra of the form $\Bbbk Q/I$ for some quiver $Q$ and admissible ideal $I$.
\end{definition}
\begin{remark}
Admissiblity ensures there is no redundant arrows (which appears if there is a relation like, for example, $\alpha-\beta\gamma\in I$ for some $\alpha\neq \beta,\gamma\in Q_1$) and there is enough vertices (trivial paths may not be primitive if there is a loop $x$ at a vertex with relation $x^2-x\in I$).
\end{remark}

\begin{lemma}
A bounded path algebra is finite-dimensional.
\end{lemma}
\begin{proof}
There exists a surjective algebra homomorphism $\Bbbk Q/(\Bbbk Q_1)^k \twoheadrightarrow \Bbbk Q/I$; the former is finite-dimensional.
\end{proof}

\begin{example}
Let $Q$ be the Jordan quiver with unique arrow $\alpha$.
Let $I$ be the ideal of $\Bbbk Q$ generated by $\alpha^k$ for some $k\geq 2$.
Then $I$ is an admissible ideal and $\Bbbk Q/I \cong \Bbbk [x]/(x^k)$ is a \index{truncated polynomial ring}\defn{truncated polynomial ring}.
\end{example}

\begin{definition}
A \defn{representation}\index{representation!of bouned quiver} $M$ of a bounded quiver $(Q,I)$ is a representation $M=(M_i,M_\alpha)_{i,\alpha}$ of $Q$ such that $M_a=0$ for all $a\in I$; here $M_a:=\sum_{p}\lambda_pM_{p}$ for $a=\sum_p \lambda_p p$ written as a linear combinations of paths $p$.

A \defn{homomorphism}\index{homomorphism!of bounded quiver representations} $f:M\to N$ of representations of $(Q,I)$ is a collection of linear maps $f_i:M_i\to N_i$ that intertwines arrows' action.
\end{definition}

As before, representations are really just synonyms of modules.
\begin{lemma}
A representation of a bounded quiver $(Q,I)$ is equivalent to a $\Bbbk Q/I$-module, and homomorphisms between representations are equivalent to those between $\Bbbk Q/I$-modules.
\end{lemma}


We have seen that it is easy to write down the indecomposable decomposition of the free $\Bbbk Q$-module $\Bbbk Q_{\Bbbk Q}$, we would like such nice thing to carry over to bounded path algebras.

\begin{theorem}(Idempotent lifting)
If $I$ is a nilpotent ideal of $A$ (i.e. $I^n=0$ for some $n\geq 1$) and $\overline{e}=\overline{e}^2\in A/I$, then there is a \defn{lift} $e=e^2\in A$ of $\overline{e}$, i.e. $\overline{e}=e+I$.
\end{theorem}
Proof omitted.  


\begin{corollary}\label{cor:idem lifting}
Let $I$ be an nilpotent ideal in $A$.  Suppose that
\[
1_{A/I}=f_1+\cdots + f_n
\]
for $f_i\in A/I$ are primitive orthogonal idempotents.
Then we have
\[
1_A= e_1+\cdots e_n
\]
where each $e_i\in A$ is a primitive orthogonal idempotent that lifts $f_i$.
\end{corollary}
%\begin{proof}
%Define idempotents ${e_i}'$ inductively.
%
%Set $e_1'=1$.  For each $i>1$, take $e_i'$ as any lift of $f_i+\cdots +f_n$ in the ring $e_{i-1}' A e_{i-1}'$.  Then for any $j\geq i$, $e_j'$ is an idempotent in the ring , and so $e_i' e_j' = e_j' = e_j'e_i'$.
%
%Define $e_i := e_i' - e_{i+1}'$ and so we have $e_i+I=f_i$.  Now we need to check orthogonality.
%If $j>i$, then by using $e_{i+1}'e_j'=e_j'$ and $e_{i+1}'e_{j+1}'=e_{j+1}'$ we have 
%\[
%e_{i+1}' e_j = e_{i+1}'(e_j'-e_{j+1}') = e_{i+1}'e_j' - e_{i+1}'e_{j+1}' = e_j'-e_{j+1}' = e_j,\]
%and so \[
%e_i e_j = (e_i'-e_{i+1}') e_{i+1}' e_j = e_{i+1}'e_j - e_{i+1}'e_j = 0.
%\]
%By a dual argument we have $e_je_i=0$.
%\end{proof}

\begin{corollary}\label{cor:prim idem on bound path algebra}
Let $A$ be a bound path algebra.
The primitive orthogonal idempotents of $A$ are given by the trivial paths.
\end{corollary}
\begin{proof}
Apply Corollary \ref{cor:idem lifting} with $I=\rad A$.  Since $A/I=\Bbbk Q_0$ is semisimple and so we have primitive orthogonal idempotents given by the trivial paths.
\end{proof}

\begin{notation}
As in the case of path algebra, denote by $S_x$ or $S(x)$ the simple $\Bbbk Q/I$-module given by placing a one-dimensional vector space at vertex $x\in Q_0$ and zero everywhere else.

Similarly, denote by $P_x$ or $P(x)$ the indecomposable $\Bbbk Q/I$-module $e_x \Bbbk Q/I$.
Likewise, by $I_x$ or $I(x)$ the indecomposable $D((\Bbbk Q/I) e_x)$.
\end{notation}

\begin{proposition}
For a finite-dimensional quiver algebra $A=\Bbbk Q/I$, there is a decomposition of $A$-modules
\[
A_A = \bigoplus_{x\in Q_0} P_x, \text{ and } (DA)_A = \bigoplus_{x\in Q_0} I_x.
\]
Moreover, $\{ S_x\cong \top(P_x) \cong \soc(I_x)\mid x\in Q_0\}$ form the complete set of isoclasses representatives of simple $A$-modules.
\end{proposition}
\begin{proof}
%Each arrow $\alpha\in Q_1$ generates a maximal right ideal of $A$ with quotient $S_x$ for $x=s(\alpha)$.  So we have $A/\rad(A)\cong \Bbbk Q_0 = \prod_{x\in Q_0}\Bbbk$.  
%As primitive orthogonal decomposition of the identity element of $A$ lifts to that of the identity element of  $A/\rad(A)$ by Corollary \ref{cor:idem lifting}, we have $e_x$ primitive, and so $P_x$ and $I_x$ are indecomposable.
By Corollary \ref{cor:prim idem on bound path algebra}, the trivial path $e_x$ is a primitive idempotent of $A$, and so $P_x = e_x A$ and $I_x=D(A e_x)$ are indecomposable.

The simple $A$-modules (up to isomorphisms) correspond to those over the semisimple quotient algebra $A/\rad(A)$ by Corollary \ref{cor:simple corresp}.  Hence, there are precisely $|Q_0|$ simple modules (up to isomorphism), given by the simple top of $P_x$, which is also isomorphic to the simple socle of $I_x$.
\end{proof}

We give a brief justification of why quiver representations provide a good way to construct lots of algebras.

\begin{theorem}
Suppose $\Bbbk$ is algebraically closed.
Then every finite-dimensional $\Bbbk$-algebra $A$ is \defn{Morita equivalent} to a bounded path algebra $\Bbbk Q/I$.
More precisely, $\Bbbk Q/I$ is given by $\End_A(\bigoplus_{e} eA)$ where $e$ varies over the set of representative of equivalence classes of primitive idempotents of $A$.
\end{theorem}

We defer the precise meaning of Morita equivalent; it roughly translates to saying that understanding $A$-modules and homomorphisms between them is equivalently (but not necessarily `equal to') to understanding modules and homomorphisms between a Morita equivalent bounded path algebra.

\begin{example}
Let $A=\Mat_n(\Bbbk)$ be a matrix ring.
Then the elementary matrix $e:=E_{1,1}$ is a primitive idempotent and $eA\cong E_{j,j}A$ for all $1\leq j\leq n$.
So $A$ is Morita equivalent to $\Bbbk \cong \Bbbk Q \cong \End_A(eA)$ where $Q$ is a one-vertex-no-arrow quiver.
\end{example}

Primitive idempotent decomposition, say, $1=\sum_{i=1}^ne_i$, allows us to write an algebra $A$ in matrix form $( e_iAe_j )_{1\leq i,j\leq n}$, where the `row spaces' form the indecomposable direct summands $e_iA$ and the dual of the `column space' form the indecomposable direct summands $D(Ae_i)$.
It could be a helpful mental exercise to think about the meaning of $eAe\cong \End_A(eA)$ from Yoneda lemma - this maybe a useful idea to keep in mind when one tries to understand the above theorem.

%Now we apply the above corollary to $I=\rad(A)$.
%We usually use the following convention of notation: \[
%A/\rad(A) = S_1\oplus \cdots S_t\]
%for the decomposition corresponding to idempotent decomposition $1$ in the semisimple algebra $A/\rad(A)$.  
%Note that different $S_i$ can be isomorphic here.
%Then by idempotent lifting we have idempotent decomposition $1=e_1+\cdots +e_t$ and indecomposable projective $P_i := Ae_i$.
%
%\begin{lemma}
%We have 
%$\Hom_A(P_i,S_j)\cong \begin{cases}
%\End_A(S_i), \text{if }S_i\cong S_j;\\
%0, \text{otherwise}.
%\end{cases}$
%\end{lemma}
%\begin{proof}
%If non-zero homomorphism $\theta:P_i\rightarrow S_j$ then $P_i/\ker\theta$ is a non-trivial submodule of $S_j$ and so by simplicity of $S_j$ we have $P_i/\ker\theta \cong S_j$ itself.
%By Corollary \ref{cor:idem lifting}, we have $P_i/\rad(A)P_i \cong S_i$. As $P_i/\ker\theta$ surjects onto $P_i/\rad(A)P_i = \cong S_i$, we have $S_i\isom S_j$ and $\theta$ lifts to an endomorphism of $S_i$.
%\end{proof}
%
%\begin{lemma}
%Suppose $K$ is algebraically closed.
%For any $M\in \mod A$, we have $\dim_K \Hom_A(P_i, M) = [M:S_i] := $ number of composition factors of $M$ that is isomorphic to $S_i$.
%\end{lemma}
%\begin{proof}
%Consider a Jordan-H\"{o}lder filtration $M\supset M_1\supset \cdots M_\ell\supset 0$.  $P_i$ being projective implies that $\Hom_A(P_i,M_j/M_{j+1})\cong \Hom_A(P_i, M_j)/\Hom_A(P_i,M_{j+1})$ by Remark \ref{rem:exact hom}.
%
%Note that $M_j/M_{j+1}$ is simple, and algebraically closed implies that $\End_A(S_i)\cong K$, so inductively applying the previous lemma yields the claim.
%\end{proof}


\iffalse
\pagebreak
%%%%%%%%%%%%%%%%%%%%%%%%%%%%%%%%%%%%%%%%%%%%%%%%%%%%%%%%%%%%%%%%%%%
%%%%%%%%%%%%%%%%%%%%%%%%%%%%%%%%%%%%%%%%%%%%%%%%%%%%%%%%%%%%%%%%%%%
%%%%%%%%%%%%%%%%%%%%%%%%%%%%%%%%%%%%%%%%%%%%%%%%%%%%%%%%%%%%%%%%%%%

{\bf Module diagram}

It is convenient to display the structure of a module via a more combinatorial form (a diagram) -- if possible. \footnote{There is no widely agreed name to these diagrams; for convenience, we just call them `module diagram' in this notes.}
This is (as of today technology) a better way to display module structure -- at least compare to composition series, or lattice diagram of the submodule lattice, or even, quiver representations, in some cases. 

%We remark first that it is not always possible for a module to have a `sufficient meaningful diagram', but when it does, calculations concerning the modules usually become much easier.

\begin{definition}
%Let $G,Q$ be two quivers.
%A \defn{quiver morphism} $\theta: G\to Q$ is a pair $(\theta_0:G_0\to Q_0, \theta_1:G_1\to Q_1)$ that intertwines with the source and target functions, i.e. $s_Q(\theta_1(\alpha))=\theta_0(s_G(\alpha))$ and likewise for $s_G$ replaced by $t_G$.
%A \defn{$Q$-coloured (directed) graph} is just a quiver morphism from a directed graph (i.e. quiver) $G$ to $Q$.

%A \defn{$\Bbbk Q$-coloured graph} is a quiver morphism $\theta: G\to Q$ such that there is no linear combination $\sum_{p:\text{ path in $G$}}\lambda_p p$ in $\Bbbk G$ such that $\sum_p\lambda_p\theta(p)\in I$, and there is no $Q$-coloured subgraph of the form
%\[
%\xymatrix@R=10pt@C=18pt{x \ar[rd]^a& & x\ar[ld]_{a} \\ & y & .}
%\]
Let $M\in \mod A$ be a finite-dimensional $A$-module for $A=\Bbbk Q/I$ a bounded path algebra.
The \defn{module diagram} is a (directed) graph with vertices labelled by composition factors of $M$ (in particular, there are $\dim_{\Bbbk} Me_x$ many vertices labelled by $x$), and arrows labelled by those in $Q_1$ in such a way that $x\xrightarrow{a}y$ if for an arrow $a\in Q$ that sends (the lift of) an element in the composition factor at $x$ to (the lift of) an element in the composition factor at $y$.
\end{definition}

Module diagram drawn in this way is not invariant under isomorphism.  A connected diagram may not even implies indecomposability in general (c.f. Homework 2).  Nevertheless, when the algebras or modules are well-behaved, then these diagram provide a very efficient combinatorial way to perform a lot of (linear algebra) calculation.

It is customary to draw the the module diagram flowing from top to bottom; in particular, the top (semisimple quotient) of $M$ is placed on the top of the diagram and the socle at the bottom.
We may omit a connecting line if there is no ambiguity.

\begin{example}
The indecomposable $U_{i,j}$ of $\Bbbk \linA_n$ is just a column of numbers labelled from $i$ down to $j$.  For a concrete example, the module diagram of $U_{4,6}$ is just $\begin{smallmatrix}
4\\5\\6\end{smallmatrix}$.
\end{example}

%
%Another convenience brought by module diagram is homological calculation.  We highlight one here.  Recall that a module is \emph{finitely generated} if there is some $n>0$ such that $A^{\oplus} \twoheadrightarrow M$.
%Note that the kernel of such a map is usually quite large -- many of the direct summands of $A$ will appear.  One is often interests in a more optimised version.
%
%\begin{definition}
%A surjective homomorphism $\pi_M: P_M\twoheadrightarrow M$ from a projective $A$-module $P_M$ is called the \defn{projective cover}\index{projective!cover} of $M$ if no indecomposable direct summand of $P_M$ is in the kernel of $\pi_M$.  By abuse of terminology, we often say `$P_M$ is the projective cover of $M$' to mean the existence of $\pi_M$.
%
%The \defn{syzygy} of $M$, denoted by $\Omega(M)$, is the kernel of the projective cover of $M$.
%\end{definition}
%
%Finding projective cover is extremely easy with module diagram, namely, the projective module covering $M(w)$ is given by $\bigoplus_{x} P_x$ where $x$ varies over all peaks of the module diagram of $M(w)$ (counted \emph{with} multiplicity), and the map is given by mapping $e_x$ to the corresponding element in the canonical basis of $M(w)$.
%
%\begin{example}
%We continue with the previous Example \ref{eg:module dgm}.
%The module $M(w)$ have projective cover $P_M = P_1\oplus P_5$ with $e_1 \mapsto v_6$ and $e_5\mapsto v_1$; or flexing out the full details:
%\[
%\left\{\begin{array}{rcl}
%P_1 &\to & M(w) \\
%e_1 &\mapsto& v_6 \\
%(2|1) &\mapsto& v_5 \\
%(4|1) &\mapsto& v_4 \\
%(3|1) &\mapsto& v_3; \\
%\text{everything else} & \mapsto & 0, 
%\end{array} \right. \quad \text{ and } \quad 
%\left\{\begin{array}{rcl}
%P_5 & \to & M(w)\\
%(3|6) &\mapsto& v_3 \\
%(5|6) &\mapsto& v_2 \\
%e_5 &\mapsto& v_1 \\
%(5|5)_v &\mapsto& v_0 \\
%\text{everything else} & \mapsto & 0, 
%\end{array}\right.
%\]
%We can visualise this map by highlighting the involved parts in the module diagram.  
%\begin{center}
%\begin{tikzpicture}
%
%\node at (0,0) {$\begin{matrix}
%1 \\ 2 \\ 3\\ 1 \end{matrix}$};
%\fill[blue!50,opacity=0.5,rounded corners={2pt}]  (-0.2,0.9) rectangle (0.2,-0.4);
%
%\begin{scope}
%\node[inner sep=1pt] (v2) at (1.5,0) {$\begin{matrix}6 \\3\\ 4 \end{matrix}$};
%\node (v1) at (2,1) {$5$};
%\node (v3) at (2.5,0) {$5$};
%\node (v4) at (2,-1) {$5$};
%\draw  (v1) edge (v2.north east);
%\draw  (v1) edge (v3);
%\draw  (v3) edge (v4);
%\draw  (v2.south east) edge (v4);
%
%\fill[yellow,opacity=0.5,rounded corners={2pt}] (1.2,-0.2) -- (1.4,0.8) -- (1.8,1.2) -- (2.2,1.2) -- (2.8,-0.2) -- cycle ;
%\end{scope}
%
%\begin{scope}
%\node at (4,0.5) {$1$};
%\node at (4.4,0.1) {$2$};
%\node at (4.8,-0.3) {$3$};
%\node at (5.2,0.1) {$6$};
%\node at (5.6,0.5) {$5$};
%\node at (6,0.1) {$5$};
%\end{scope}
%
%\node at (5,2) {$M(w)$};
%\node at (0,2) {$P_1$};
%\node at (2,2) {$P_5$};
%
%\fill[blue!50,opacity=0.5,rounded corners={2pt}]  (3.6,0.7) -- (4.2,0.7) -- (5.4,-0.5) -- (4.6,-0.5) -- cycle;
%\fill[yellow,opacity=0.5,rounded corners={2pt},shift={(3.4,-0.4)}] (1,-0.2) --  (2,1.1) -- (2.4,1.1) -- (3.2,-0.2) -- cycle ;
%\draw[->]  plot[smooth, tension=.7] coordinates {(0.3,0.7) (0.8,0.9) (1.4,1.6) (2.8,1.6) (3.8,0.8)};
%\draw[->]  plot[smooth, tension=.7] coordinates {(2.4,1.1) (4.8,1.4) (5.55,.9)};
%\end{tikzpicture}
%\end{center}
%Then it is easy to see what the kernel is -- the parts that are not highlighted, and the part with overlapping highlight, which is the (space spanned by) $(3|1)-(3|6)$ in this example.
%Now, the module diagram of $\Omega(M(w))$ can be obtained by gluing the unhighlighted parts in $P_1, P_5$ along the overlapping highlight parts.
%So in this example we have
%\[
%\Omega(M(w)) = \begin{smallmatrix}
%&3& & \\
%1&&4& \\
%& & &5
%\end{smallmatrix}.
%\]
%\end{example}
%
%An interesting example is what is called the \defn{Green's walk around Brauer tree} discovered by J. A. Green in the finite group representation context.  
%In modern setting, this is a special $\Omega$-orbit.
%We demonstrate an example of this in the following.
%
%\begin{example}
%We continue with the Brauer tree used in Example \ref{eg:Br tree 6 edges}.
%Let us start with the simple $S_1$ (=module diagram $1$).
%The projective cover is $P_1=\begin{smallmatrix}1\\2\\3\\1\end{smallmatrix}$, with kernel $\rad P_1 = \begin{smallmatrix}2\\3\\1\end{smallmatrix}$.
%
%This diagram has peak $2$, so it has projective cover $P_2 = \begin{smallmatrix}2\\3\\1\\2\end{smallmatrix}$ with kernel $S_2$.
%The next syzygy is then $\rad P_2 = \begin{smallmatrix}3\\1\\2\end{smallmatrix}$, and the next is $\begin{smallmatrix}
%4 \\5 \\6
%\end{smallmatrix}$, etc.  The full $\Omega$-orbit is then
%\[
%1,\; \begin{smallmatrix}2\\3\\1\end{smallmatrix},\; 2,\; \begin{smallmatrix}3\\1\\2\end{smallmatrix},\; \begin{smallmatrix}4\\5\\6\end{smallmatrix},\; 4,\; 
%\begin{smallmatrix}5\\6\\3\\4\end{smallmatrix},\; 
%\begin{smallmatrix}5\\5\end{smallmatrix},\; 
%\begin{smallmatrix}6\\3\\4\\5\end{smallmatrix},\; 6,\;
%\begin{smallmatrix}3\\4\\5\\6\end{smallmatrix},\; 
%\begin{smallmatrix}1\\2\\3\end{smallmatrix}, 1,\; \ldots 
%\]
%\end{example}

\begin{example}
Consider the following bounded quiver:
\[
Q: \xymatrix@1{ 1 \ar@/^/[r]^{\alpha_1} & 2 \ar@/^/[r]^{\alpha_2}\ar@/^/[l]^{\beta_1} & 3 \ar@/^/[l]^{\beta_2} },  \qquad I=(\alpha_1\alpha_2, \beta_1\beta_2, \beta_1\alpha_1-\alpha_2\beta_2).
\]
Then we have
\[
P_1 = \begin{smallmatrix}
\Bbbk e_1\\\Bbbk \alpha_1\\ \Bbbk \alpha_1\beta_1
\end{smallmatrix} = \begin{smallmatrix}
1 \\ 2 \\ 1
\end{smallmatrix} \quad,\quad 
P_2 = \begin{smallmatrix}
\Bbbk e_1\\\Bbbk \alpha_2 \oplus\Bbbk \beta_1\\ \Bbbk \alpha_2\beta_2
\end{smallmatrix} = \begin{smallmatrix}
2 \\ 1 \,\, 3 \\ 2
\end{smallmatrix} \quad,\quad 
P_3 = \begin{smallmatrix}
\Bbbk e_3\\\Bbbk \beta_2\\ \Bbbk \beta_2\alpha_2
\end{smallmatrix} = \begin{smallmatrix}
3 \\ 2 \\ 3
\end{smallmatrix}
\]

Let us consider the two-sided ideal $Ae_1A$.
This is spanned by all paths (`up to $I$') that passes through the vertex $1$.
As a right module, we can find its manifestation in the module diagram by picking everything below any appearance of the label $1$ -- in this case, it is all of $P_1$ and the $\begin{smallmatrix}1\\2\end{smallmatrix}$ part submodule of $P_2$.
In particular, the quotient algebra $(A/Ae_1A)_{A/Ae_1A}$ has module diagram:
\[
P_2^{A/Ae_1A} = e_2A/e_2Ae_1A = P_2/P_2e_1A = \begin{smallmatrix}
2 \\ 3 \end{smallmatrix} \quad,\quad 
P_3^{A/Ae_1A} = P_3/P_3e_1A = P_3^A = \begin{smallmatrix}
3 \\ 2 \\ 3
\end{smallmatrix}
\]
The bounded quiver presentation of $A/Ae_1A$ is given by 
\[
Q: \xymatrix@1{ 2 \ar@/^/[r]^{\alpha} & 3 \ar@/^/[l]^{\beta} },  \qquad I=(\alpha\beta).
\]

On the other hand, for $eAe$ with $e=e_2+e_3$, the module diagram is given by removing all composition factors that are not $S_2, S_3$, i.e.
\[
e_2Ae = \begin{smallmatrix}
2 \\ 3 \\2 \end{smallmatrix} \quad,\quad 
e_3Ae = \begin{smallmatrix}
3 \\ 2 \\ 3
\end{smallmatrix}
\]
and the bounded quiver presentation of $eAe$ is given by 
\[
Q: \xymatrix@1{ 2 \ar@/^/[r]^{\alpha} & 3 \ar@/^/[l]^{\beta} },  \qquad I=(\alpha\beta\alpha, \beta\alpha\beta).
\]
\end{example}
\fi

\pagebreak 
%%%%%%%%%%%%%%%%%%%%%%%%%%%%%%%%%%%%%%%%%%%%%%%%%%%%%%%%%%%%%%%%%%%%%%%%%%%%%%%
%%%%%%%%%%%%%%%%%%%%%%%%%%%%%%%%%%%%%%%%%%%%%%%%%%%%%%%%%%%%%%%%%%%%%%%%%%%%%%%
%%%%%%%%%%%%%%%%%%%%%%%%%%%%%%%%%%%%%%%%%%%%%%%%%%%%%%%%%%%%%%%%%%%%%%%%%%%%%%%
%%%%%%%%%%%%%%%%%%%%%%%%%%%%%%%%%%%%%%%%%%%%%%%%%%%%%%%%%%%%%%%%%%%%%%%%%%%%%%%

\section{Some basic category theory}

We briefly recall some language from category theory that are commonly used in representation theory.

\medskip

A \defn{category} $\CC=(\mathrm{ob}\CC,\hom\CC,\circ)$ consists of the following data.
\begin{enumerate}
\item A collection $\mathrm{ob}\CC$ of \defn{objects}.  It is common to write $X\in \CC$ instead of $X\in \mathrm{ob}\CC$.

\item For any pair $X,Y\in \CC$, there is a collection of \defn{morphisms from $X$ to $Y$}.  Such an element is often written $f:X\to Y$ if context clear, or $f\in \CC(X,Y)$ or $f\in\Hom_\CC(X,Y)$.

\item A binary operation $-\circ -: \CC(Y,Z)\times \CC(X,Y)\to \CC(X,Z)$ that takes $(g:Y\to Z, f:X\to Y)$ to the  \defn{composition} $g\circ f : X\to Z$, and this operations satisfies the following:

\begin{itemize}
\item Composition is \defn{associative}, i.e. $h(gf)=(hg)f$ for all (meaningful) $h,g,f$.

\item There are \defn{identity morphisms} $\id_X$ that are left and right units with respect to composition, i.e. $f\id_X=f$ and $\id_Yf=f$ for all (meaningful) $f$.
\end{itemize}
\end{enumerate}
For simplicity, we assume always that $\CC(X,Y)$ are sets.

\begin{example}
Some common categories of interest in algebras.
\begin{enumerate}
\item The category $\mod A$ of \emph{finitely generated} $A$-modules.

\item The category $\Mod A$ of \emph{all} $A$-modules.

\item The category $\Ab$ of abelian groups and group homomorphisms.

\item Let $Q$ be a quiver and $R$ be a ring.  Then we have a \defn{$\Bbbk$-linear path category} whose objects are the vertices of $Q$, and morphisms are $\Bbbk$-linear combinations of paths (directed, and possibly trivial, walks) of $Q$.  Morphism compositions are induced by path concatenation in the same way as path algebra.  Indeed, what we are doing is essentially just viewing the ring $\Bbbk Q$ as a category.

\item The category $\mathsf{coh} X$ of coherent sheaves over a scheme $X$.
\end{enumerate}
\end{example}

Suppose we have categories $\CC$ and $\CC'$.
We say that $\CC'$ is a \defn{subcategory} of $\CC$ if the following are satisfied.
\begin{enumerate}
\item $\mathrm{ob}\CC'$ is a subcollection of $\mathrm{ob}\CC$.
\item $\CC'(X,Y)\subseteq \CC(X,Y)$ for all $X,Y$.
\item Composition and identity morphisms in $\CC$ and in $\CC'$ coincide.
\end{enumerate}
If, moreover, $\CC'(X,Y)=\CC(X,Y)$, then we say that $\CC'$ is a \defn{full subcategory} of $\CC$.

\begin{example}
$\mod A$ is a full subcategory of $\Mod A$.
\end{example}

\begin{definition}
Let $\Bbbk$ be a field, a category $\CC$ is \defn{$\Bbbk$-linear} if $\CC(X,Y)$ are $\Bbbk$-vector spaces for all $X,Y\in\CC$ and compositions are $\Bbbk$-bilinear maps.  A functor $F:\CC\to\CD$ between $\Bbbk$-linear categories is \defn{$\Bbbk$-linear} if $F_{X,Y}:\CC(X,Y)\to \CD(F(X),F(Y))$ is $\Bbbk$-linear.
\end{definition}

\begin{example}
When $A$ is a $\Bbbk$-algebra.
Then the module category (finitely generated or not) is $\Bbbk$-linear.
\end{example}

A \defn{covariant functor} (resp. \defn{contravariant functor}) $F:\CC\to \CD$ consists of 
\begin{itemize}
\item an assignment of object $F(X)\in\CD$ for each $X\in\CC$, and
\item an assignment of morphism $F(f)\in\CD(F(X),F(Y))$ (resp. $F(f)\in\CD(F(Y),F(X))$) for each $f\in \CC(X,Y)$, such that
\item $F(\id_X)=\id_{F(X)}$, and
\item $F(gf)=F(g)F(f)$ (resp. $F(gf)=F(f)F(g)$).
\end{itemize}
Note that, a contravariant functor from $\CC$ to $\CD$ is the same as a convariant functor from the \defn{opposite category} $\CC^\op$ to $\CD$.  Here, $\CC^\op$ given by the same collection of objects as $\CC$, but the morphisms and their composition are in reverse direction, i.e. $\CC^\op(X,Y):=\CC(Y,X)$. 

%Functor allows us to change from the representation theory of one algebra to another.  The key point is that it preserves identity and compositions.

\begin{example}
The \defn{identity functor} $1_\CC:\CC\to \CC$ is the functor given by mapping every object and homomorphism to itself.
\end{example}

In practice (in representation theory and the like), a functor tells us how we can transform from the theory of modules over one algebra to those over another algebra.

\begin{example}
The ($\Bbbk$-linear) duality $D=\Hom_{\Bbbk}(-,\Bbbk):\mod A\to \mod A^\op$ is a contravariant functor.
\end{example}

To compare two functors (or compare how a pair of functors is close/far away from the identity functor), one uses \defn{natural transformations}.
More precisely, a natural transformation $\eta:F\Rightarrow G$ of functors $F,G:\CC\to \CD$ is a collection of morphisms $\eta_X:F(X)\to G(X)$ such that there is the following commutative diagram
\[
\xymatrix@C=35pt@R=25pt{ 
F(X) \ar[r]^{\eta_X} \ar[d]_{F(f)} & G(X) \ar[d]^{G(f)} \\
F(Y) \ar[r]^{\eta_Y} & G(Y) 
}
\]
If we say that a map $\eta_X:F(X)\to G(X)$ is \defn{natural in $X$}, then we mean that $\{\eta_X\}_{X\in \mod A}$ defines a natural transformation.

A \defn{natural isomorphism} is a natural transformation $\eta$ such that $\eta_X$ is an isomorphism for all $X$.  We simply write $F\cong G$ if there is a natural isomorphism between two functors $F,G$.

\begin{definition}
Let $\CC,\CD$ be two categories.
\begin{itemize}
\item $\CC$ and $\CD$ are \defn{equivalent} if there exists a pair of functors $F:\CC\to\CD, G:\CD\to\CD$ such that $1_\CC\cong GF$ and $1_{\CD}\cong GF$.

\item A functor $F:\CC\to\CD$ is \defn{dense} if, for all $D\in\CD$, there is $C\in\CC$ such that $F(C)\cong D$; i.e. ``surjective on object up to isomorphism".

\item A functor $F:\CC\to\CD$ is \defn{faithful} (resp. \defn{full}), if the induced map $F_{X,Y}:\CC(X,Y)\to \CD(F(X),F(Y))$ given by $f\mapsto F(f)$ is injective (resp. surjective).
\end{itemize}
\end{definition}

\begin{proposition}
$F:\CC\to\CD$ is an equivalence of categories if, and only if, $F$ is fully faithful dense.
\end{proposition}


\begin{exercise}
Recall that $R^\op$ is the opposite ring of $R$, whose underlying set is the same as that of $R$ with multiplication $(a\cdot^{\op}b):=b\cdot a$.
A \defn{representation}\index{representation} of $R$ is a ring homomorphism
\[
\rho: R^\op \to \End_\Z(M), \qquad  r \mapsto \rho_r,
\]
for some abelian group $(M,+)$.
A homomorphism $f:\rho_M\to \rho_N$ of representations $\rho_M:R^\op \to \End_\Z(M), \rho_N:R^\op \to \End_\Z(N)$ given by an abelian group homomorphism $f:M\to N$ that intertwines $R$-action, i.e. $\rho_N(r)\circ f=f\circ \rho_M(r)$ for all $r\in R$.

Eplain why a representation of $R$ is equivalent to a right $R$-module; and why homomorphisms correspond.
\end{exercise}

%\pagebreak
%%%%%%%%%%%%%%%%%%%%%%%%%%%%%%%%%%%%%%%%%%%%%%%%%%%%%%%%%%%%%%%%%%%
%%%%%%%%%%%%%%%%%%%%%%%%%%%%%%%%%%%%%%%%%%%%%%%%%%%%%%%%%%%%%%%%%%%
%%%%%%%%%%%%%%%%%%%%%%%%%%%%%%%%%%%%%%%%%%%%%%%%%%%%%%%%%%%%%%%%%%%

\pagebreak
\section{Extra: Additive and abelian categories}

In brief, additive categories are the categories where it make sense to talk about direct products and has the same backbone as the abelian groups.
Abelian categories are additive categories where it makes sense (on the categorical level) to talk about kernel and cokernel, and that these behave like what we expect in many typical situations like representation theory and algebraic geometry.

\begin{definition}
Consider an object $O\in \CC$ in a category $\CC$.
\begin{itemize}
\item $O$ is an \defn{initial object} if for all $X\in \CC$, there is a \emph{unique} morphism $O \to X$. 

\item $O$ is a \defn{terminal object} if for all $Y\in \CC$, there is a \emph{unique} morphism $Y\to O$.

\item $O$ is a \defn{zero object} if it is both initial and terminal.
\end{itemize}
\end{definition}
Note that, if it exists, then it is unique up to unique isomorphism.  For example, the zero module $0\in \mod A$ in the module category of a ring is a zero object.


\begin{definition}
Let $I$ be a(n indexing) set and $(X_i)_{i\in I}$ a family of objects in $\CC$.  
\begin{enumerate}
\item An object $P$ equipped with morphisms $(p_i:P\to X_i)_{i\in I}$ is a \defn{product} if for all $Z\in \CC$ and all $(f_i:Z\to X_i)_{i\in I}$, there is a unique morphism $f$ such that $p_if = f_i$ for all $i\in I$.

\item Dually, an object $C$ equipped with morphisms $(\iota_i:X_i\to C)_{i\in I}$ is a \defn{coproduct} if for all $Z\in \CC$ and all $(f_i:X_i\to Z)_{i\in I}$, there is a unique morphism $f$ such that $f\iota_i = f_i$ for all $i\in I$.

\item A \defn{biproduct} $B\in \CC$ is an object that is isomorphic to a product and also to a coproduct.
\end{enumerate}
\end{definition}
When a product/coproduct exist, it is unique up to unique isomorphism.  It is often written as $\prod_{i\in I} X_i$ for product, $\coprod_{i\in I} X_i$ for coproduct, and $\bigoplus_{i\in I} X_i$ for biproduct.

Product and coproduct are special instances of what-are-called limit and colimit respectively.  As a consequence of abstract nonsense (category theory), there are natural bijections
\[
\CC(X,\prod_{i\in I}Y_i) \cong \prod_{i\in I} \CC(X,Y_i), \quad \CC(\coprod_{i\in I}X_i,Y) \cong \prod_{i\in I} \CC(X_i,Y).
\]
Note that coproduct get `extract' out to a product here.
If the indexing set $I$ is finite, then all the $\prod$ and $\coprod$ in the above formula can be replaced by $\bigoplus$.

\begin{definition}
A category $\CC$ is \defn{additive} if the following are satisfied.
\begin{itemize}
\item It is \defn{pre-additive}, i.e. $\CC(X,Y)$ are abelian groups for all $X,Y\in\CC$ with bilinear composition.
\item Finite products $\prod_{i=1}^n X_i$ (in the categorical sense) exists for all $X_1,\ldots, X_n\in\CC$.
\item There is a zero object $0\in \CC$ (in the categorical sense).
\end{itemize}
A functor $F:\CC\to\CD$ between additive categories is \defn{additive} if $F(0)\cong 0$ and $F$ preserves (finite) products.
\end{definition}
In an additive category, we have (see the additive category page on nLab, or Kashiwara-Schapiro's book Corollary 8.2.4)
\[
\text{finite product = finite direct sum = finite coproduct.}
\]
Note that the terminology `direct sum' matches the one used in module theory.
Note also that being additive is a \emph{property} not a structure, since products and zero objects are uniquely determined by universal property, there is no choice involved; c.f. Kashiwara-Schapiro's book Theorem 8.2.14.

%\begin{exercise}
%An \defn{ideal} $\mathcal{I}$ of a category $\CC$ is a collection of morphisms $\mathcal{I}$ such that for all $f:X\to Y\in \mathcal{I}$ and all $g:Y\to Z,h:W\to X\in \CC$, we have $gfh:W\to Z \in \mathcal{I}$.  
%Suppose that $\CC$ is pre-additive.  Show that $\mathcal{I}(X,Y):=\mathcal{I}\cap \CC(X,Y)$ is a normal subgroup of $\CC(X,Y)$.  In particular, show that the quotient category $\CC/\mathcal{I}$, whose objects are the same as $\CC$ but morphism spaces are given by $(\CC/\mathcal{I})(X,Y):=\CC(X,Y)/\mathcal{I}(X,Y)$, is pre-additive.
%\end{exercise}

In a pre-additive category $\CC$, since Hom-sets are abelian groups, there is always a distinguished zero morphism $0\in \CC(X,Y)$ for any $X,Y\in\CC$.  This zero morphism allows us to formulate kernels and cokernels in the categorical level.

\begin{definition}
Let $\CC$ be a pre-additive category and $f:X\to Y$ a morphism in $\CC$.
\begin{enumerate}
\item A \defn{kernel} of $f$ is a morphism $i:K\to X$ such that $(K\xrightarrow{fi}Y)=0$ and for any $i':K'\to X$ with $fi'=0$, there exists a unique morphism $j:K'\to K$ such that $i'=ij$.
\[
\xymatrix@C=40pt{ & &  Y \\
K' \ar@/_1pc/[rrd]_{\forall\,i'} \ar@/^1pc/@{.>}[rru]|{0} \ar@{-->}[r]^{\exists! j} & K \ar[rd]^{i}\ar@{.>}[ru]|{0} & \\
& & X \ar[uu]^{f} \\ 
}
\]
In the case when the kernel of $f$ exists, then denote by $\Ker(f)$ the object $K$, and by $\ker(f)$ the morphism $i$, as appeared above.

\item Dually, a \defn{cokernel} of $f$ is a morphism $p:Y\to C$ such that $(X\xrightarrow{pf}C)=0$ and for any $p':Y\to C'$ with $p'f=0$, there exists a unique morphism $q:C\to C'$ such that $p'=qp$.
\[
\xymatrix@C=40pt{ Y \ar[rd]_{p} \ar@/^1pc/[rrd]^{\forall\,p'}& &  \\
& C \ar@{.>}[ld]|{0} \ar@{-->}[r]^{\exists! q} & C' \ar@/^1pc/@{.>}[lld]|{0} \\
X \ar[uu]^{f} & & \\ 
}
\]
In the case when the cokernel of $f$ exists, then denote by $\Cok(f)$ the object $C$, and by $\mathrm{cok}(f)$ the morphism $p$, as appeared above.

\item If $\ker(f)$ exists, then a \defn{coimage} of $f$ is a cokernel of $\ker(f)$.  When this exists, the cokernel object is denoted by $\mathrm{Coim}(f)$.

\item Dually, if $\mathrm{cok}(f)$ exists, then an \defn{image} of $f$ is a kernel of $\mathrm{cok}(f)$.  When this exists, the kernel object is denoted by $\mathrm{Im}(f)$.
\end{enumerate}
\end{definition}

\begin{lemma}\label{lem:ker cok im coim}
Let $f:X\to Y$ be a morphism in a pre-additive category $\CC$.
\begin{enumerate}
\item $\ker(f)$ is a \defn{monomorphism}, i.e. a left-cancellative morphism (meaning that, for any morphisms $g_1,g_2$, $\ker(f)\circ g_1 = \ker(f)\circ g_2 \Rightarrow g_1=g_2$).

\item $\mathrm{cok}(f)$ is an \defn{epimorphism}, i.e. a right-cancellative morphism (meaning that, for any morphism $g_1,g_2$, $g_1\circ \mathrm{cok}(f)=g_2\circ\mathrm{cok}(f) \Rightarrow g_1=g_2$).

\item If kernel, cokernel, coimage, image all exist for $f$, then $f$ can be factored uniquely as $X\to \mathrm{Coim}(f)\to \mathrm{Im}(f) \to Y$.
\end{enumerate}
\end{lemma}
\begin{proof}
(1), (2):  Follows from the universal property (definition).

(3) We have $(\Ker(f)\xrightarrow{\ker(f)}X\xrightarrow{f}Y)=0$. Hence, being a cokernel of $\ker(f)$, we have a unique morphism $\mathrm{Coim}(f)\to Y$ for which $f$ factors through.  

Since coimage is a cokernel, $X\to \mathrm{Coim}(f)$ is an epimorphism, which means that any morphism $g$ with $(X\to \mathrm{Coim}(f) \xrightarrow{g} Z)=0$ implies $g=0$.
Now we have
\begin{align*}
0 &= (X\xrightarrow{f} Y \xrightarrow{\mathrm{cok}(f)} \Cok(f)) \\
&= (X\to \mathrm{Coim}(f) \to  Y \xrightarrow{\mathrm{cok}(f)} \Cok(f)).
\end{align*} 
Hence, the morphism $(\mathrm{Coim}(f) \to Y \to \Cok(f))=0$.
On the other hand, $\mathrm{Im}(f)$ is the kernel of $Y\to \Cok(f)$, and so $(\mathrm{Coim}(f) \to Y \to \Cok(f))=0$ implies that $\mathrm{Coim}(f)\to Y$ factors uniquely through $\mathrm{Im}(f)\to Y$.  This finishes the construction of the desired morphism.
\end{proof}

\begin{exercise}
\begin{enumerate}
\item Show that in $\mod A$, injective (resp. surjective) module homomorphisms coincide with monomorphisms (resp. epimorphisms).
\item Show that in the category of rings (morphisms being ring homomorphisms), the natural embedding $\mathbb{Z}\to\mathbb{Q}$ is an epimorphism.
\end{enumerate}
\end{exercise}

\begin{definition}
A category $\CC$ is \defn{abelian} if
\begin{itemize}
\item $\CC$ is additive, and
\item any morphism $f:X\to Y$ admits a kernel $\ker(f):\Ker(f)\to X$ and a cokernel $\mathrm{cok}(f):Y\to \Cok(f)$, such that the induced morphism $\mathrm{Coim}(f):=\Cok(\ker(f))\to \Ker(\mathrm{cok}(f))=:\mathrm{Im}(f)$ of Lemma \ref{lem:ker cok im coim} (3) is an isomorphism.
\end{itemize}
\end{definition}


\begin{exercise}
$\mod A$ and $\proj A = \{\text{modules isomorphic to a finite direct sum of direct summands of $A$}\}$ are $\Bbbk$-linear additive categories, but only the former is abelian.
\end{exercise}

\begin{exercise}
Consider the category $\underline{\mod} A$ whose objects are the same as those in $\mod A$ (i.e. finitely generated $A$-modules), but with morphism space given by
\[\Hom_{\underline{\mod}A}(X,Y):=\Hom_A(X,Y)/I(X,Y),\]
where $I(X,Y)$ consists of all morphisms that factors through some $P\in \proj A$.
Show that
\begin{enumerate}
\item Any $P\in \proj A$ is isomorphic to $0$ in $\underline{\mod}A$.
\item $\underline{\mod}A$ is additive but not abelian.
\end{enumerate}
\end{exercise}

\begin{definition}
Let $f:X\to Y$ be a morphism in an abelian category $\mathcal{A}$.
\begin{itemize}
\item $f$ is said to be \defn{injective} if $\Ker(f)=0$.  In which case, we say that $X$ is a \defn{subobject} of $Y$.

\item $f$ is said to be \defn{surjective} if $\Cok(f)=0$.  In which case, we say that $Y$ is a \defn{quotient} (object) of $X$.
\end{itemize}
\end{definition}

\begin{lemma}
Let $f:X\to Y$ be a morphism in an abelian category $\mathcal{A}$.
\begin{enumerate}
\item $f$ is injective $\Leftrightarrow\; f$ is a monomorphism.
\item $f$ is surjetive $\Leftrightarrow\; f$ is an epimorphism.
\item $f$ is injective and surjective $\Leftrightarrow\; f$ is an isomorphism.
\end{enumerate}
\end{lemma}
\begin{proof}
\begin{align*}
\Ker(f)=0 &\Leftrightarrow \forall i:Z\to X, ( fi=0 \Rightarrow i \text{ factors through } 0 ) \\
&\Leftrightarrow \forall i:Z\to X, ( fi=0 \Rightarrow i=0 ) \\
%&\Leftrightarrow \Ker(\mathcal{A}(Z,X) \xrightarrow{f\circ -} \mathcal{A}(Z,Y)) = 0 \;\;\forall Z\in \mathcal{A}\\
& \Leftrightarrow f \text{ mono.}
\end{align*}
This proves (1); (2) can be proved similarly.

(3) $f$ isom implies both mono and epi, and so by (1) and (2) we get injective and surjective.  Conversely, when $f$ is both injective and surjective, then we have $X= \mathrm{Coim}(f)\cong \mathrm{Im}(f) = Y$.
\end{proof}

%TODO:
%L/R adjoint is R/L exact.


\pagebreak
%%%%%%%%%%%%%%%%%%%%%%%%%%%%%%%%%%%%%%%%%%%%%%%%%%%%%%%%%%%%%%%%%%%
%%%%%%%%%%%%%%%%%%%%%%%%%%%%%%%%%%%%%%%%%%%%%%%%%%%%%%%%%%%%%%%%%%%
%%%%%%%%%%%%%%%%%%%%%%%%%%%%%%%%%%%%%%%%%%%%%%%%%%%%%%%%%%%%%%%%%%%


\section{Bimodule, tensor and Hom}

The most typical method in constructing functors between module categories is through `tensoring' and `hom-ing', which we explain in details now.  For simplicity, we always work over a field $\Bbbk$.

\subsection{Bimodule}
\begin{definition}
Let $A,B$ be two $\Bbbk$-algebras.
An $A$-$B$-\defn{bimodule} is a $\Bbbk$-vector space $M$ that has the structure of a left $A$-module and also the structure of a right $B$-module, such that $(am)b=a(mb)$ for all $a\in A, b\in B, m\in M$.
In such a case, we may write $M\in A\mod B$ or ${}_AM_B$ to specify $M$ is an $A$-$B$-bimodule.

For simplicity, we assume all bimodules are \defn{$\Bbbk$-central}\index{central}, i.e. $\lambda m=m\lambda$ for all $\lambda\in\Bbbk$.  We will omit the adjective $\Bbbk$-central from now on.
\end{definition}

\begin{example}
For any algebra $A$, both $A$ and $D(A)$ are naturally an $A$-$A$-bimodule.
Note that the right/left module structure on $D(A)$ is induced by the left/right module structure on $A$.  (The direction of action has swapped!)
\end{example}

\begin{example}
$\Hom_A(X,Y)$ is naturally a $\End_A(Y)$-$\End_A(X)$-bimodule with action given by composition of homomorphisms.
\end{example}

%\begin{definition}
%Let $M,N$ be $\Z$-modules (i.e. abelian groups).  Then the \defn{tensor product}\index{tensor product} $M\otimes_\Z N$ is the abelian group given by the quotient $\Z(M\times N)/\sim$ of the free $\Z$-module with basis $M\times N$ (the Cartesian product of underlying sets) by the equivalence relation $\sim$.  Here, $\sim$ is generated by 
%\[(m+m',n) \sim (m,n)+(m',n)\quad \text{ and } \quad (m,n+n')\sim (m,n)+(m,n'),\]
%for all $m,m'\in M$ and all $n,n'\in N$.
%The image of $(m,n)\in \Z(M\times N)/\sim$ in the quotient $M\otimes_\Z N$ is denoted by $m\otimes n$.
%
%Let $R$ be a ring, $M$ a \emph{right} $R$-module and $N$ a \emph{left} $R$-module.  Then the tensor product $M\otimes_RN$ of $M$ and $N$ over $R$ is the $\Z$-module given by the quotient $M\otimes_\Z N/\sim'$, where $\sim'$ is the equivalence relation $(mr,n)\sim (m,rn)$ for all $m\in M, n\in N$.
%\end{definition}

\subsection{Tensor product}
\begin{definition}
Let $V,W$ be finite-dimensional $\Bbbk$-vector space with bases, say, $\CB, \CC$ respectively.
Then the \defn{tensor product}\index{tensor product!over a field} $V\otimes_{\Bbbk} W$ (or simplifies to $V\otimes W$ if context is clear) is the finite-dimensional $\Bbbk$-vector space with bases given by 
\[
\{ v\otimes w \mid v\in \CB, w\in \CC\}.
\]
%Let $A$ be a $K$-algebra, $M$ be a right $A$-module, and $N$ be a left $A$-module.  Then the \defn{tensor product}\index{tensor product!over an algebra} $M\otimes_A N$ of $M$ and $N$ over $A$ is the quotient $K$-vector space $M\otimes_K N / R$, where 
%\[
%R=\{ma\otimes n-m\otimes an\mid m\in M, a\in A, n\in N\}.
%\]
\end{definition}

In particular, note that $\dim_{\Bbbk}V\otimes W=(\dim_{\Bbbk}V)\times(\dim_{\Bbbk}W)$.
%WARNING: $M\otimes_A N$ is generally not an $A$-module - this is only the case when $A$ is .  We will say more about this later.

%Possibly the most important remark for people who see tensor product for the first time is that the elements of $V\otimes_\Z W$ are given by $\Z$-linear combinations of symbols $v\otimes w$ where $v\in V$ and $w\in W$.

%\begin{example}
%Show that $\Z/p\Z\otimes_\Z \Z/q\Z =0$ for relative primes $p,q$.
%\end{example}
%\begin{proof}
%We show that for any $a\in \Z/p\Z$ and $b\in \Z/q\Z$, $a\otimes b=0$ in $\Z/p\Z\otimes_\Z \Z/q\Z$.  Recall from Bezout's identity that we have integers $x,y\in \Z$ so that $px+qy=1$.  So we have
%\[
%a\otimes b = (a\cdot 1)\otimes b = a\cdot(px+qy)\otimes b = pax\otimes b+qay\otimes b = 0 + ay\otimes qb = 0
%\]
%This completes the proof.
%\end{proof}

%To prove things using tensor product, one often uses the characterisation by its \defn{universal property}.
%\begin{proposition}
%Suppose that $R$ is a ring, $M$ is a right $R$-module and $N$ is a left $R$-module.
%If $L$ is an abelian group and $f:M\times N \to L$ is a map of sets such that 
%\begin{align*}
%f(m+m',n) &= f(m,n)+f(m',n), \;\; \forall m,m'\in M \text{ and } n\in N;\\
%f(m,n+n') &= f(m,n)+f(m,n'), \;\; \forall m\in M \text{ and } n,n'\in N;\\
%f(mr,n) &= f(m,rn), \;\;\forall m\in N, r\in R, n\in N,
%\end{align*}
%then there exists a unique homomorphism $\hat{f}:M\otimes_R N\to L$ of abelian groups satisfying $\hat{f}(m\otimes n)=f(m,n)$.
%\end{proposition}
%One usually memorise this via the commutative diagram:
%\[
%\xymatrix@C=50pt{ M\times N \ar[r]^{{}^\forall f} \ar[d] & L \\ M\otimes_R N \ar@{-->}[ru]_{\exists!\,\hat{f}} & }
%\]
%\begin{proof}
%By definition, $M\otimes_R N = \Z(M\times N)/X$, where $X$ consists of all elements of the form $(m+m',n)-(m,n)-(m',n)$, $(m,n+n')-(m,n)-(m,n')$, and $(mr,n)-(m,rn)$. 
%
%Extend the map $f$ of sets by linearity (in $\Z$) yield a $\Z$-module homomorphism $\tilde{f}:\Z(M\times N)\to L$.
%By the assumption on $f$, we have $\tilde{f}(X)=0$, and so it induces the required homomorphism $\hat{f}$.  
%
%Since $\hat{f}(m\otimes n)$ is fixed by $f$ for all $m\in M$ and $n\in N$, and all elements in $M\otimes_R N$ are $\Z$-linear combinations of elements of the form $m\otimes n$.  The image of $\hat{f}$ for all elements are fixed by $f$.  This shows uniqueness of $\hat{f}$.
%\end{proof}

%\begin{definition}
%Let $R,S$ be rings.  An abelian group $M$ is an \defn{$R$-$S$-bimodule}\index{bimodule} if it is a left $R$-module and right $S$-module with commuting $R$- and $S$-action, i.e. $r(ms)=(rm)s$ for all $r\in R, m\in M, s\in S$.  In other words, it is a left module over $R\times S^\op$ (equivalently, right module over $S\times R^\op$).
%\end{definition}

%\begin{lemma}\label{lem:bimod on tensor Hom}
%Consider rings $R,S,T$.  Let $M$ be an $R$-$S$-bimodule, $N$ be an $S$-$T$-bimodule, and $L$ be an $R$-$T$-bimodule.
%\begin{enumerate}
%\item $M\otimes_S N$ is a $R$-$T$-bimodule given by $r\cdot(m\otimes n):=(rm)\times n$ and $(m\otimes n)\cdot t:=m\otimes(nt)$.
%
%\item $\Hom_R(L,M)$ is a $T$-$S$-bimodule given by $(t\cdot f)(l) := (f(lt))$ and $(f\cdot s)(l):=f(l)s$.
%\end{enumerate}
%\end{lemma}
%\begin{proof}
%Exercise.
%\end{proof}
%
%The above lemma tells us that tensor and Hom can be used to transfer modules (in fact, even homomorphisms) between different rings.
%Another consequence of Lemma \ref{lem:bimod on tensor Hom} is that, if $R$ is a commutative ring, then $R$-modules are the same as $R$-$R$-bimodules, and so $M\otimes_R N$ are automatically $R$-modules for $R$-modules $M$ and $N$.
%Similarly, as left (resp. right) $A$-modules over a $K$-algebra, say $A$, are really $A$-$K$-bimodules (resp. $K$-$A$-bimodules), and so tensoring over $R$ $M\otimes_A N$ is automatically a $K$-vector space.

%\begin{exercise}\label{ex:trivial tensor}
%Show that $R\otimes_R M\cong M$ for any left $R$-module $M$.
%\end{exercise}

%Now we can talk about the following result, which makes the structure of tensor products over a field much easier to understand.
%\begin{lemma}
%Let $\mathcal{B}, \mathcal{C}$ be bases of finite dimensional $K$-vector spaces $V, W$.
%Then $V\otimes_K W$ is a finite-dimensional $K$-vector space with basis $\{m\otimes n\mid m\in \mathcal{B}, n\in\mathcal{C}\}$.
%In particular, we have $\dim_K V\otimes_K W = (\dim_K V) (\dim_K W)$.
%\end{lemma}
%\begin{proof}
%By Exercise \ref{ex:trivial tensor}, we have $R^{\oplus n}\otimes_R M\cong M^{\oplus n}$.  The claim follows by taking $R=K$ and $M=W$.
%\end{proof}

%\begin{proposition}\label{prop:tensor decomp}
%There is a $\Bbbk$-vector space isomorphism
%\[
%\xymatrix@C=40pt@R=0pt{\left(\bigoplus_{i=1}^m M_i\right)\otimes _{\Bbbk} \left(\bigoplus_{j=1}^n N_j\right) 
%\ar[r]^(.55){\cong}& \bigoplus_{i,j}M_i\otimes_\Bbbk N_j \\
%u\otimes v \ar@{|->}[r]& (u_i\otimes v_j)_{i,j}
%}
%\]
%for $u=\sum_i u_i$ with $u_i\in M_i$ for all $i$, and $v=\sum_j v_j$ with $v_j\in N_j$ for all $j$.
%\end{proposition}

\begin{proposition}
Let $A, B$ be $\Bbbk$-algebras.  Then $A\otimes_\Bbbk B$ is also a $\Bbbk$-algebra with multiplication given by extending $(a\otimes b)(a'\otimes b')\mapsto aa'\otimes bb'$ linearly.  For $M\in \mod A$ and $N\in \mod B$, we have $M\otimes_\Bbbk N\in \mod A\otimes_\Bbbk B$.
\end{proposition}
\begin{proof}
Routine checking.
\end{proof}

%\begin{example}\label{eg:matrix tensor}
%Consider $A=(a_{i,j})_{1\leq i,j\leq m} \in \Mat_m(\Bbbk)$ and $B\in \Mat_n(\Bbbk)$ and defines (what is sometimes called Kronecker product of matrices)
%\[
%A\otimes B := \begin{pmatrix}
%a_{1,1} B & a_{1,2}B & \cdots & a_{1,m}B \\
%a_{2,1} B & \ddots &  & a_{2,m}B \\
%\vdots & & \ddots & \vdots \\
%a_{m,1} B & a_{m,2}B & \cdots & a_{m,m}B \\
%\end{pmatrix}.
%\]
%Then we have an isomorphism of algebras
%\[
%\Mat_m(\Bbbk)\otimes_\Bbbk \Mat_n(\Bbbk)\to \Mat_{mn}(\Bbbk), \quad (A,B)\mapsto A\otimes B.
%\]
%%From this, we can see that $(A\otimes B)^{-1} = A^{-1}\otimes B^{-1}$, if (and only if) both $A,B$ are invertible.  Thus, the isomorphism restricts to a group isomorphism $\GL(K^{\oplus m})\otimes_K \GL(K^{\oplus n})\cong \GL(K^{\oplus mn})$.
%\end{example}


\begin{lemma}
An idempotent $e\in A\otimes_{\Bbbk} B$ is primitive if and only if $e=e_l\otimes e_r$ for some primitive idempotents $e_l\in A$ and $e_r\in B$.
In particular, we have $\mathrm{Sim}(A\otimes_{\Bbbk}B) = \{S\otimes T\mid S\in\mathrm{Sim}(A), T\in\mathrm{Sim}(B)\}$.
\end{lemma}

Note that not all $A\otimes_{\Bbbk}B$-module is of the form $M\otimes N$.
\begin{example}
Let $A=\Bbbk[x]/(x^2)$ and $A':=\Bbbk[y]/(y^2)$.
Then $B:=A\otimes_{\Bbbk}A' = \Bbbk[x,y]/(x^2,y^2)$.
Then we have an indecomposable 2-dimensional $B$-module $V=\Bbbk u+\Bbbk v$ (top $S=B/\rad(B)$ and socle $S$) where both $x,y$ acts by $u\mapsto v$.
This cannot be of the form $M\otimes N$ for some $M\in\mod A$ and $N\in \mod A'$. Indeed, as both $x,y$ acts non-trivially, if $V=M\otimes N$ then both $M,N$ must have dimension at least 2, and so the tensor product has dimension at least 4; but $\dim_{\Bbbk}V=2$.
\end{example}

\begin{proposition}
An $A\otimes_{\Bbbk}B^\op$-module is the same as a ($\Bbbk$-central) $B$-$A$-bimodule.
Moreover, homomorphisms of $A\otimes B^\op$-modules correspond to ($\Bbbk$-linear) homomorphisms of $B$-$A$-bimodule.
\end{proposition}

\begin{definition}
Let $X\in \mod A$ be a right $A$-module and $Y\in \mod A^\op$ be a left $A$-module.
Then define $X\otimes_A Y$ to be the vector space $X\otimes_{\Bbbk} Y/U$ where $U$ is the subspace consisting of $xa\otimes y-x\otimes ay$ for all $x\in X, y\in Y, a\in A$.
\end{definition}

In the above, if ${}_AY_B$ is, in addition, an $A$-$B$-bimodule, then $X\otimes_AY$ has a natural right $B$-module structure: $(x\otimes y)b:=x\otimes(yb)$.  
In fact, as any left $A$-module is also a $A$-$\Bbbk$-bimodule, we can $X\otimes_A Y$ being a $\Bbbk$-vector space as a special case of this observation.

Suppose we have a homomorphism $f:M\to N$ of right $A$-modules.
Then for an $A$-$B$-bimodule ${}_AY_B$ we get a homomorphism of
\begin{align*}
\xymatrix@R=0pt@C=30pt{ M\otimes_A Y_B \ar[r]^{f\otimes_AY}&  N\otimes_AY_B \\
m\otimes y  \ar@{|->}[r] & f(m)\otimes y }
\end{align*}
Note that $(gf)\otimes_AY = (g\otimes_AY)(f\otimes_AY)$, that is, $-\otimes_AY:\Mod A\to \Mod B$ is a (covariant) \defn{functor}\index{functor}.
It is also \defn{$\Bbbk$-linear additive}\index{additive} in the sense that $(\lambda f+\mu g)\otimes_AY=\lambda (f\otimes_AY)+\mu (g\otimes_AY)$ for all homomorphisms $f,g$ and scalar $\lambda,\mu\in\Bbbk$.
Note that this naturally restricts to a functor $-\otimes_AY: \mod A\to \mod B$ on the finitely generated modules so long as $Y$ is so.

Likewise, if $X$ is a bimodule, then ${}_BX\otimes_A Y$ has a left module structure; mutatis mutantis.

%
%\begin{example}
%Consider the bounded quiver
%\[
%Q: \xymatrix@1{ 1 \ar@/^/[r]^{\alpha} & 2 \ar@/^/[l]^{\beta} },  \qquad I=(\beta\alpha).
%\]
%We look at the ($A\otimes A^\op$-module) structure of $Ae\otimes eA, AeA, Ae\otimes_{eAe}eA$ for $e=e_1$ in the following.
%\[
%Ae\otimes eA=\begin{smallmatrix}&& 11&& \\ &21&&12& \\ 11 && 22 && 11 \\ &12&&21&\\&&11&& \end{smallmatrix}
%= \quad \vcenter{
%\xymatrix@dr@=35pt{e_1\otimes e_1 \ar@{-}[r]^{\alpha} \ar@{-}[d]_{\beta^\op} & e_1\otimes \alpha \ar@{-}[r]^{\beta} \ar@{-}[d]_{\beta^\op}& e_1\otimes \alpha\beta \ar@{-}[d]_{\beta^\op} & \\ 
%\beta\otimes e_1 \ar@{-}[r]^{\alpha} \ar@{-}[d]_{\alpha^\op} & \beta\otimes \alpha \ar@{-}[r]^{\beta} \ar@{-}[d]_{\alpha^\op}& \beta\otimes \alpha\beta \ar@{-}[d]_{\alpha^\op} & \\ 
%\alpha\beta\otimes e_1 \ar@{-}[r]^{\alpha}  & \alpha\beta\otimes \alpha \ar@{-}[r]^{\beta} & \alpha\beta\otimes \alpha\beta
%}}
%\]
%As a right $A$-module, this is $\dim_\Bbbk Ae=3$ copies of $eA=P_1=\begin{smallmatrix}
%1\\2\\1\end{smallmatrix}$. 
%\[
%AeA = \begin{smallmatrix}
%&11& \\ 21 && 12 \\ &11&
%\end{smallmatrix}=\quad \vcenter{
%\xymatrix@dr@=30pt{e_1\ar@{-}[r]^{\alpha} \ar@{-}[d]_{\beta^\op} &  \alpha  \ar@{-}[d]^{\beta}\\
%\beta \ar@{-}[r]_{\alpha^\op} &\alpha\beta 
%}
%}
%\]
%As right $A$-module, we have $AeA=eA\oplus \Bbbk \beta \cong P_1\oplus S_2$.
%
%For $Ae\otimes_{eAe} eA$, first note that $eAe=\Bbbk \{e=e_1,\alpha\beta\}$, and so $\alpha\beta\otimes e_1 = e_1\otimes \alpha\beta$.  In particular, so basis elements in $Ae\otimes eA$ vanishes, for example, $\alpha\beta\otimes \alpha =e_1\otimes \alpha\beta\alpha=0$ as $\beta\alpha=0$ in $A$.
%\[
%Ae\otimes_{eAe} eA = \xymatrix@R=8pt@C=3pt{ &11\ar@{-}[ld] \ar@{-}[rd]& \\ 21 \ar@{-}[rrd]\ar@{-}[d] && 12 \ar@{-}[d]\ar@{-}[lld] \\ 22 && 11}\quad =\quad  \xymatrix{
%&e\otimes e\ar@{-}[ld]_{\beta^\op} \ar@{-}[rd]^{\alpha}& \\ \beta\otimes e \ar@{-}[rrd]^(.3){\alpha^\op}\ar@{-}[d]^{\alpha} && e\otimes\alpha \ar@{-}[d]_{\beta}\ar@{-}[lld]_(.3){\beta^\op} \\ \beta\otimes\alpha && e\otimes\alpha\beta
%}
%\]   
%The right $A$-module structure of $Ae\otimes_{eAe}eA$ is isomorphic to $P_1\oplus P_1/\soc(P_1)$.
%\end{example}



\subsection{Hom}

Suppose now that we have ${}_BX_A$ a $B$-$A$-bimodule and $M$ a right $A$-module.
Then the space $\Hom_A(X,Y)$ has a natural \emph{right} $B$-module structure:
\[
(f:X\to Y)\cdot b := (x\mapsto f(bx))
\]
Indeed, we have \[
((f\cdot b)\cdot b')(x) = (f\cdot b)(b'x) = f(bb'x) = (f\cdot(bb'))(x),
\]
and other axioms are even easier to verify.
%$\Hom_A({}_BX_A,-)$ is also a $\Bbbk$-linear additive covariant functor: for $f:M\to N$ a homomorphism of $A$-modules, we have
%\[
%\xymatrix@R=0pt@C=50pt{\Hom_A(X,M) \ar[r]^{f\circ -} & \Hom_A(X,N)\\
%\alpha \ar@{|->}[r] & f\circ \alpha }
%\]

Similarly, in the same setting, $\Hom_A(Y,{}_BX_A)$ also has a \emph{left} $B$-module structure:
\[
(b'\cdot (b\cdot f))(x) = b'((b\cdot f)(x)) = b'( bf(x) ) = (b'b)f(x)=((b'b)\cdot f)(x).
\]  
%However, note that $\Hom_A(f,X)=-\circ f:\Hom_A(N,X)\to \Hom_A(M,X)$ for any $f:M\to N$, i.e. $\Hom_A(-,X)$ is a ($\Bbbk$-linear additive) \defn{contravariantly functor}, meaning that it reverse the direction of homomorphisms.

\begin{exercise}
Show that $\Hom_A({}_BX_A,-)$ defines a $\Bbbk$-linear additive covariant functor from $\Mod A$ to $\Mod B$ (and also $\mod A\to \mod B$ when $X$ is finitely generated).
Likewise, show that $\Hom_A(-,{}_BX_A)$ defines a $\Bbbk$-linear additive contravariant functor on both the (big) module category and the finitely-generated-module category.
\end{exercise}

%The following is just direct application of the categorical notion of co/products.
%\begin{lemma}
%Hom functor commutes with finite direct sum in both variables, i.e. there is a commutative diagram:
%\[
%\xymatrix@C=45pt{ 
%\Hom_A(\bigoplus_{j=1}^\ell L_j,N) \ar[r]^{-\circ \theta}\ar[d]^{\cong} & \Hom_A(\bigoplus_{i=1}^m M_i, N) \ar[d]^{\cong} \\ 
%\bigoplus_{j=1}^\ell\Hom_A(L_j,N) \ar[r]^{(-\circ \theta\iota_j)_{i,j}} & \bigoplus_{i=1}^m \Hom_A(M_i, N)
%}
%\]
%where $\iota_j:L_j\hookrightarrow \bigoplus_k L_k$ and $\pi_i:\bigoplus_{k} M_k\to M_i$ are the canonical maps,
%and there is also a similar commutative diagram arising from $\Hom_A(M,\bigoplus_{i=1}^\ell L_i)\cong \bigoplus_{i=1}^\ell \Hom_A(M,L_i)$.
%\end{lemma}
%\begin{remark}
%In proper functorial language, this is saying that there are natural isomorphisms
%\[
%\Hom_A(\bigoplus_{j=1}^\ell L_j,-)\cong \bigoplus_{j=1}^\ell \Hom_A(L_j,-)\text{ and }\Hom_A(-,\bigoplus_{i=1}^\ell L_i)\cong \bigoplus_{i=1}^\ell(-,L_i).
%\]
%\end{remark}

\begin{lemma}
For any $A$-module $M$, we have (natural) $A$-module isomorphisms 
\[
M\otimes_AA\cong M, \quad\text{and } \Hom_A(A,M)\cong M
\]
given by $m\otimes 1 \mapsto m$ and $f\mapsto f(1)$.  Moreover, $-\otimes_A A$ and $\Hom_A(A,-)$ are both naturally isomorphic to the identity functor.
\end{lemma}
\begin{proof}
First one follows from the construction that $ma\otimes 1 = m\otimes a$.  The second one is just special case of Yoneda lemma.
\end{proof}


\subsection{Tensor-Hom adjunction}
Suppose ${}_AM_B$ is a (f.g.) $A$-$B$-bimodule, then we have two functors:
\[
\xymatrix@1@C=55pt{ \mod A \ar@/^/[r]^{-\otimes_AM} &\mod B \ar@/^/[l]^{\Hom_B(M_B,-)}}.
\]
These are not inverse to each other; but they form a so-called \defn{adjoint pair}\index{adjoint pair}, which is equivalent to saying that there is the following natural isomorphisms.

\begin{theorem}[Tensor-Hom adjunction]
Let $X\in\mod A$, $Y\in\mod B$, ${}_AM_B\in A\mod B$.  Then there is a canonical isomorphism of $\Bbbk$-vector spaces
\[
\theta_{X,M,Y}:\xymatrix@R=0pt@C=0pt{**[l]\Hom_B(X\otimes_AM, Y) \ar[r]^{\cong} & **[r] \Hom_A(X,\Hom_B(M,Y)) \\ 
**[l]{f} \ar@{|->}[r] & **[r] (x\mapsto (m\mapsto f(x\otimes m))) \\
**[l](x\otimes m \mapsto (g(x))(m))& **[r]{g}\ar@{|->}[l] }
\]
that is natural in each of $X,M,Y$.
\end{theorem}
%Here, naturality means that, for example, in the case of $X$, we can (symbolically re)write the isomorphism $\theta_X=\theta_{X,M,Y}:F_{M,Y}(X)\to G_{M,Y}(X)$ for all $X\in \mod B$ so that $F_{M,Y}(f)\circ\theta_X = \theta_{W}\circ G_{M,Y}(f)$ for all $f\in\Hom_B(X,W)$.  
\begin{proof}
Check that the maps written are ($\Bbbk$-linear and) mutual inverse of each other.
\end{proof}

In computer science, the map $\theta_{X,M,Y}$ is also called ``currying".

As innocence as it looks, this isomorphism is fundamental in (homological algebra and) representation theory.

\begin{example}[Adjoint triple (RHS)]
$eA$ is naturally an $eAe$-$A$-bimodule.
Hence, we have an adjoint pair $(-\otimes_{eAe}eA, \Hom_{A}(eA,-))$.

On the other hand, $Ae$ is naturally an $A$-$eAe$-bimodule, and so we have another adjoint pair $(-\otimes_A Ae, \Hom_{eAe}(Ae,-))$.
Note that we have $\Hom_{A}(eA,-)\cong -\otimes_AAe$ by Yoneda lemma.
\end{example}

\begin{example}[Adjoint triple (LHS)]
$A/I$ is naturally an $A$-$A/I$-bimodule for any two-sided ideal $I$ of $A$, and so we have an adjoint pair $(-\otimes_A A/I, \Hom_{A/I}(A/I,-))$.

$A/I$ is also an $A/I$-$A$-bimodule, and so there is another adjoint pair $(-\otimes_{A/I}A/I, \Hom_A(A/I,-))$.
Note that both $\Hom_{A/I}(A/I,-)$ and $\otimes_{A/I}A/I$ sends an $A/I$-module to itself (up to isomorphism) and acts identically on morphisms, i.e. $\Hom_{A/I}(A/I,-)\cong \mathrm{Id} \cong -\otimes_{A/I}A/I$.
\end{example}

\begin{exercise}
Let $F=-\otimes_A M$ and $G=\Hom_B(M,-)$ for an $A$-$B$-bimodule ${}_AM_B$.  Show (without appeal to category theory) that we have natural transformations $\epsilon: FG\to \mathrm{Id}$ and $\eta:\mathrm{Id}\to GF$.  Show also that the following holds:
\[
\id_F = (\epsilon F) \circ (F\eta) \text{ and } \id_G = (G\epsilon)\circ (\eta G).
\]
Here, $\epsilon F$ (and similarly for $\eta G$) denotes the natural transformation $\epsilon F: FGF \to F$ given by $(\epsilon F)_X = \epsilon_{FX}$, whereas $F\eta$ (and similarly for $G\epsilon$) denotes the natural transformation $F\to FGF$ given by $(F\eta)_X = F(\eta_X)$.
\end{exercise}



%
%The following innocent looking isomorphisms are arguably the most used isomorphisms in homological algebra.
%\begin{lemma}\label{lem:adjoint isom K-mod}
%For any finite-dimensional $K$-vector spaces $U,V,W$, the following hold.
%\begin{enumerate}
%\item $V^* \otimes_K W \cong \Hom_K(V,W)$.
%\item $\Hom_K(U\otimes_K V, W)\cong \Hom_K(U, \Hom_K(V,W))$.
%\end{enumerate}
%\end{lemma}
%\begin{proof}
%(1) Let $\CB=\{v_1,\ldots, v_m\},\CC=\{w_1,\ldots, w_n\}$ be bases of $V,W$ respectively.  Let $\CB^*=\{f_1,\ldots, f_m\}$ be the canonical dual basis, i.e. $f_i(v_j)=\delta_{i,j}$ for all $1\leq i,j\leq m$.
%
%Define $\theta(f_i\otimes w_j)$ to be the $K$-linear map that extends $v_k\mapsto f_i(v_k)w_j\in W$ and check that $\theta$ is $K$-linear.
%
%Conversely, for $\alpha\in\Hom_K(V,W)$, let $\phi(\alpha):=\sum_i f_i\otimes \alpha(v_i)$.  Check that $\phi$ and $\theta$ are inverse to each other.
%
%(2) Define
%\[
%\theta : \Hom_K(U\otimes V, W)\to \Hom_K(U, \Hom_K(V,W)), \quad f\mapsto \theta_f,
%\]
%where $\theta_f(u):V\to W$ is the map that sends $v\in V$ to $f(u\otimes v)\in W$.
%
%Define also
%\[\phi : \Hom_K(U, \Hom_K(V,W))\to \Hom_K(U\otimes V, W) , \quad f\mapsto \phi_f,\]
%where $\phi_f(u\otimes v):=(f(u))(v)$.  Check that $\phi$ and $\theta$ are inverse to each other.
%\end{proof}
%\begin{remark}
%The isomorphism (1) absolutely require finite-dimensionality.  The isomorphism (2) is called `currying' in computer science, coined from Curry-Howard correspondence.  This isomorphism is actually natural, and yields an adjoint pair $(-\otimes_K V,\Hom_K(V,-))$ of functors.
%\end{remark}

\pagebreak
%%%%%%%%%%%%%%%%%%%%%%%%%%%%%%%%%%%%%%%%%%%%%%%%%%%%%%%%%%%%%%%%%%%
%%%%%%%%%%%%%%%%%%%%%%%%%%%%%%%%%%%%%%%%%%%%%%%%%%%%%%%%%%%%%%%%%%%
%%%%%%%%%%%%%%%%%%%%%%%%%%%%%%%%%%%%%%%%%%%%%%%%%%%%%%%%%%%%%%%%%%%

\section{Exactness}

\begin{definition}
Consider a sequence\footnote{Superscript/subscript indexing formalism only matters to topologist; I will be liberal in these notations.} $M_\bullet = (M_i, d_i)_{i\in \Z}$ of modules and homomorphisms of modules
\[
M_\bullet: \cdots \xrightarrow{d_{i-2}} M_{i-1}\xrightarrow{d_{i-1}} M_i \xrightarrow{d_{i}} M_{i+1} \xrightarrow{d_{i+1}}\cdots.
\]
We say that the sequence $M_\bullet$ is
\begin{itemize}
\item a (cochain)\footnote{Since we do not deal with any true topological theory, cochain just means the indices increase as we go along the sequence.  We always use cochain convention except perhaps when dealing with projective resolution later.} \defn{complex}\index{complex} if $d_{i+1}d_{i}=0$ for all $i\in\Z$.  In such a case, we have $\Im(d_i)\subset \Ker(d_{i+1})$ for all $i\in \Z$ and the $i$-th \defn{cohomology} of $M_\bullet$ is 
\[
H^i(M_\bullet):=\Ker(d_{i})/\Im(d_{i-1}).
\]

\item \defn{exact}\index{exact} at $M_k$ for some $k\in\Z$ if $\Im(d_{k-1})=\Ker(d_k)$.  Note that this implies $d_k\circ d_{k-1}=0$.

\item \defn{exact} if it is so at every term.

\item \defn{short exact} (often abbreviated as \defn{s.e.s.} or \defn{ses}) if it is a 5-term exact sequence that starts and ends at the trivial module, i.e., of the form
\begin{align}\label{eq:ses}
0\to L\xrightarrow{f} M \xrightarrow{g} N\to 0 
\end{align}
such that $f$ is injective, $g$ is surjective, and $\Ker(g)=\Im(f)$.  In this case, $M$ is also called an \defn{extension}\index{extension} of $N$ by $L$.
\end{itemize}
\end{definition}

\begin{definition}
A (covariant) functor $F:\mod A\to \mod B$ is 
\begin{itemize}
\item \defn{left exact} if it maps short exact sequence (such as \eqref{eq:ses}) to an exact sequence
\[
0\to F(L)\xrightarrow{F(f)} F(M)\xrightarrow{F(g)} F(N).
\]
In other words, it preserves kernel.

\item \defn{right exact} if it maps short exact sequence (such as \eqref{eq:ses}) to an exact sequence
\[
F(L)\xrightarrow{F(f)} F(M)\xrightarrow{F(g)} F(N) \to 0.
\]
In other words, it preserves cokernel.

\item \defn{exact} if it is both left exact and right exact, i.e. maps ses to ses.
\end{itemize}
We define left/right exactness for contravariant functor analogously.  In particular, left exact contravariant fucntor turns cokernel into kernel.
\end{definition}

\begin{lemma}
Let ${}_BX_A$ be an $A$-$B$-bimodule.  Then the following hold.

\begin{enumerate}
\item $\Hom_A(X,-)$ maps an exact sequence $0\to L\xrightarrow{f}M\xrightarrow{g} N$ to an exact sequence 
\[
0\to \Hom_A(X,L)\xrightarrow{f\circ-} \Hom_A(X,M) \xrightarrow{g\circ-} \Hom_A(X,N).
\]
In particular, $\Hom_A(X,-)$ is left exact.

\item $\Hom_A(-,X)$ maps an exact sequence $L\xrightarrow{f}M\xrightarrow{g} N\to 0$ to an exact sequence
\[
0\to \Hom_A(N,X)\xrightarrow{-\circ g} \Hom_A(M,X) \xrightarrow{-\circ f} \Hom_A(L,X).
\]
In particular, the contravariant functor $\Hom_A(-,X)$ is left exact.
\end{enumerate}
\end{lemma}
\begin{proof}
We show (1) and leave (2) for the reader.

\underline{Exactness at $\Hom_A(X,L)$}:  we need $f\circ-$ to be injective.  Indeed, if $f\circ \theta=0$ for some $\theta:X\to L$, then $f(\theta(x))$ for all $x\in X$.  This means that $\theta(x)\in \Ker(f)=0$, and so $\theta=0$.

\underline{$\Im(f\circ-)\subset \Ker(g\circ-)$}:  Suppose that $\theta:X\to M$ is given by $f\circ \phi$ for some $\phi:X\to L$.  Then $g\phi(x)=g(f\phi(x))=(gf)\phi(x)=0$, which means that $\theta\in\Ker(g\circ -)$.

\underline{$\Ker(g\circ-)\subset \Im(f\circ-)$}:  Suppose that $g\theta=0$ for some $\theta:X\to M$.  Then for every $x\in X$, we have $\theta(x)\in \Ker(g)=\Im(f)$, and so we can write $\theta(x)=f(\phi(x))$ for some $\phi(x)\in L$.
Since $f$ is injective, $\phi(x)\in L$ is uniquely determined, and so we have a well-defined function $\phi:X\to L$.
We check that $\phi\in\Hom_A(X,L)$:
\begin{itemize}
\item $f(\phi(x+x'))=\theta(x+x')=\theta(x)+\theta(x')=f(\phi(x))+f(\phi(x'))=f(\phi(x)+\phi'(x))$.  Hence, $f$ being injective implies that $\phi(x+x')=\phi(x)+\phi(x')$.

\item Suppose that $\lambda\in\Bbbk$.  Then $f(\phi(\lambda x))=\theta(\lambda x)=\lambda\theta(x)=\lambda f(\phi(x))=f(\lambda \phi(x))$.  Hence, $f$ being injective implies that $\lambda\phi(x)=\phi(\lambda x)$.
\end{itemize}
Now we have $\theta=f\phi$ as $A$-module homomorphism, and so $\theta\in\Im(f\circ-)$.
\end{proof}

A similar lemma for tensor product exists, and can be proved by direct verification as in the Hom functor case.  Instead, we use another trick involving tensor-Hom adjunction, but first we need one more tool.

\begin{lemma}[Yoneda embedding reflects exactness]\label{lem:Yoneda lemma exactness}
Consider a sequence $L\xrightarrow{f}M\xrightarrow{g} N$ in $\mod A$.
If the sequence
\[
\Hom_A(X,L)\xrightarrow{f\circ-} \Hom_A(X,M) \xrightarrow{g\circ -} \Hom(X,N)
\]
is exact for all $X\in\mod A$, then $L\xrightarrow{f}M\xrightarrow{g} N$ is also exact.
Similarly, if 
\[
\Hom_A(N,X)\xrightarrow{-\circ g} \Hom_A(M,X) \xrightarrow{-\circ f} \Hom(N,X)
\]
is exact for all $X\in\mod A$, then so is the original sequence.
\end{lemma}
\begin{proof}
We show the first one.

\underline{$\Im(f)\subset \Ker(g)$}: Take $X=L$, then we have $gf=(g\circ -)(f\circ -)(\id_L)=0$.

\underline{$\Ker(g)\subset\Im(f)$}:  Consider $X=\Ker(g)$ and inclusion $\iota:\Ker(g)\hookrightarrow M$.
Then $(g\circ-)(\iota)=g\iota = 0$, so exactness implies that $\iota=f\phi$ for some $\phi\in\Hom_A(\Ker(g),M)$.  Hence, $\Ker(g)=\Im(\iota) \subset \Im(f)$.
\end{proof}
%\begin{remark}
%(For people who knows about Yoneda lemma)
%The Yoneda lemma $\Hom_A(eA,M)\cong Me$ is about hom from the representable `functor' $eA$ to a `$\Bbbk$-linear presheaf' $M$; it is really the original Yoneda lemma in a special case.  In contrast, the lemma here is about the Yoneda embedding $X\mapsto \Hom_A(X,-)$ being a left exact functor.
%\end{remark}

\begin{lemma}
$-\otimes_A X$ maps an exact sequence $L\xrightarrow{f}M\xrightarrow{g} N\to 0$ to an exact sequence
\[
L\otimes_AX\xrightarrow{f\otimes_AX} M\otimes_AX \xrightarrow{g\otimes_AX} N\otimes_AX \to 0.
\]
In particular, $-\otimes_A X$ is right exact.
\end{lemma}
\begin{proof}
We apply $\Hom_B(-,Y)$ to the sequence (after tensoring $X$).
By the naturality of the adjoint isomorphism, we have a commutative diagram:
\[
\xymatrix{
0\ar[r]\ar@{=}[d]\ar@{}[rd]|{\circlearrowright} & \Hom_B(N\otimes_A X,Y)\ar[r]^{-\circ g\otimes_AX}\ar[d]|{\cong}\ar@{}[rd]|{\circlearrowright} & \Hom_A(M\otimes_AX,Y) \ar[r]^{-\circ f\otimes_AX}\ar[d]|{\cong} \ar@{}[rd]|{\circlearrowright} & \Hom_A(L\otimes_AX,Y) \ar[d]|{\cong}\\
0\ar[r] & \Hom_A(N,\Hom_B(X,Y))\ar[r]^{-\circ g} & \Hom_A(M,\Hom_B(X,Y)) \ar[r]^{-\circ f} & \Hom_A(L,\Hom_B(X,Y))
}
\]
The second row is exact since it is given by applying the left exact functor $\Hom_A(-,Z)$ for $Z=\Hom_B(X,Y)$.
Hence, (by careful diagram chasing) the first row is also exact.  
Since Yoneda embedding reflects exactness, we get the claimed exactness.
\end{proof}

\pagebreak
%%%%%%%%%%%%%%%%%%%%%%%%%%%%%%%%%%%%%%%%%%%%%%%%%%%%%%%%%%%%%%%%%%%
%%%%%%%%%%%%%%%%%%%%%%%%%%%%%%%%%%%%%%%%%%%%%%%%%%%%%%%%%%%%%%%%%%%
%%%%%%%%%%%%%%%%%%%%%%%%%%%%%%%%%%%%%%%%%%%%%%%%%%%%%%%%%%%%%%%%%%%

\section{Projective and injective modules}

%\begin{definition}
%Consider two ses's $0\to L\xrightarrow{f} M \xrightarrow{g} N \to 0$ and $0\to L'\xrightarrow{f'} M' \xrightarrow{g'} N' \to 0$.
%We say that they are \defn{equivalent} if there is the following commutative diagram:
%\[
%\xymatrix{
%0\ar[r] & L \ar[r]^{f}\ar[d]^{\cong} & M\ar[r]^{g}\ar[d]^{\cong}& N\ar[r]\ar[d]^{\cong}&0 \\
%0\ar[r] & L' \ar[r]^{f'} & M'\ar[r]^{g'}& N'\ar[r]&0 \\
%}
%\]
%\end{definition}


\begin{definition}
An $A$-module $P$ is \defn{projective}\index{projective module}\index{module!projective} if for any given surjective homomorphism $f:M\twoheadrightarrow M'$ and any homomorphism $p:P\rightarrow M$, we have $p$ factors through $f$, i.e. $\exists q:P\rightarrow M'$ s.t. $fq=p$ there is the following commutative diagram
\[
\xymatrix{ & P\ar@{-->}[ld]_{\exists q} \ar[d]^{\forall p}  \\
M' \ar@{->>}[r]^{f} & M.
}
\]
In other words, $f\circ- =\Hom_A(P,f):\Hom_A(P,M')\to \Hom_A(P,M)$ is surjective, i.e. $\Hom_A(P,-)$ is exact.
Denote by $\proj A$ the category of finitely generated projective $A$-modules.

Dually, an $A$-module $I$ is \defn{injective}\index{injective module}\index{module!injective} if for any given injective homomorphism $f:M'\hookrightarrow M$ and any homomorphism $i:M\to I$, $i$ factors through $f$.  This is equivalent to saying that $\Hom_A(f,I):\Hom_A(M,I)\to\Hom_A(M',I)$ is surjective, i.e. $\Hom_A(-,I)$ is exact.
Denote by $\inj A$ the category of finitely generated injective $A$-modules.
\end{definition}

\begin{example}
Take $P=A$.  Then we know that $\Hom_A(A,Y)\cong Y$ via $\alpha\mapsto \alpha(1)$ for any $Y\in \mod A$.  Hence, for any surjective $f:M'\to M$ and any $p:A\to M$, to find $q$ we only need to show that $f(q(1))=p(1)$, but
\[
p(1)=f({}^\exists x)=f({}^\exists q(1)).
\]
That is, the free $A$-module $A_A$ is projective.  Note that this does not require finite-dimensionality of $A$.
Consequently, any free $A$-module (of \emph{any} rank) is also projective.

Dually, using $\Hom_A(X,DA)\cong \Hom_{A^\op}(A,DX)$ and the same argument, we get that $DA$ is injective.  Note that this \emph{DOES} require the finite-dimensionality of $A$ since we need to the isomorphism between the Hom-space under duality.
\end{example}


\begin{lemma}\label{lem:split ses}
The following are equivalent for a ses $0\to L\xrightarrow{f} M \xrightarrow{g} N \to 0$.
\begin{enumerate}
\item There is some $h:N\to N$ such that $gh=\id_N$.
\item There is some $e:M\to L$ such that $ef=\id_M$.
\item There is a commutative diagram
\[
\xymatrix@C=35pt{
0\ar[r] & L \ar[r]^{f}\ar@{=}[d] & M\ar[r]^{g}\ar[d]^{\cong}_{u}& N\ar[r]\ar@{=}[d]&0 \\
0\ar[r] & L \ar[r]_{(1,0)^T} & L\oplus N\ar[r]_{(0,1)}& N\ar[r]&0 \\
}
\]
\end{enumerate}
In the case when any of these conditions is satisfied, we say that the ses \defn{splits}.
\end{lemma}
\begin{proof}
See `Splitting lemma' on Wikipedia.
\end{proof}
\begin{remark}
Note that (3) is strictly stronger than just having $M\cong L\oplus N$ for general modules.  However, in our setting\footnote{also OK for $L,N$ finitely generated over a Noetherian $A$, see \url{https://mathoverflow.net/questions/167701/}}, having $M\cong L\oplus N$ is enough for splitness.  Indeed, applying $\Hom_A(-,L)$ yields an exact sequence 
\[
0\to \Hom_A(N,L)\to\Hom_A(L\oplus N, N)\to \Hom_A(N,N)
\]
of left $\End_A(N)$-modules.
Now the original ses splits is equivalent to having $hf=\id_N$, and so is equivalent to the last map of this induced sequence to be surjective.
Since everything is finite-dimensional in our setting, and $\dim_\Bbbk\Hom_A(L\oplus N,N)=\dim_\Bbbk\Hom_A(L,N)+\dim_\Bbbk\Hom_A(N,N)$, exactness at $\Hom_A(L\oplus N,N)$ means that the last map must be surjective.
\end{remark}

The following justifies why we called $eA$ projective before.

\begin{lemma}\label{lem:projective equiv cond}
The following are equivalent of a finitely genearted $A$-module $P$.
\begin{enumerate}
\item $P$ is projective, i.e. $\Hom_A(P,-)$ is an exact functor.
\item Any ses $0\to L\xrightarrow{f} M\xrightarrow{g} P\to 0$ splits.
%\item Every surjective $A$-module homomorphism $f:M\to P$ splits, i.e. $M=\Ker(f)\oplus P$.
\item $P$ is a direct summand of a free module of finite rank.
\end{enumerate}
\end{lemma}
\begin{proof}
(1) $\Rightarrow$ (2): We have a surjective map $\Hom_A(P,M)\xrightarrow{f\circ-}\Hom_A(P,P)$, and so $\id_P=fq$ for some $q:P\to M$.

(2) $\Rightarrow$ (3): Since $P$ is finitely generated, there is a surjective $A$-module homomorphism $\pi:A^{\oplus n}\to P$ for some $n$.  So we have a short exact sequence \[
0 \to \Ker\pi \to A^{\oplus n} \xrightarrow{\pi} P\to 0.
\]
Hence, it follows by (2) and Lemma \ref{lem:split ses} that $P$ is a direct summand of $A^{\oplus n}$.

(3) $\Rightarrow$ (1):  We have learnt that indecomposable direct summands of $A_A$ is given by the right ideal $eA$ of some primitive idempotent $e=e^2\in A$.  Hence, by the assumption and Krull-Schmidt property $P=\bigoplus_{i=1}^n e_iA$ with $e_i$ primitive idempotents.  Now we have a natural projection $\pi:A^{\oplus n}\to P$ given by sending the $i$-th identity $1_A$ to $e_i$, and a natural inclusion $\iota: P\to A^{\oplus n}$ given by $\iota|_{e_iA}=(e_i A\hookrightarrow A)$.
Note that $\pi\iota =\id_P$.

Consider a surjective $A$-module homomorphism $f:M\to N$ and take any $A$-module homomorphism $p:P\to N$.
This yields $p\pi : A^{\oplus n}\to N$, which can be lifted to some $q':A^{\oplus n}\to M$ as $A^{\oplus n}$ is projective.
%By Yoneda lemma we have $\Hom_A(e_iA,N)\cong Ne_i$ for all $i$, and so $p=(p_i:e_iA\to N)_{i=1,\ldots, n}$ is determined by a tuple $(x_i=p(e_i))$ of elements of $N$. 
%Take $m_i\in M$ so that $f(m_i)=x_i$ and define an $A$-module homomorphism $q=(q_i)_{i=1,\ldots, n}:P\to M$ given by $q(e_i):=m_i$, then we have $fq_i=p_i$ for all $i$, as required.
Now we have
\[
(fq')\iota=(p\pi)\iota = p,
\]
which means that taking $q=q'\iota$ give the required lift of $p$.
\end{proof}
\begin{remark}
This result do not require finiteness anywhere, nor Krull-Schmidt; but this special case yields an easier proof.  For proof of the general case, see Rotman's book Prop 3.3 and Thm 3.5.
\end{remark}

There is a dual result under some restriction.
\begin{lemma}
Suppose $A$ is finite dimensional and $I$ is a finitely generated $A$-module.
Then the following are equivalent.
\begin{enumerate}
\item $I$ is injective, i.e. $\Hom_A(-,I)$ is an exact functor.
\item An ses $0\to I\to M\to N\to 0$ splits.
\item $I$ is a direct summand of finite direct sum of $DA$.
\end{enumerate}
\end{lemma}

%\begin{definition}
%For $M\in\mod A$ and simple module $S\in\mod A$, denote by $[M:S]$ the multiplicity of $S$ as a composition factor of $M$.
%\end{definition}
%
%\begin{lemma}\label{lem:lift simple endo}
%An endomorphism $\theta$ of a simple module $S$ lifts to an endomorphism $\hat{\theta}$ on its projective cover $P$ so that $\theta p=p\hat{\theta}$ for any non-trivial projection $p:P\to S$.
%\end{lemma}
%\begin{proof}
%By Schur's lemma, a non-zero endomorphism $\theta$ of $S$ has an inverse $\phi$.
%As $\phi$ is an isomorphism, it is also surjective, and so the projection $p$ lifts to $q:P\to S$ so that $\phi q=p$.  Now we can lift $q$ to $\hat{\theta}$ so that $p\hat{\theta}=q$.
%Hence, we have a commutative diagram
%\[
%\xymatrix{
%P \ar@{->>}[r]^{p}\ar[d]_{\hat{\theta}} \ar[rd]^{q}& S\ar[d]^{\theta=\phi^{-1}}\\
%P \ar@{->>}[r]^{p}& S}
%\]
%as required.
%\end{proof}


\pagebreak
%%%%%%%%%%%%%%%%%%%%%%%%%%%%%%%%%%%%%%%%%%%%%%%%%%%%%%%%%%%%%%%%%%%
%%%%%%%%%%%%%%%%%%%%%%%%%%%%%%%%%%%%%%%%%%%%%%%%%%%%%%%%%%%%%%%%%%%
%%%%%%%%%%%%%%%%%%%%%%%%%%%%%%%%%%%%%%%%%%%%%%%%%%%%%%%%%%%%%%%%%%%

\section{Projective cover and injective hull}

By definition, any finitely generated module $M$ comes a canonical surjective $A$-module homomorphism $A^{\oplus n}\twoheadrightarrow M$.
One can expect the kernel of this map is `too large', meaning that many direct summands of the domain appear in the kernel.  
For more efficient calculation, we often use the most optimal direct summand of $A^{\oplus n}$.

\begin{definition}
A \defn{projective cover} of an $A$-module $M$ is a projective $A$-module $P$ along with a surjective $A$-module homomorphism $p:P_M\to M$ such that the restriction $p|_Q$ for every proper submodule $Q\subset P_M$ is non-surjective.

Dually, an \defn{injective hull} of $M$ is an injective module $I$ along with an injective $A$-module homomorphism $i:M\to I_M$ such that any proper quotient $q:I_M\twoheadrightarrow J$ yields a non-injective map $qi$.
\end{definition}

\begin{lemma}
The following hold for all $M\in \mod A$.
\begin{enumerate}[(1)]
\item Projective cover $P_M$ of $M$ exists and is unique up to isomorphism.  Moreover, it is characterised by $\top(P_M)\cong \top(M)$.

\item Injective hull $I_M$ of $M$ exists and is unique up to isomorphism.  Moreover, it is characterised by $\soc(I_M)\cong \soc(M)$.
\end{enumerate}
\end{lemma}
\begin{proof}
We show (1); (2) can be shown dually.

Suppose $\top(M)=M/\rad(M)\cong S_1^{\oplus m_1}\oplus\cdots \oplus S_n^{\oplus m_n}$.  By consequence of Artin-Wedderburn, we have $S_i = P_i/\rad P_i$ for each $i$.  Take $P_M=P_1^{\oplus m_1}\oplus \cdots \oplus P_n^{\oplus m_n}$.

Since $M\twoheadrightarrow M/\rad(M)$, the canonical surjection $P_M\twoheadrightarrow M/\rad(M)$ lifts to $p:P_M\to M$. 
As $M\twoheadrightarrow M/\rad(M)$, we have $\Im(p)+\rad(M)=M$, and so it follows from Nakayama lemma (Proposition \ref{prop:rad(M)} (4)) that $\Im(p)=M$, meaning that $p$ is surjective.

Let $Q\subset P_M$ be a submodule; we show that $p|_Q$ is surjective implies $Q\cong P_M$.
Indeed, $p|_Q$ surjective implies that $\top(\Im(p|_Q))=\top(M)$.  
Hence, using the definition of $P_M$ being projective we have a commutative diagram
\[
\xymatrix{ & P_M\ar@{-->}[ld]_{\exists q} \ar[d]^{\overline{p}}  \\
Q \ar@{->>}[r]^{g} & \top(M).
}
\]
Since $Q$ surjects onto $\top(M)$, for $\iota:Q\hookrightarrow P_M$ the canonical inclusion we get that $g\iota q = \top(M)=\top(P)$.
Hence, we have $\Im(\iota q)+\rad(P)= P$.
By Nakayama lemma, we have that $\Im(\iota q)=P$, which means that $\iota$ is also surjective; thus, $\iota$ is an isomorphism, as required.
\end{proof}
\begin{remark}
The claim for projective cover is still true for artinian algebras; but the claim for injective hull really needs finite-dimensionality of $A$.
\end{remark}
%
%\begin{definition}
%For $M\in\mod A$ and simple module $S\in\mod A$, denote by $[M:S]$ the multiplicity of $S$ as a composition factor of $M$.
%\end{definition}
%
%\begin{lemma}
%Let $P_x$ be the (indecomposable) projective cover of simple module $S_x$ whose endomorphism ring is $D_x:=\End_A(S_x)$.
%Then we have
%\[
%\dim_{D_x}\Hom_A(P_x,M)=[M:S_x]=\dim_{D_x}\Hom_A(M,I_x).
%\]
%\end{lemma}
%\begin{proof}
%We show by induction on the length of $M$.
%
%Since $P_x$ has a simple top $S_x$, we have $\Hom_A(P_x,S_y)\cong\Hom_A(S_x,S_y)\cong \delta_{x,y}D_x$.
%Hence, the claim holds for $M=S_y$ a simple module.
%
%In general, we have suppose $S_y$ is a direct summand of the top of $M$, then we have a short exact sequence $0\to N\to M\to S_y\to 0$ where the length of $N$ is strictly less than that of $M$.
%Applying $\Hom_A(P_x,-)$ and using exactness of the Hom-functor we have a short exact sequence
%\[
%0\to \Hom_A(P_x, N)\to \Hom_A(P_x,M) \to \Hom_A(P_x,S_y)\to 0.
%\]
%Note that this is an exact sequence of (right) $\End_A(P_x)$-modules; hence, also an exact sequence of right $D_x=\End_A(S)$-module where $\theta\in D_x$ acts by the lift $\hat{\theta}\in \End_A(P_x)$ as shown in Lemma \ref{lem:lift simple endo}.
%Now, we have $\dim_{D_x}\Hom_A(P_x,M)=\dim_{D_x}\Hom_A(P_x,S_y)+\dim_{D_x}\Hom_A(P_x,N)$, and the proof is completed by applying induction hypothesis.
%
%The proof for the $\dim_{D_x}\Hom_A(M,I_x)$ side is dual.
%\end{proof}



\pagebreak
%%%%%%%%%%%%%%%%%%%%%%%%%%%%%%%%%%%%%%%%%%%%%%%%%%%%%%%%%%%%%%%%%%%
%%%%%%%%%%%%%%%%%%%%%%%%%%%%%%%%%%%%%%%%%%%%%%%%%%%%%%%%%%%%%%%%%%%
%%%%%%%%%%%%%%%%%%%%%%%%%%%%%%%%%%%%%%%%%%%%%%%%%%%%%%%%%%%%%%%%%%%
\section{Resolution and Ext-group}

\begin{definition}
A \defn{projective resolution} $P_\bullet$ of an $A$-module $M$ is a sequence
\[
\cdots \to P_{2} \xrightarrow{d_2} P_1 \xrightarrow{d_1} P_0\xrightarrow{d_0} M\to 0
\]
that is exact everywhere with $P_k$ projective for all $k\geq 0$.
It is \defn{minimal} if $P_k\twoheadrightarrow \Ker(d_{k-1})$ is a projective cover for all $k\geq 1$.
The \defn{$n$-th syzygy} of $M$ is $\Ker(d_n)$ for $(P_\bullet,d_\bullet)$ the minimal projective resolution of $M$.

Dually, an \defn{injective coresolution} $I_\bullet$ of $M$ is a sequence 
\[
0\to M \xrightarrow{d_0} I_0 \xrightarrow{d_1} I_1 \xrightarrow{d_2} I_2 \to \cdots
\]
that is exact everywhere with $I_k$ injective for all $k\geq 0$.
It is \defn{minimal} if $\Cok(d_{k-1})\hookrightarrow I_k$ is an injective hull for all $k\geq 1$.
The \defn{$n$-th cosyzygy} of $M$ is $\Cok(d_{n-1})$ for $(I_\bullet,d_\bullet)$ the minimal projective resolution of $M$.
\end{definition}

%
%\begin{example}
%*** see lecture ***
%\end{example}



\begin{definition}
For $A$-modules $M,N$, let $P_\bullet$ be a projective resolution of $M$.  Define for $k\geq 0$
\begin{align*}
\Ext_A^k(M,N) := & H^k( \Hom_A(P_\bullet,N) ) \\
 =& H^k( \cdots \xleftarrow{-\circ d} \Hom_A(P_{k+1},N) \xleftarrow{-\circ d} \Hom_A(P_k,N) \leftarrow \cdots ) \\ 
= & \frac{\{ f: P_k\to N\mid (fd:P_{k+1}\to N)=0\}}{\{f:P_k\to N\mid f=gd\text{ some }g:P_{k-1}\to N\}}
\end{align*}
Note that $\Ext_A^0(M,N)=\Hom_A(M,N)$.
\end{definition}

Similar to $\Hom_A(-,-)$, $\Ext_A^k(-,-)$ also commutes with finite direct sum in both variables.

There are some other ways to calculate the Ext-groups.

\begin{proposition}
For any $A$-modules $M,N$ and any $k\geq 0$, we have
\[
\Ext_A^k(M,N) = H^k (\Hom_A(M,I_\bullet))
\]
where $I_\bullet$ is an injective coresolution of $N$.
\end{proposition}


\begin{proposition}[Dimension shifting]
For each $k\geq 1$, there are natural isomorphisms 
\[\Ext_A^{k}(\Omega(M),N) \cong \Ext_A^{k+1}(M,N)\cong \Ext_A^k(M,\Omega^{-1}N),
\]
where $\underline{\Hom}_A(X,Y)$ (resp. $\overline{\Hom}_A(X,Y)$) is the quotient of $\Hom_A(X,Y)$ by the subspace consisting of $f:M\to N$ that factors through a projective (resp. injective) $A$-module, i.e. there is a commutative diagram
\[
\xymatrix@R=5pt{X \ar[rr]^{f}\ar[rd] & & Y\\ & Z\ar[ru] & }
\]
for some projective (resp. injective) module $Z$.
\end{proposition}
\begin{proof}
Consider the space $Z_k:=\{f:P_k\to N\mid fd_{k+1}=0\}$ in the definition of $\Ext_A^k(M,N)$.
Since we have a exact sequence $P_{k+1} \to P_k\xrightarrow{p} \Omega^k(M)\to 0$, applying $\Hom_A(-,N)$ yields an exact sequence
\[
0\to \Hom_A(\Omega^k(M), N)\xrightarrow{-\circ p} \Hom_A(P_k,N) \xrightarrow{-\circ d_{k+1}} \Hom_A(P_{k+1},N).
\]
By exactness, we have $Z_k=\Ker(-\circ d_{k+1})\cong \Hom_A(\Omega^k(M),N)$ sending each $f\in Z_k$ to $\overline{f}$ so that $\overline{f}p=f$.

It remains to show that this isomorphisms restricts to one between $B_k:=\Im(-\circ d_k)$ and $\mathcal{P}:=\{f\in\Hom_A(\Omega^k(M),N)\text{ that factors through projective}\}$.
Clearly, any $f\in B_k$ (by definition) factors through a projective $P_{k-1}$ and so $B_k\subset \mathcal{P}$.
For $\overline{f}:\Omega^k(M)\to N$ that factors through a projective, say, $P$, we want $\overline{f}p = gd_k$ some $g$.
Consider $0\to \Omega^k(M) \to P_{k} \to \Omega^{k-1}(M)\to 0$ and apply $\Hom_A(-,N)$ yields
\[
0\to \Hom_A(\Omega^{k-1}(M), N)\xrightarrow{} \Hom_A(P_k,N) \xrightarrow{-\circ d_{k+1}} \Hom_A(P_{k+1},N).
\]
\end{proof}

\begin{exercise}
Suppose that $M,N\in \mod A$.
\begin{enumerate}
\item Show that when $M$ and $N$ are simple modules, then we have $\Hom_A(\Omega(M),N)\cong \Ext_A^1(M,N)\cong \Hom_A(M,\Omega^{-1}(N))$.

\item Show that when every projective $A$-module is injective, then $\Ext^1(M,N)\cong \underline{\Hom}_A(\Omega(M),N)$ where $\underline{\Hom}_A(M,N)$ is the quotient of $\Hom_A(M,N)$ by all homomorphisms factoring through a projective module.
\end{enumerate}
\end{exercise}

\begin{proposition}\label{prop:ext simple}
Consider indecomposable projective modules $P_x, P_y$ with simple tops $S_x, S_y$ respectively.
Then we have an isomorphism of $\Bbbk$-vector spaces $\Ext_A^1(S_x,S_y)\cong \Hom_A(\rad(P_x)/\rad^2(P_x), S_y)$.
Moreover, the $\Bbbk$-dimension of this space is the same as that of $e_x\frac{\rad(A)}{\rad^2(A)}e_y$.
\end{proposition}
\begin{proof}
By the previous exercise, we have \[
\Ext_A^1(S_x,S_y)\cong \Hom_A(\Omega(S_x),S_y)\cong \Hom_A(\rad(P_x),S_y) \cong \Hom_A(\rad(P_x)/\rad^2(P_x), S_y).\]
For the last part, first we have by Schur's lemma 
\[\Hom_A(\rad(P_x)/\rad^2(P_x), S_y)\cong \Hom_A(S_y,\rad(P_x)/\rad^2(P_x))\] as $\Bbbk$-vector space, which then yields
\begin{align*}
\Hom_A(S_y,\rad(P_x)/\rad^2(P_x)) & \cong \Hom_A(P_y,\rad(P_x)/\rad^2(P_x)) \\
&\cong \Hom_A(e_yA, e_x\frac{\rad(A)}{\rad^2(A)}) \cong e_x\frac{\rad(A)}{\rad^2(A)}e_y
\end{align*}
where the last isomorphism uses Yoneda's lemma.
\end{proof}
\begin{remark}
Note that when $A=\Bbbk Q/I$ a bounded path algebra, then arrows from $x$ to $y$ in $Q$ correspond bijectively to basis elements of $\Ext^1_A(S_x,S_y)$.
\end{remark}

\section{Induced long exact sequence}



\begin{definition}
Suppose $C_\bullet=(C_k,d_k)_k, C_\bullet'=(C_k',d_k')_k$ are (chain) complexes of $A$-modules.
A \defn{chain map} is $f_\bullet:C_\bullet\to C_\bullet'$ given by $A$-module homomorphisms $f_k:C_k\to C_k'$ over all $k\in\Z$ such that $d_k'f_k=f_{k+1}d_k$.

Two chain maps $f_\bullet, g_\bullet: C_\bullet\to C_\bullet'$ are \defn{homotopic} (or ``the same up to homotopy") if there exists a sequence of homomorphisms $(h_n:C_{n}\to C_{n-1})_{n\in \Z}$ such that $f_n-g_n = d_{n-1}'h_n+h_{n+1}d_n$ for all $n\in \Z$.
A chain map homotopic to $0$ is said to be \defn{null-homotopic}, or just call it a \defn{null-homotopy}.
\end{definition}

\begin{theorem}[Comparison theorem]\label{comparison thm}
An $A$-module homomorphism $f:M\to N$ extends to a chain map on their projective resolutions, as well as a chain map on their injective coresolutions.
Moreover, changing the choice of (co)resolution results in a homotopic chain map.
\end{theorem}
\begin{proof}
Suppose $P_\bullet, P_\bullet'$ are projective resolutions of $M$ and $N$ respectively.
Define the desired chain map $f_\bullet: (P_\bullet\to M\to 0)\to (P_\bullet'\to N\to 0)$ starting from $f_{-1}=f:M\to N$ inductively as follows.
We take $P_{-1}=M$ and $P_{-1}'=N$.

Given $f_n:P_n\to P_n'$ defined, using the fact that $P_n$ is projective we can lift $f_nd_{n+1}$, which yields a commutative diagram
\[
\xymatrix@C=50pt{ & P_{n+1} \ar@{-->}[ld]_{\exists f_{n+1}} \ar[d]^{f_nd_{n+1}}  \\
P_{n+1}' \ar@{->>}[r]^{d_{n+1}'} & \Im(d_{n+1}'),
}
\]
with the desired chain map property $d_{n+1}'f_{n+1}=f_nd_{n+1}$.

Independence up to homotopy:  EXERCISE.

The claim for injective coresolution can be shown analogously.
\end{proof}

\begin{notation}
For a complex $C_\bullet=(\cdots C_k\xrightarrow{d_k} C_{k+1}\to\cdots)$, and $z_k\in \Ker(d_k)$, denote by $[z_k]:=z_k+\Im(d_{k-1})$.
\end{notation}


\begin{lemma}\label{lem:homology functor}
Suppose $C_\bullet, C_\bullet'$ are complexes of $A$-modules and $f_\bullet:C_\bullet\to C_\bullet'$ is a chain map.
Then for each $k\in \Z$, we have an induced $A$-module homomorphism $H^k(f_\bullet):H^k(C_\bullet)\to H^k(C_\bullet')$ given by $[z_k]\mapsto [f_k(z_k)]$ for any $z_k\in \Ker(d_k:C_k\to C_{k+1})$.
Moreover, $H^k$ preserves identity map and additive, as well as intertwines with composition, i.e. $H^k$ is a functor from the category of complexes of $A$-modules to the category of $A$-modules.
\end{lemma}
\begin{proof}
Since $d_k(z_k)=0$, we have 
\[
d_k'(f_k(z_k)) = f_{k+1}d_k(z_k)=f_{k+1}(0)=0,
\]
i.e. $f_k$ restricts to a map $\Ker(d_k)\to \Ker(d_k')$.

Suppose now that $z_k\in \Im(d_{k-1})$, say, $z_k=d_{k-1}(x_{k-1})$.
Then we have \[
f_k(z_k)=f_kd_{k-1}(x_{k-1})=d_{k-1}'f_{k-1}(x_{k-1}),
\]
i.e. $\Im(f_k|_{\Im(d_{k-1})}) \subset \Im(d_{k-1}')$.
Hence, $H^k(f_\bullet):H^k(C_\bullet)\to H^k(C_\bullet')$ is well-defined.

We leave the rest as exercise.
\end{proof}

The following is a powerful tool in computing Ext-groups.

\begin{theorem}[Induced long exact sequence]
Suppose $0\to X\to Y\to Z\to 0$ is a short exact sequence of $A$-modules.
For any $A$-module $M$, there is the following \defn{long exact sequence}:
\begin{align*}
& 0 \to \Hom_A(M,X) \to \Hom_A(M,Y)\to \Hom_A(M,Z) \to \\
& \phantom{0\to}  \Ext_A^1(M,X) \to \Ext_A^1(M,Y) \to \Ext_A^1(M,Z) \to \\
& \qquad \cdots \to \Ext_A^k(M,X) \to \Ext_A^k(M,Y) \to \Ext^k(M,Z) \to \cdots 
\end{align*}
\end{theorem}
Proof omitted.
%\begin{proof}
%Our first goal is to show that there are \defn{connecting homomorphisms} $\delta_k:\Ext_A^k(M,Z)\to \Ext_A^{k+1}(M,X)$ for all $k\geq 0$.
%
%\underline{Setup}:  Let $P_\bullet$ be a projective resolution of $M$.
%Then for each $k\geq 0$, we have a short exact sequence $0\to \Hom_A(P_k,X)\to \Hom_A(P_k,Y)\to \Hom_A(P_k,Z)\to 0$.
%Denote by $C_N^k:=\Hom_A(P_k,N)$ for $N\in\{X,Y,Z\}$, and $i_k:=\Hom_A(P_k,X\to Y)$ and $p_k:=\Hom_A(P_k,Y\to Z)$.  Then we have the following commutative (check!) grid with exact rows:
%\[
%\xymatrix@R=12pt{
%0 \ar[r] & C^0_X \ar[r]^{i_0}\ar[d] & C^0_Y\ar[r]^{p_0}\ar[d] & C^0_Z \ar[r] \ar[d]& 0 \\
%0 \ar[r] & C^1_X \ar[r]^{i_1}\ar[d] & C^1_Y\ar[r]^{p_1}\ar[d] & C^1_Z \ar[r]\ar[d] & 0 \\
%& \vdots\ar[d] & \vdots \ar[d]& \vdots \ar[d]& \\ 
%0 \ar[r] & C^k_X \ar[r]^{i_k}\ar[d] & C^k_Y\ar[r]^{p_k}\ar[d] & C^k_Z \ar[r]\ar[d] & 0 \\
%& \vdots & \vdots & \vdots& \\ 
%}
%\]
%where every column $C^\bullet_N=(C^k_N, \partial^k_N)_k$ is a complex.  (Be careful that the superscript does not mean taking exponent here.)
%Note that $\Ext_A^k(M,N)=H^k(C^\bullet_N)$ by definition.
%
%\underline{Defining an assignment $\delta^k:\Ker(\partial^k_Z)\to \Ker(\partial^{k+1}_X)$}:
%\begin{itemize}
%\item Start with $z_k\in \Ker(\partial^k_Z)$.  Since $p_k$ is surjective, $z_k=p_k(y_k)$ for some $y_k\in C_k^Y$.
%
%\item By commutativity of the grid, we have $p_{k+1}\partial^k(y_k)=\partial^kp_k(y_k)=\partial^k(z_k)=0$, i.e. $y_{k+1}:=\partial^k(y_k)\in\Ker(p_{k+1})$.
%
%\item By the exactness of the $(k+1)$-st row, we have a unique $x_{k+1}\in C_{k+1}^X$ such that $i_{k+1}(x_{k+1})=y_{k+1}$.
%
%\item  Since $\partial^{k+1}(y_{k+1})=\partial^{k+1}\partial_k(y_k)=0$, it follows from the commutativity of the grid that \[i_{k+2}\partial^{k+1}(x_{k+1})=\partial^{k+1}i_{k+1}(x_{k+1})=\partial^{k+1}(y_{k+1})=0.\]
%Hence, we can now define an assignment $[z_k]\mapsto [x_{k+1}]$.
%\end{itemize}
%
%
%\underline{Showing that $\delta^k:\Ker(p_k)\to H^{k+1}(C^\bullet_X)$ given by $z_k\mapsto [x_{k+1}]$ is a well-defined homomorphism}:
%\begin{itemize}
%\item Recall that we went through $z_k\leadsto y_k\leadsto y_{k+1}\leadsto x_{k+1}$ and that they are given by 
%\[z_k=p_k(y_k)\quad \text{ and }\quad  y_{k+1}=\partial^k(y_k)=i_{k+1}(x_{k+1}).
%\]
%Hence, to show that the assignment is a well-defined map of sets, we need to show that $[x_{k+1}]=[x_{k+1}']$ for $i_{k+1}(x_{k+1}')=\partial^k(y_k')$ for some other $y_k'\in C^k_Y$ with $p_k(y_k')=z_k$.
%
%\item By assumption, we have $p_k(y_k-y_k')=0$.  Hence, we have $y_k-y_k'\in \Ker(p_k)=\Im(i_k)$ by exactness of the $k$-th row at $C^k_Y$.  Thus, we have some $x_k\in C^k_X$ so that $i_k(x_k)=y_k-y_k'$.
%
%\item By commutativity of the grid, we have \[
%i_{k+1}\partial^{k}(x_k)=\partial^ki_k(x_k)=\partial^k(y_k-y_k')=\partial^k(y_k)-\partial^k(y_k')= i_{k+1}(x_{k+1}-x_{k+1}').
%\]
%
%\item As $i_{k+1}$ is injective, we have $x_{k+1}-x_{k+1}'=\partial^k(x_k)$; hence, $[x_{k+1}]=[x_{k+1}']$ as required.
%
%\item To show that the map is actually a homomorphism, we need to check that it preserves zero, scalar multiple, and additivity.  These are all routine check and we leave these as exercise.
%\end{itemize}
%
%\underline{Showing that $\Im(\delta^k|_{\Im(\partial^{k-1}_Z)})=0$, i.e. $\delta^k$ lifts to a homomorphism on $H^k(C^\bullet_Z)$}:
%\begin{itemize}
%\item Suppose $z_k=\partial^{k-1}(z_{{k-1}})$ for some $z_{k-1}\in C^{k-1}_Z$.
%
%\item As $p_{k-1}$ is surjective, we have $z_{k-1}=p_{k-1}(y_{k-1})$ some $y_{k-1}\in C^{k-1}_Y$.
%
%\item By commutativity of the grid, we have $p_k\partial^{k-1}(y_{k-1})=\partial^{k-1}p_{k-1}(y_{k-1})=\partial^{k-1}(z_{k-1})=z_k$.
%
%\item This gives rise to $y_{k+1}=\partial^k(\partial^{k-1}(y_{k-1}))=0$, and so $x_{k+1}=0$.
%\end{itemize} 
%
%\underline{The sequence of the claim is exact everywhere}:
%\begin{itemize}
%\item Combining $\delta_k$ with Lemma \ref{lem:homology functor}, we obtain a sequence
%\[
%0\to H_X^0 \xrightarrow{H^0(i_0)} H_Y^0 \xrightarrow{H^0(p_0)} H_Z^0 \xrightarrow{\delta^0} H_X^1 \to \cdots \xrightarrow{\delta^{k-1}} H_X^{k} \xrightarrow{H^k(i_k)} H_Y^k \xrightarrow{H^k(p_k)} H_Z^k \xrightarrow{\delta^k} \cdots.
%\]
%Note this is exact at $H_X^0$ and $H_Y^0$ by exactness of the Hom-functor.
%For convenience, we write $p^k_*:=H^k(p_k), i^k_*:=H^k(i_k)$.
%
%\item $H^k(p_k)H^k(i_k)=H^k(p_ki_k)=0$ by Lemma \ref{lem:homology functor}.
%
%\item Exactness at $H^k_Y$:  Suppose that $p^k_*([y])=[p_k(y)]=0$.  So we have $p_k(y)=\partial^{k-1}(z)$ for some $z\in C^{k-1}_Z$.  On the other hand, $p_{k-1}$ being surjective yields some $y'\in C^{k-1}_Y$ with $p_{k-1}(y')=z$.  By commutativity of the grid, we have
%\[
%p_k(y)=\partial^{k-1}(p_{k-1}(y'))=p_k\partial^{k-1}(y').
%\]
%Hence, we have $p_k(y-\partial^{k-1}(y'))=0$.
%Exactness at $C^k_Y$ yields $y-\partial^{k-1}(y')=i_k(x)$ for some $x\in C^k_X$.
%By commutativity of the gird, we have 
%\[
%i_{k+1}\partial^k(x)=\partial^ki_k(x)=\partial^k(y)-\partial^k\partial^{k-1}(y')=\partial^k(y)=0.
%\]
%As $i_{k+1}$ is injective, we have $\partial^k(x)=0$, and so $i^k_*([x])=[i_k(x)]=[y-\partial^{k-1}(y')]=[y]$, i.e. $[y]\in \Im(i^k_*)$ as required.
%
%\item $\delta^kp^k_*=0$: $\delta^kp^k_*([y])=\delta^k([p_k(y)])=[x]$ with $i_{k+1}(x)=\partial^k(y')$ for some $y'$ such that $p_k(y')=p_k(y)$.  Since $[x]$ is independent of choice of $y'$, we can just take $y'=y$.  As $y\in \Ker(\partial^{k})$ by assumption, we have $i_{k+1}(x)=0$, and so $x=0$ by injectivity of $i_{k+1}$.
%
%\item Exactness at $H^k_Z$: Suppose that $\delta^k([z])=0$.  Then we have some $x'\in C^k_X$ so that $i_{k+1}\partial^k(x')=\partial^k(y)$ with $p_k(y)=z$.  By commutativity of the grid, we have 
%\[
%\partial^k(y) = i_{k+1}\partial^k(x')=\partial^ki_k(x')
%\]
%and so $y-i_k(x')\in \Ker(\partial^k_Y)$.
%Hence, we have $p^k_*([y-i_k(x')])=[p_k(y)-p_ki_k(x')]=[p_k(y)]=[z]$ as required.
%
%\item $i^k_*\delta^k=0$: $i^k_*\delta^k([z])=[i_k(x)]$ for $[x]=\delta^k([z])$.  By definition of $\delta^k$, we have $i_k(x)=\partial^{k-1}(y)$ with $p_{k-1}(y)=z$.  Hence, we have $[i_k(x)]=[\partial^{k-1}(y)]=0$ by the definition of $H^k_Y$.
%
%\item Exactness at $H^k_X$:  Suppose that $i^k_*([x])=[i_k(x)]=0$.  This means that there is some $y\in C^{k-1}_Y$ with $\partial^{k-1}(y)=i_k(x)$.  By commutativity of the grid and $p_ki_k=0$, we have
%\[
%\partial^{k-1}p_{k-1}(y)=p_k\partial^{k-1}(y)=p_ki_k(x)=0.
%\]
%Hence, $p_{k-1}(y)\in \Ker(\partial^{k-1}_Z)$.  By definition we have $\delta^{k-1}([p_{k-1}(y)])=[x']$ with $i_k(x')=\partial^{k-1}(y)=i_k(x)$.  Hence, $[x']=[x]$ and $[x]\in\Im(\delta^{k-1})$ as required.
%\end{itemize}
%This finishes the proof.
%\end{proof}
%\begin{remark}
%Careful reader should have noticed that there exists a connecting homomorphism long exact sequence by applying homology to any short exact sequence of \emph{complexes}; there we are only consider the case when the complexes involved are all Hom-complexes-of-projective-resolutions.
%\end{remark}


\subsection{Other homological lemmata}



\begin{lemma}[Horseshoe lemma]
Suppose $0\to L\to M\to N\to 0$ is a short exact sequence.  Then a projective resolution $P_\bullet$ of $L$ and a projective resolution $Q_\bullet$ induces a projective resolution of $M$ given by with degree $k\geq 0$ term given by $P_k\oplus Q_k$.
\end{lemma}
In pictorial form:
\[
\xymatrix@C=35pt@R=20pt{ 
& \vdots\ar[d] & \vdots\ar[d] & \vdots\ar[d] & \\
& P^L_1 \ar[d] & P^L_1\oplus P^N_1 \ar[d] & P^N_1 \ar[d] & \\
& P^L_0 \ar[d] & P^L_0\oplus P^N_0 \ar[d] & P^N_0 \ar[d] & \\
0 \ar[r] & L \ar[r] & M \ar[r] & N \ar[r] & 0 \\ 
}
\]

\begin{lemma}
\begin{enumerate}
\item Suppose there are exact rows and homomorphisms $w,u$ such that the left-hand square commutes:
\[
\xymatrix@C=35pt{
 L \ar[r]^{f}\ar[d]^{w}\ar@{}[rd]|{\circlearrowleft} & M\ar[r]^{g}\ar[d]^{u}& N\ar[r]\ar@{-->}[d]^{\exists v}&0 \\
 L' \ar[r]_{f'} &M'\ar[r]_{g'}& N'\ar[r]&0 \\
}
\]
Then there exists a unique homomorphism $v:N\to N'$ such that the right-hand square commutes.  Moreover, if $w,u$ are isomorphisms, then so is $v$.

\item Suppose there are exact rows and homomorphisms $u,v$ such that the right-hand square commutes:
\[
\xymatrix@C=35pt{
 0\ar[r] & L \ar[r]^{f}\ar@{-->}[d]^{\exists w} & M\ar[r]^{g}\ar@{}[rd]|{\circlearrowleft} \ar[d]^{u}& N\ar[d]^{v} \\
 0\ar[r] & L' \ar[r]_{f'} &M'\ar[r]_{g'}& N' \\
}
\]
Then there exists a unique homomorphism $w:L\to L'$ such that the left-hand square commutes.  Moreover, if $u,v$ are isomorphisms, then so is $w$.
\end{enumerate}
\end{lemma}
\begin{proof}
Diagram chasing.
\end{proof}

\begin{lemma}[Short 5-lemma]\label{lem:short 5lem}
Suppose there is a commutative diagram
\[
\xymatrix@C=35pt{
0\ar[r] & L \ar[r]^{f}\ar[d]^{w}\ar@{}[rd]|{\circlearrowleft} & M\ar[r]^{g}\ar[d]^{u}\ar@{}[rd]|{\circlearrowleft}& N\ar[r]\ar[d]^{v}&0 \\
0\ar[r] & L' \ar[r]_{f'} &M'\ar[r]_{g'}& N'\ar[r]&0 \\
}
\]
with exact rows.
Then the following hold.
\begin{itemize}
\item If $w,v$ are both injective, then so is $u$.
\item If $w,v$ are both surjective, then so is $v$.
\end{itemize}
\end{lemma}
\begin{proof}
Diagram chasing.
\end{proof}

\pagebreak
%%%%%%%%%%%%%%%%%%%%%%%%%%%%%%%%%%%%%%%%%%%%%%%%%%%%%%%%%%%%%%%%%%%
%%%%%%%%%%%%%%%%%%%%%%%%%%%%%%%%%%%%%%%%%%%%%%%%%%%%%%%%%%%%%%%%%%%
%%%%%%%%%%%%%%%%%%%%%%%%%%%%%%%%%%%%%%%%%%%%%%%%%%%%%%%%%%%%%%%%%%%

\subsection{Ext-group versus Extensions}

The previous proposition has a better intuition using another manifestation of the Ext-groups.

\begin{definition}
Consider two extensions of $N$ by $L$ given by short exact sequences $0\to L\xrightarrow{f} M\xrightarrow{g} N\to 0$ and $0\to L\xrightarrow{f'} M'\xrightarrow{g'} N\to 0$ are \defn{equivalent} if there is a commutative diagram
\[
\xymatrix@C=35pt{
0\ar[r] & L \ar[r]^{f}\ar@{=}[d] & M\ar[r]^{g}\ar[d]^{u}& N\ar[r]\ar@{=}[d]&0 \\
0\ar[r] & L \ar[r]_{f'} &M'\ar[r]_{g'}& N\ar[r]&0 
}
\]
\end{definition}
\begin{remark}
The map $u$ is necessarily an isomorphism (as a consequence of 5-lemma (Lemma \ref{lem:short 5lem}) or snake lemma).
\end{remark}

%Consider an element $\theta\in \Ext_A^1(N,L)$ and let $P_\bullet$ be the minimal projective resolution of $N$.
%\[
%\xymatrix@R=10pt{ P_1 \ar[rr]\ar@{->>}[rd] \ar[dd]_{\theta} & & P_0 \ar[rr] && N \ar[r] & 0. \\
%& \Omega(N)\ar@{^{(}->}[ru] \ar[ld]^{\theta'} & &&& \\
%L}
%\]
%By definition, $\theta$ is given by some $\hat{\theta}: P_1\to L$ such that $\hat{\theta} d=0$, which means that it induces a well-defined homomorphism from $\overline{\theta}:\Omega(N)\to L$.

\begin{theorem}
There is a bijective correspondence
\[
\left\{
\begin{array}{c}
\text{equivalence classes of}\\
\text{extensions of $N$ by $L$}
\end{array}\right\} \leftrightarrow \Ext_A^1(N,L)
\]
such that the split exact sequence corresponds to $0\in \Ext_A^1(N,L)$.
\end{theorem}
\begin{proof}
Let $\mathcal{E}(N,L)$ be the left-hand set.  Let us first define a map $\phi:\mathcal{E}(N,L)\to \Ext^1_A(N,L)$.
Consider an extension $\xi: 0\to L\xrightarrow{f} M\xrightarrow{g} N\to 0$.
Suppose that $P_\bullet$ is a projective resolution of $N$.
Then, by the same yoga as in the proof of Comparison Theorem (Theorem \ref{comparison thm}), we can lift $\id_N:N\to N$ to a chain map:
\[
\xymatrix@C=35pt{
& P_2\ar[r]^{d_2}\ar@{.>}[d]^0 & P_1 \ar[r]^{d_1}\ar[d]^{\alpha_1} & P_0\ar[r]^{d_0} \ar[d]^{\alpha_0}& N\ar[r]\ar@{=}[d]&0 \\
\xi: &0\ar[r] & L \ar[r]^{f} & M\ar[r]^{g} & N\ar[r]&0
}
\]
In particular, we have $d_2^*(\alpha)=\alpha_1d_2=0$, so we have an element $[\alpha_1]\in \Ext^1(N,L):=\Ker(d_2^*)/\Im(d_1^*)$. 

We claim that $[\xi]\mapsto [\alpha_1]$ is a well-defined map.
\begin{itemize}
\item Changing projective resolution:  By the Comparison Theorem, we get a new chain map $(\alpha_n')_n$ that is homotopic to $(\alpha_n)_n$.  Spelling this out means that $\alpha_1'-\alpha_1=0\cdot h_1+h_0d_1=h_0d_1 = d_1^*(h_0) \in \Im(d_1^*)$.  Hence, $[\alpha_1']=[\alpha_1]$.

\item Changing equivalence extension: By the definition of equivalence betwen extensions, we have a commutative diagram
\[
\xymatrix@C=35pt{
& P_2\ar[r]^{d_2} & P_1 \ar[r]^{d_1}\ar[d]^{\alpha_1} & P_0\ar[r]^{d_0} \ar[d]^{\alpha_0}& N\ar[r]\ar@{=}[d]&0 \\
\xi: &0\ar[r] & L \ar[r]^{f}\ar@{=}[d] & M\ar[r]^{g}\ar[d]^{\beta} & N\ar[r]\ar@{=}[d]&0\\
\xi': &0\ar[r] & L \ar[r]^{f} & M'\ar[r]^{g} & N\ar[r]&0
}
\]
This means that $\xi'$ defines the same element $[\alpha_1]\in \Ext_A^1(N,L)$.
\end{itemize}

Now we show that if $\xi$ is a split extension, then $\phi([\xi])=0$.  Indeed, since $\phi$ is well-defined map to equivalence of extensions, we can just take the canonical split extension $L\oplus N$, then the canonical inclusion $\iota:N\to L\oplus N$ yields the following commutative diagram:
\[
\xymatrix@C=35pt{
& P_2\ar[r]^{d_2}\ar[d]^0 & P_1 \ar[r]^{d_1}\ar[d]^{0} & P_0\ar[r]^{d_0} \ar[d]^{\iota d_0}& N\ar[r]\ar@{=}[d]&0 \\
\xi: &0\ar[r] & L \ar[r]^{(1,0)^t} & L\oplus N \ar[r]^{(0,1)} & N\ar[r]&0
}
\]
Hence, $\phi([\xi])=0$ as required.

Now we need to construct the inverse of $\phi$.
\begin{itemize}
\item Construction of assignment $[\alpha]\mapsto [\xi]$: 
%Without loss of generality, we can assume $P_\bullet$ is taken to be minimal.\footnote{If we work with more general abelian category where existence of projective cover cannot be guarantee, then it is necessary to also argue that construction is independent of the choice of resolution.} Then we have a canonical short exact sequence $0\to \Omega(N)\xrightarrow{i} P_0 \xrightarrow{d_0} N\to 0$ where $P$ is the projective cover of $N$.  
Consider the short exact sequence $0\to \Ker d_0\xrightarrow{i} P_0\xrightarrow{d_0} N\to 0$
Applying $\Hom_A(-,L)$ and using $\Ext_A^1(P,-)=0$ for any projective module $P$, we get an exact sequence
\[
\Hom_A(P,L)\to \Hom_A(\Ker d_0,L)\xrightarrow{\partial} \Ext_A^1(N,L) \to 0
\]
from the induced long exact sequence.
Hence, for each $[\alpha] \in \Ext_A^1(N,L)$, we have some $\overline{\alpha}\in \Hom_A(\Ker d_0,L)$ such that $\partial(\overline{\alpha}) = \alpha$ \footnote{This is, in fact, the homomorphism induced by some $P_1\to L$ by looking at the definition of $\Ext_A^1(N,L)$ as a homology group of the Hom-complex.}
We form the \defn{pushout} $M$ of $i$ and $f$, i.e. 
\[
M:=\Cok\big((i,-\overline{\alpha}):\Ker d_0\to P_0\oplus L\big) = \frac{P_0\oplus L}{S}, \text{ where } S:=\{(i(x), -f(x))\in P_0\oplus L\mid x \in \Omega(N)\},\]
then we get a commutative diagram
\[
\xymatrix@C=35pt{
& 0\ar[r] & \Ker d_0 \ar[r]^{i}\ar[d]^{\overline{\alpha}} & P_0\ar[r]^{d_0} \ar[d]^{\alpha_0}& N\ar[r]\ar@{=}[d]&0 \\
\xi: &0\ar[r] & L \ar[r]^{f} & M\ar[r]^{g} & N\ar[r]&0
}
\]
where $\alpha_0(p):=(0,p)+S$, $f(l):=(l,0)+S$, and $g([p,l]):=j(p)$.
Note that the bottom row is also exact (Exercise: check this!).
We now have an assignment $[\alpha]\mapsto [\xi]$.

\item Well-definedness:  Suppose we choose another $\alpha'\in[\alpha]$ and associated $\overline{\alpha}':\Ker d_0\to L$.  Then following the same procedure we have another extension $0\to L\xrightarrow{f'} M'\xrightarrow{g'} N\to 0$ and homomorphism $\alpha_0':P_0\to M'$.  Define $\theta: M\to M'$ by $(l,p)+S \mapsto f'(l)+\alpha_0'(p) \in M'$ yields a commutative diagram:
\[
\xymatrix@C=35pt{
\xi:& 0\ar[r] & L \ar[r]^{f}\ar@{=}[d] & M\ar[r]^{d_0} \ar[d]^{\theta}& N\ar[r]\ar@{=}[d]&0 \\
\xi': &0\ar[r] & L \ar[r]^{f'} & M'\ar[r]^{g'} & N\ar[r]&0
}
\]
Hence, we have $[\xi]=[\xi']$.

\item Inverse to $\phi$:  This is straightforward to check that $\phi\theta=\id$ by construction; for $\theta\phi=\id$, note that $\alpha_1$ induces naturally $\overline{\alpha_1}:\Ker(d_0)\to L$ which allows us to check it is compatible with $\theta$.
\end{itemize}
\end{proof}

$\Ext_A^1(N,L)$ is an abelian group, so there is a binary operation.  There is indeed a corresponding operation on short exact sequences (which we omit in this text).

\begin{theorem}
The set of equivalence classes of short exact sequence with first term $L$ and last term $N$ form an abelian group under \defn{Baer sum}, and this abelian group is isomorphic to $\Ext_A^1(N,L)$, with the zero element corresponding to the equivalence class of split short exact sequences.
\end{theorem}

There exists similar description for $\Ext_A^n(N,L)$ but the notion of splitness is not as nice as in the case of ses.
In any case, for us, we only need to keep in mind that $\Ext_A^1(N,L)$ contains information about short exact sequence of the form $0\to L\to M\to N\to 0$; c.f. Proposition \ref{prop:ext simple} and relation with arrows of quiver.
Having said that, we should warn that equivalence classes of ses is not the same as isomorphism classes of the middle term, i.e. there exists non-equivalent ses with the same middle term.


\pagebreak
%%%%%%%%%%%%%%%%%%%%%%%%%%%%%%%%%%%%%%%%%%%%%%%%%%%%%%%%%%%%%%%%%%%
%%%%%%%%%%%%%%%%%%%%%%%%%%%%%%%%%%%%%%%%%%%%%%%%%%%%%%%%%%%%%%%%%%%
%%%%%%%%%%%%%%%%%%%%%%%%%%%%%%%%%%%%%%%%%%%%%%%%%%%%%%%%%%%%%%%%%%%


\section{Nakayama algebras}

Recall that a module is \defn{uniserial} if it admits a unique composition series.
We are going to fully classify the indecomposable modules over certain algebras; to achieve this, we determine the possible form of the quiver of such algebra -- this requires use of projective modules.

\begin{lemma}\label{lem:uniserial basic}
Let $M$ be an $A$-module.

\begin{enumerate}
\item If $M$ is uniserial, then so is every submodule and every quotient of $M$.

\item The following are equivalent.
\begin{enumerate}[\rm (i)]
\item $M$ is uniserial.
\item The radical series of $M$ is a composition series.
\item The socle series of $M$ is a composition series.
\end{enumerate}
In particular, the Loewy length of a uniserial module $M$ is the same as the (composition) length of $M$.
\end{enumerate}
\end{lemma}
\begin{proof}
(1) Immediate from definition.

(2) We show (i) $\Leftrightarrow$ (ii); showing (i) $\Leftrightarrow$ (iii) can be done dually.

We prove by induction on the length of $M$.
First note that (i) $\Leftrightarrow$ (ii) holds for all simple modules -- this form the base of the induction.  Now suppose that $M$ has a composition series $0\subset M_1 \subset\cdots \subset M_\ell = M$.

Then using Proposition \ref{prop:rad(M)}, we have
\begin{align*}
M/\rad M \text{ is simple} &\Leftrightarrow M \text{ has a unique maximal submodule} \\ 
& \Leftrightarrow M_{\ell-1} = \rad M
\end{align*}
Now we can apply the induction hypothesis and (1) to finish.
\end{proof}

\begin{lemma}\label{lem:uniserial projective<=>cyclic/linear quiver}
The following are equivalent for a finite-dimensional algebra $A$.
\begin{enumerate}
\item The $A$-module $eA$ is uniserial for all primitive idempotent $e\in A$. 

\item $\rad(eA)/\rad^2(eA)$ is simple or zero for all primitive idempotent $e\in A$.
\end{enumerate}
\end{lemma}
\begin{proof}
\underline{(1) $\Rightarrow$ (2)}:  Since $eA$ is uniserial, the claim follows from Proposition \ref{prop:rad(M)} (2).

\underline{(2) $\Leftarrow$ (1)}:  Fix any primitive idempotent $e$. We show that the radical series of $P:=eA$ is a composition series -- so uniseriality follows from Lemma \ref{lem:uniserial basic}.

Let $r$ be the length of the radical series of $P$, then it suffices to show that $\rad^{i-1}(P)/\rad^i(P)$ is simple or zero for all $i \ge 1$.  We do this by induction on $i$.  We already know from Example \ref{eg:P/radP simple} that $P/\rad P$ is simple, which handles the case when $i=1$.  Note that Lemma \ref{lem:uniserial basic} says that $\rad P/\rad^2 P$ is simple or zero, which handles the case when $i=2$.  

Suppose that $S:=\rad^{i-1}(P)/\rad^i(P)$ is simple.
Let $q:P'\to \rad^{i-1}(P)$ be a projective cover.
Then by Lemma \ref{lem:proj cover characterise}, $q$ induces a projective cover $q':P'\to \rad^{i-1}(P)/\rad^i(P)$.  Hence, $P' = P_S = e'A$ is indecomposable (i.e. $e'$ is an primitive idempotent).

By Proposition \ref{prop:rad(M)}, we have $q|_{\rad P'}$ is surjective onto $\rad^i(P)$, and $q|_{\rad^2(P')}$ is surjective onto $\rad^{i+1}(P)$.
Hence, passing to cokernels, we get a surjective homomorphism
$\rad(P')/\rad^2(P')\twoheadrightarrow \rad^i(P)/\rad^{i+1}(P)$.
Since $P'$ is indecomposable, the assumption implies that $\rad (P')/\rad^2 (P')$ is simple or zero -- hence, so is $\rad^i(P)/\rad^{i+1}(P)$.  This completes the proof.
\end{proof}

\begin{definition}
A finite-dimensional algebra $A$ is a \defn{Nakayama algebra} if every indecomposable direct summand of $A$ and of $DA$ is uniserial.
Equivalently, every indecomposable direct summand of the right $A$-module $A_A$ and of the left $A$-module ${}_AA$ is uniserial.
\end{definition}

\begin{lemma}\label{lem:max uni proj/Naka is inj}
Let $P$ be an indecomposable projective $A$-module.
If $\mathrm{LL}(P)=\mathrm{LL}(A)$, then $P$ is also indecomposable injective.
\end{lemma}
\begin{proof}
Since $P$ is uniserial, we have $S:=\soc(P)$ simple, and so the injective envelop $I_P$ of $P$ is given by $I_S$.  Since $I$ is uniserial, we have $\LL(I)=\mathrm{length}(I)$.  Hence, we have
\[
\LL(P) = \mathrm{length}(P) \le \mathrm{length}(I) = \LL(I) \le \LL(A'),\]
which means that the lengths of $P$ and $I$ are the same, hence $P\cong I$ is injective.
\end{proof}

\begin{theorem}
Let $A$ be a Nakayama algebra.
Then every indecomposable $A$-module is uniserial.  Moreover, every indecomposable module is of the form $P/\rad^t(P)$ for some $t\ge 0$ and some indecomposable projective $A$-module $P$.
\end{theorem}
\begin{proof}
Let $t$ be the Loewy length of $M$.
Since $\rad^t(M)=0$, $M$ has a structure of $A/\rad^t(A)$-module.
Since $\rad^{t-1}(M)\neq 0$, $\rad^{t-1}(A)\neq 0$, which means that the Loewy length of $A/\rad^t(A)$ is $t$.  We claim that $A':=A/\rad A$ is a Nakayama algebra.

Indeed, suppose that $A=\bigoplus_{i=1}^n P_i$ is a decomposition.  Then we have
\[
A' = \frac{A}{\rad^t A} = \bigoplus_{i=1}^n \frac{P_i}{P_i \rad^t A} = \bigoplus_{i=1}^n \frac{P_i}{P_i \rad^t P_i}
\]
which is a decomposition.
Since uniseriality is inherited by quotient modules, each $P_i/\rad^tP_i$ is uniserial.  Now arguing on the left-hand side, we get that every indecomposable left $A'$-module direct summand of $A'$ is also uniserial -- thus showing the claim.

Let $f:\bigoplus_{j=1}^r P_j' \to M$ be a projective of $M$ as $A'$-module, where each $P_j'$ is indecomposable.
Then we have
\[
t = \mathrm{LL}(A') \ge \max\{\mathrm{LL}(P_1'), \ldots, \mathrm{LL}(P_r') \} \ge \mathrm{LL}(M) = t,
\]
where $\mathrm{LL}(X)$ denotes the Loewy length of a module $X$.
Hence, there is some $j$ such that $\mathrm{LL}(P_j') = t$.

Note that there is such a $j$ with $f|_{P_j'}:P_j'\to M$ an injective map -- otherwise, we have $\mathrm{LL}(\Im(f|_{P_j'}))< t$ for all $j$, but
\[
t=\mathrm{LL}(M) = \mathrm{LL}(\Im(f)) = \mathrm{LL}(\bigoplus_{j} \Im(f|_{P_j'})) = \max\{\mathrm{LL}(\Im(f|_{P_j'}))\}_j < t
\]
-- a contradiction.

For $P_j'$ such that $f|_{P_j'}:P_j'\to M$ is injective, we claim that $P_j'$ is an injective $A'$-module.  Indeed, 
since $\mathrm{LL}(A/\rad^t A) = t = \mathrm{LL}(P_j')$, it follows from Lemma \ref{lem:max uni proj/Naka is inj} that $P_j'$ is also indecomposable injective.  Now, we have that $f|_{P_j'}:P_j'\to M$ is a surjective homomorphism from an injective module, and so it has a left inverse.  In particular, we have $P_j'\cong M$, which means that $M\cong P_i/\rad^t(P_i)$ for some $t$.
\end{proof}


\pagebreak
%%%%%%%%%%%%%%%%%%%%%%%%%%%%%%%%%%%%%%%%%%%%%%%%%%%%%%%%%%%%%%%%%%%
%%%%%%%%%%%%%%%%%%%%%%%%%%%%%%%%%%%%%%%%%%%%%%%%%%%%%%%%%%%%%%%%%%%
%%%%%%%%%%%%%%%%%%%%%%%%%%%%%%%%%%%%%%%%%%%%%%%%%%%%%%%%%%%%%%%%%%%


\section{Morita theory}

Roughly speaking, Morita theory is a characterisation of module categories that are equivalent to a given module category.  The beauty of this theory is that, in its strongest form, it only require equivalence between two categories of modules, without \emph{any} extra assumption (namely, linearity, additive, exact) on the functor.
To spell out the main result, we need some terminologies.

\begin{definition}
Let $\CC$ be a category.  An object $X\in \CC$ is called a \defn{generator} if the functor $\Hom_\CC(G,-):\CC\to \mathcal{S}et$ is faithful.
\end{definition}
There is also the dual notion of cogenerator which we will omit.

\begin{example}
Let $R$ be a ring and $\CC=\Mod R$ the (big) module category of $R$.
\begin{enumerate}
\item Since homomorphisms form an abelian group, the definition says that for a homomorphism $f$, we have ``$f\neq 0\Rightarrow \Hom_R(X,f) = f\circ  - \neq 0$".  Clearly, since $\Hom_R(R,-)\cong \mathrm{Id}_{\Mod R}$, the free module $R$ is a generator.
Note that, if $X$ is a generator, since $\Hom_R(R,Y)\cong Y$ is non-zero (whenever $Y$ is so), there is always a homomorphism from $X$ to any non-zero $Y\in \CC$.

\item Suppose that $X$ is a generator of $\Mod R$.  Then for any $Y\in \Mod R$, the module $X\oplus Y$ is also a generator.
\end{enumerate}
\end{example}

\begin{definition}
A \defn{progenerator} of $\Mod R$ (or $\mod R$) is a finitely generated projective generator.
\end{definition}
\begin{example}
Any finite direct sum $R^{\oplus n}$ of $R$ is a progenerator of $\Mod R$.
\end{example}

Our goal is the following.
\begin{theorem}\label{thm:Morita}
\begin{enumerate}[\rm(I)]
\item Suppose that $F:\Mod R\to \Mod S$ is an equivalence of categories.
Then the following hold:
\begin{enumerate}
\item $S\cong \End_R(P)$ for some progenerator $P$ of $\Mod R$.  In particular, $P$ is naturally a $S$-$R$-bimodule.

\item $F \cong \Hom_R(P,-)\cong -\otimes_RQ$ where $Q:=F(R)=\Hom_R(P,R)$ is a progenerator of $\Mod S$.

\item The inverse of $F$ is given by $-\otimes_S P$.

\item $P\otimes_R Q\cong S$ as $S$-$S$-bimodule, and $Q\otimes_S P\cong R$ as $R$-$R$-bimodule.
\end{enumerate}

\item The following are equivalent.
\begin{enumerate}[\rm(1)]
\item $\Mod R\simeq \Mod S$.
\item $\mathsf{Proj}R \simeq \mathsf{Proj} S$, where $\mathsf{Proj}\Lambda$ is the full subcategory of all projective $\Lambda$-modules.
\item $\proj R\simeq \proj S$, where $\proj\Lambda$ is the full subcategory of all finitely generated projective $\Lambda$-modules.
\item $\mod R\simeq \mod S$.
\item $S\cong \End_R(P)$ for some progenerator $P$ of $\Mod R$ (or $\mod R$).
\item $R\cong \End_S(Q)$ for some progenerator $Q$ of $\Mod S$ (or $\mod S$).
\end{enumerate}

\end{enumerate}
\end{theorem}


Let us give more ways to understand the notion of generator before moving to the proof of Morita theory.


\begin{proposition}\label{prop:generator}
The following are equivalent for $X\in \Mod R$.
\begin{enumerate}
\item $X$ is a generator.
\item For all $M\in \Mod R$, there is an indexing set $I$ and a surjective homomorphism $X^{\oplus I} \onto M$.
\item $R$ is a direct summand of a direct sum of copies of $X$.
\end{enumerate}
\end{proposition}
\begin{proof}
\underline{(1)$\Rightarrow$(2)}:  Consider the module $Y$ given by the direct sum of copies of $X$, one for each element of $\Hom_R(X,M)$, i.e.
\[
Y:=X^{\oplus\Hom_R(X,M)}:=\bigoplus_{f\in \Hom_R(X,M)} X_f, \text{ where } X_f:=X.
\]
Then we have an $R$-module homomorphism 
\[
\theta: Y\to M, \text{ given by } (x_f)_f \mapsto \sum_{f\in \Hom_R(X,M)} f(x_f).
\]
Note that since we took the direct sum, meaning that only finitely many $x_f$'s are non-zero, the sum on the right-hand side is well-defined.  

It remains to show that $\theta$ is surjective.  Suppose the contrary.  Then $M/\Im(\theta)\neq 0$ and so the canonical projection $\pi:M\to M/\Im(\theta)$ satisfies $\pi\theta =0$.  Since $X$ is a generator, we must have some $f:X\to M$ such that $\pi \circ f \neq 0$.  By construction, we have $\Im(f)\subset \Im(\theta)$, and so $\Im(\pi f)\subset \Im(\pi\theta)=0$ -- a contradiction.

\underline{(2)$\Rightarrow$(3)}: We have a surjective $f:X^{\oplus I}\to R$, but $R$ is projective as an $R$-module, and so by Lemma \ref{lem:projective equiv cond} the ses $0\to \Ker(f)\to X^{\oplus I} \to R\to 0$ splits, meaning that $R$ is isomorphic to a direct summand of $X^{\oplus I}$.

\underline{(3)$\Rightarrow$(1)}:  This is explained in the example above.
\end{proof}
\begin{remark}\label{rem:generator}
Note that for the category $\mod A$ of finitely generated modules over a finite-dimensional algebra $A$ (in place of $R$), then we can simplifies $Y=X^{\oplus \mathcal{B}}$ for a basis $\mathcal{B}$ of $\Hom_A(X,M)$, i.e. in (2), $I$ can be assumed as finite.  Also, in (3), we can assume the number of copies of direct sum is finite.
\end{remark}

%%%%%%%%%%%%%%%%%%%%%%%%%%%%%%%%%%%%%%
\subsection{Categorical properties}

In the following, we say that a condition (P) is a \defn{categorical property} if the condition (P) is preserved by any category equivalence.

\begin{lemma}\label{lem:cat properties}
%Suppose that $F:\Mod R\to \Mod S$ is an equivalence.
\begin{enumerate}
\item A homomorphism $f:X\to Y$ in $\Mod R$ is injective if and only if for any homomorphism $e,e':Z\to X$, $fe = fe'$ implies $g=g'$.  In particular, being injective is a categorical property.

\item A homomorphism $f:X\to Y$ in $\Mod R$ is surjective if and only if for any homomorphism $g,g':Y\to Z$, $gf=g'f$ implies  $g=g'$.  In particular, being surjective is a categorical property.
%If a homomorphism $g:M\to N$ in $\Mod R$ is surjective, then so is $F(g)$.

\item If $F:\Mod R\xrightarrow{\sim}\Mod S$ is an equivalence, then ses in $\Mod R$ are sent to ses in $\Mod S$.  Moreover, if an ses splits in $\Mod R$, then its image in $\Mod S$ also splits.

\item Being a projective is a categorical property.

\item Being a generator is a categorical property.
\end{enumerate}
\end{lemma}
\begin{remark}
The categorical formulation of injective and surjective can be applied for any category; these are called monomorphisms and epimorphisms.  Note that this depends on the category one works in and may not coincide with the the usual meaning in set theory -- e.g. as an exercise, one can show that the natural injective map $\Z\hookrightarrow \mathbb{Q}$ is an epimorphism in the category of rings (and ring homomorphisms).
\end{remark}
\begin{proof}
We leave (1), (2) as exercise.
%By the previous exercise, we have $fe=fe'$ implies $e=e'$ for any homomorphism $e,e'$.
%Suppose we have homomorphisms $h,h'$ in $\Mod S$.  Since $F$ is an equivalence we can write $h=F(e)$ and $h=F(e')$ for some homomorphisms $e,e'\in \Mod R$.  In particular, we have
%\[
%F(f) h = F(f) h' \Rightarrow F(f)F(e) = F(f)F(e') \Rightarrow F(fe)=F(fe') \Rightarrow fe=fe' \Rightarrow e=e' \Rightarrow F(e)=F(e').
%\]

The first part of (3) follows immediately from (1) and (2).
For the splitting part, recall from Lemma \ref{lem:projective equiv cond} that splitting means the existence of one-sided inverse in either of the middle two morphisms in the sequence. This condition is preserved under category equivalence and so being split is a categorical property.

(4) is evident from the definition of projective.

For (5), suppose that $F:\CC\xrightarrow{\sim} \mathcal{D}$ is an equivalence. 
Since $F$ is essentially surjective, every object in $\mathcal{D}$ is of the form $F(Y)$ with $Y\in \CC$.
Since $\Hom_{\mathcal{D}}(F(X),F(Y))\cong \Hom_\CC(X,Y)$, it then follows from Lemma \ref{lem:Yoneda lemma exactness} that $\Hom_\CC(X,-)$ faithful if and only if so is $\Hom_{\mathcal{D}}(F(X),-)$.
\end{proof}

Recall that a module $M$ is \defn{finitely presented} if it admits a projective presentation
\[
P_1\to P_0 \to M\to 0
\]
such that $P_1, P_0$ are finitely generated (projective).

\begin{lemma}\label{lem:progen is categorical}
Suppose that $F:\Mod R\to \Mod S$ is an equivalence.
\begin{enumerate}
\item For any $M\in \Mod R$, if $M$ is finitely presented, then so is $F(M)$.

\item $F(P)$ is finitely generated projective if so is $P$.

\item $F(P)$ is a progenerator if so is $P$.
\end{enumerate}
\end{lemma}
\begin{proof}
(1): This follows from being finitely presented is equivalent to being \defn{compact} in $\Mod R$, meaning that $\Hom_R(M,)$ commutes with arbitrary coproduct -- in formula: $\Hom_R(M, \bigoplus_{i\in I}X_i) \cong \bigoplus_{i\in I} \Hom_R(M,X_i)$\footnote{Note that $\Hom_R(M, \bigoplus_{i\in I}X_i)\cong \prod_{i\in I}\Hom_R(M,X_i)$ always hold; the point is that being finitely presented means that we only need to look at morphisms to finitely many $X_i$'s.} for any set $I$.  This is a categorical property.\footnote{Being additive is a categorical property; products and coproducts are unique determined by the category itself.} For the proof of this equivalence; see [Zimmermann, Lemma 3.3.13]

(2): For a projective module, being finitely presented is the same as being finitely generated.

(3): Follows from (1) and (2).
\end{proof}


%%%%%%%%%%%%%%%%%%%%%%%%%%%%%%%%%%%%%%
\subsection{Morita bimodules}

Throughout this subsection, we fix $P$ a progenerator of $\Mod R$ over some ring $R$.  Define also
\[
\text{ a ring } S:=\End_R(P) \text{ and } Q:=\Hom_R(P,R).
\]
Note that $P$ is naturally an $S$-$R$-bimodule and $Q$ is naturally a $R$-$S$-bimodule.

\begin{lemma}\label{lem:Morita bimod isom}
The $R$-$R$-bimodule homomorphism $\alpha: Q\otimes_S P\to R$ given by $f\otimes p\mapsto f(p)$ is an isomorphism.
Dually, there is an $S$-$S$-bimodule isomorphism $\beta: P\otimes_R Q \to S$ given by $p\otimes f\mapsto (q \mapsto q\cdot f(p))$.
\end{lemma}
\begin{proof}
We leave as an exercise for the reader to check that $\alpha$ is a bimodule homomorphism.
Since $P$ is a generator of $\mod R$, we have a surjective homomorphism $f:P^{\oplus n} \twoheadrightarrow R$ by Proposition \ref{prop:generator} and Remark \ref{rem:generator}, and so there is some $p=(p_i)_{i=1,\ldots, n}\in P^{\oplus n}$ such that $f(p)=1$.
Since $\Hom_R(P^{\oplus n}, R)\cong \Hom_R(P,R)^{\oplus n}$, we have some $f_i \in \Hom_R(P,R)$ so that $\sum_{i=1}^n f_i(p_i)=1$.
Now $\alpha\big ( \sum_i f_i\otimes p_i \big ) = \sum_{i} f_i(p_i) = 1$.

The claim on $\beta$ is similar.
\end{proof}


\begin{lemma}\label{lem:Hom-tensor swap Morita bimod}
There are natural isomorphisms $\Hom_R(P,-)\cong -\otimes_R Q$ and $\Hom_{S^\op}(Q,-)\cong -\otimes_S P$.
\end{lemma}
\begin{proof}
Take any $M\in \Mod R$.  Similar to the bimodule isomorphism $\beta$, we have a right module homomorphism
\[
\beta_M : M\otimes_R Q \to \Hom_R(P, M) \text{ given by } m\otimes f \mapsto ( p \mapsto m\cdot f(p) ).
\]
One can check in a similar fashion that $\beta_M$ is an isomorphism and is natural in $M$.

The argument for the other natural isomorphism is similar.
\end{proof}

\begin{lemma}
$\Hom_R(P,-):\Mod R\to \Mod S$ is an equivalence.
\end{lemma}


%%%%%%%%%%%%%%%%%%%%%%%%%%%%%%%%%%%%%%
\subsection{Proofs}

\begin{theorem}(Morita Theory I)
Suppose that $F:\Mod R\to \Mod S$ is an equivalence of categories.
Then the following hold:
\begin{enumerate}
\item $S\cong \End_R(P)$ for some progenerator $P$ of $\Mod R$.  In particular, $P$ is naturally a $S$-$R$-bimodule.

\item $F \cong \Hom_R(P,-)\cong -\otimes_RQ$ where $Q:=F(R)=\Hom_R(P,R)$ is a progenerator of $\Mod S$.

\item The inverse of $F$ is given by $-\otimes_S P$.

\item $P\otimes_R Q\cong S$ as $S$-$S$-bimodule, and $Q\otimes_S P\cong R$ as $R$-$R$-bimodule.
\end{enumerate}
\end{theorem}
\begin{proof}
Let $G$ be the inverse of $F$.  Then by Lemma \ref{lem:progen is categorical} $P:=G(S)$ is a progenerator of $\Mod R$.  Since $F,G$ are mutual inverse, $(G,F)$ is an adjoint pair, and we have an natural isomorphism for any $M\in \Mod R$:
\[
F(M) \cong \Hom_{S}(S, F(M)) \cong \Hom_R(G(S), M) = \Hom_R(P,M).
\]
Hence, we have natural isomorphism $F\cong \Hom_R(P,-)$ which also implies that $S\cong \End_R(P)$; this shows (1).

By Lemma \ref{lem:Hom-tensor swap Morita bimod} we have $F\cong -\otimes_R Q$, which shows (2).  (3) then follows by a dual argument.  (4) follows from Lemma \ref{lem:Morita bimod isom}.
\end{proof}

\begin{theorem}
The following are equivalent.
\begin{enumerate}[\rm(1)]
\item $\Mod R\simeq \Mod S$.
\item $\mathsf{Proj}R \simeq \mathsf{Proj} S$, where $\mathsf{Proj}\Lambda$ is the full subcategory of all projective $\Lambda$-modules.
\item $\proj R\simeq \proj S$, where $\proj\Lambda$ is the full subcategory of all finitely generated projective $\Lambda$-modules.
\item $\mod R\simeq \mod S$.
\item $S\cong \End_R(P)$ for some progenerator $P$ of $\Mod R$ (or $\mod R$).
\item $R\cong \End_S(Q)$ for some progenerator $Q$ of $\Mod S$ (or $\mod S$).
\end{enumerate}
\end{theorem}
\begin{proof}
(1) $\Rightarrow$ (2), and (4) $\Rightarrow$ (3):  Being projective is a categorical property.

(2) $\Rightarrow$ (3):  For projective modules, being finitely generated (equivalent to finitely presented) is a categorical property.

(2) $\Rightarrow$ (1), and (3) $\Rightarrow$ (4): Follows from $R$ being a direct summand of (finite) direct sum of any progenerator, which is the image of $S$ under the equivalence.

(3) $\Rightarrow$ (6): Since $R$ is a progenerator, which is a categorical property, the image $P$ of $R$ under the equivalence is also a progenerator of $\Mod R$; the argument to show that $R\cong \End_S(P)$ is just the proof of Morita theory (I).

(1) $\Leftrightarrow$ (5) and (1) $\Leftrightarrow$ (6):  This this Morita theory (I).
\end{proof}


%
%\begin{lemma}
%Suppose we have an adjoint pair $\xymatrix{ F:\Mod A \ar@<.7ex>[r] & \Mod B:G \ar@<.7ex>[l] }$.
%Then we have an $A$-$B$-bimodule $F(A)$ such that there exists  natural isomorphisms of functors
%\[
%F\cong -\otimes_A F(A), \qquad G\cong \Hom_B(F(A),-).
%\]
%\end{lemma}
%\begin{proof}
%Special case of Eilenberg-Watt; see Skowronski-Yamagata Lem 4.4
%\end{proof}
%
%\begin{lemma}
%Equivalence between abelian categories are exact and preserves projectives.
%\end{lemma}
%\begin{proof}
%See Zimmermann Lemma 4.2.1; also Andersen-Fuller Prop 21.2
%\end{proof}

\subsection{Application: Algebras that are Morita equivalent to bound quiver algebras}

Recall the following result we have mentioned.

\begin{proposition}
Suppose $\Bbbk$ is algebraically closed.
Then every finite-dimensional $\Bbbk$-algebra $A$ is Morita equivalent to a bounded path algebra $\Bbbk Q/I$.
More precisely, $\Bbbk Q/I$ is given by $\End_A(\bigoplus_{e} eA)$ where $e$ varies over the set of representative of equivalence classes of primitive idempotents of $A$.
\end{proposition}
\begin{proof}
Suppose that $\{e_1,\ldots, e_n\}$ is the complete list of (representatives of) pairwise orthogonal primitive idempotents. 
For each $1\le i\le n$, let $P_i:=e_iA$ be the corresponding indecomposable projective and $S_i:=P_i/\rad P_i$ be the corresponding simple.  

Take $P:=P_1\oplus \cdots \oplus P_n$, then $P$ is a progenerator of $\Mod A$, and so we have $\mod A\simeq \mod \End_A(P)$.
Since $\Bbbk$ is algebraically closed, $\End_A(S_i) \cong \End_{A/\rad A}(S_i)$ is one-dimensional for every simple $S_i$.
Hence, the algebra $\End_A(P)$ is Morita equivalent to $A$ with only one-dimensional simples.  
It remains to show that there is a bounded path algebra $\Bbbk Q/I$ that is isomorphic to $\End_A(P)$, which we do in the following Theorem \ref{thm:elementary algebra are bdd path alg}.
\end{proof}

\begin{theorem}\label{thm:elementary algebra are bdd path alg}
Suppose that $A$ is a finite-dimensional algebra over an algebraically closed field $\Bbbk$.
If $\dim_{\Bbbk} \End_A(S)=1$ for all simple modules $S$, then there is a bound path algebra $\Bbbk Q/I$ isomorphic to $A$.
\end{theorem}
\begin{proof}
Suppose that $\{e_1,\ldots, e_n\}$ is the complete list of (representatives of) pairwise orthogonal primitive idempotents. 
Define a quiver $Q$ as follows
\begin{enumerate}
\item $Q_0 = \{ 1,\ldots, n\}$.
\item For each $1\le i,j\le n$, take a basis $\mathcal{B}_{i,j}$ of $\Ext_{A}^1(S_i,S_j) \cong e_i (\rad A/\rad^2 A) e_j$ (see Proposition \ref{prop:ext simple}) and define $Q_1$ so that the arrows $i\to j$ correspond to elements of $\mathcal{B}_{i,j}$.
\end{enumerate}
This quiver is called the \defn{ordinary quiver} or \defn{Gabriel quiver} or \defn{$\Ext^1$-quiver} of $A$.  Note that this does not depend on the choice of $\{ e_1, \ldots, e_n\}$ by simple Yoneda lemma type calculation:
\[
e_i\frac{\rad A}{\rad^2 A}e_j \cong \frac{\rad (e_i A)}{\rad^2(e_i A)}e_j \cong \Hom_A\left(e_j A, \frac{\rad (e_i A)}{\rad^2(e_i A)}\right),
\]
and the last space is clearly independent on the choice of primitive idempotents.

We need to construct an isomorphism $\Bbbk Q/I \to A$ for some admissible ideal $I$ of $\Bbbk Q$.
First we define an algebra homomorphism $\phi: \Bbbk Q\to A$ that is surjective, and then check that that $I:=\Ker(\phi)$ is admissible.  

By construction, each arrow $\alpha\in Q_1$ is an element in $\rad A/\rad^2 A$, and so we can lift it to $a_\alpha + \rad^2 A$ for some $a_\alpha\in A$, which gives us the prototype of $\phi$ by
\[
\phi(e_i) = e_i \text{ and }  \phi(\alpha) = a_\alpha.
\]
At the moment, the $\phi$ we have is merely a linear map $\Bbbk Q_0\oplus \Bbbk Q_1 \to A$.
Note that $\Bbbk Q_0\to A/\rad A$ is surjective by the assumption  $\dim_{\Bbbk}\End_A(S)=1$ applied to the Artin-Wedderburn decomposition of $A/\rad A$.
In order to extend this to an algebra homomorphism, we first need to show that $\rad A$ is spanned by all (non-trivial) paths on $Q$.
\begin{lemma}
For any $i,j\in Q_0$, elements of $e_i (\rad A) e_j$ are linear combinations of elements of the form $a_{\alpha_1}a_{\alpha_2}\cdots a_{\alpha_\ell}$ where $\alpha_1\cdots \alpha)_\ell$ is a path from $i$ to $j$ on $Q$.
\end{lemma}
\begin{proof}
We have a vector space decomposition $e_i (\rad A) e_j = e_i \frac{\rad A}{\rad^2 A} e_j \oplus e_i(\rad^2 A)e_j$ which means that any element $a\in e_i (\rad A)e_j$ yields $a - \sum_{\alpha \in \mathcal{B}_{i,j} } \lambda_\alpha a_\alpha \in e_i(\rad^2 A)e_j$ for some $\lambda_\alpha\in \Bbbk$.

By induction, for each $\ell \ge 1$, we obtain a vector space decomposition:
\[
e_i (\rad^{\ell} A) e_j = \left( \sum_{k_1,\ldots, k_{\ell-1}\in Q_0} e_i \frac{\rad A}{\rad^2 A} e_{k_1} \cdot e_{k_1} \frac{\rad A}{\rad^2 A} e_{k_2} \cdots e_{k_\ell} \frac{\rad A}{\rad^2 A} e_j \right) \oplus e_i (\rad^{\ell+1} A) e_j
\]
Note that the elements in the first summand on the right-hand side are can be written as $a_{\alpha_1}a_{\alpha_2}\cdots a_{\alpha_\ell}$ with $\alpha_1\alpha_2\cdots \alpha_\ell$ a path from $i$ to $j$ of length $m$.
Since $A$ is finite-dimensional, $\rad^n A=0$ for some $n$, and the claim follows.
\end{proof}

Now we can extend $\phi$ by \[
\phi(\alpha_1\cdots \alpha_\ell) = \phi(\alpha_1)\cdots \phi(\alpha_\ell) = a_{\alpha_1} \cdots a_{\alpha_\ell}
\]
for all paths $\alpha_1\cdots \alpha_\ell$ on $Q$ and then extend this $\Bbbk$-linearly.
Hence, $\phi$ preserves multiplication on both sides by construction.
Since $1_{\Bbbk Q} = \sum_{i\in Q_0} e_i$, $\phi$ is a unital algebra homomorphism.
By the above lemma, $\{ e_i, a_\alpha \mid i\in Q_0, \alpha\in Q_1\}$ generates $A$ as a $\Bbbk$-algebra, and so $\phi$ is surjective.

It remains to show that $I:=\Ker(\phi)$ is admissible.
By construction, we have $\Im\phi_{\Bbbk Q_{\ge\ell}}\subset \rad^\ell A$ for all $\ell\ge 0$.
Since $\rad A$ is nilpotent, we have $\rad^n A=0$ for some $n\ge 1$.  Note that the case $n=1$ only occurs when $A$ is semisimple, which means that $Q$ contains only vertices and has no arrows; and the claim is automatic.
Hence, we have $n\ge 2$, which means that $\Bbbk Q_{\ge n}\subset I$.

It remains to show that $I\subset \Bbbk Q_{\ge 2}$.
For $a\in I$, similar to the proof of the above lemma, we have
\[
a-\left( \sum_{i\in Q_0} \lambda_i e_i + \sum_{\alpha\in Q_1} \mu_\alpha \alpha \in \Bbbk Q_{\ge 2}\right)
\]
for some $\lambda_i, \mu_\alpha\in \Bbbk$.
To finish the proof we just need to show that all these scalars are zero.
Applying $\phi$ and using $\phi(a)=0$ yields
\[
\sum_{i\in Q_0} \lambda_i e_i + \sum_{\alpha\in Q_1} \mu_\alpha a_\alpha \in \Im \Bbbk Q_{\ge 2} \subset \rad^2 A.
\]
Since $e_i$ are idempotents, we have $\lambda_i=0$ for all $i$.
Hence, we have $\sum_{\alpha} \mu_\alpha a_\alpha \in \rad^2 A$, meaning that
\[
\sum_\alpha \mu_\alpha(a_\alpha + \rad^2 A) =0
\]
but $a_\alpha$'s form the basis of $\rad A/\rad^2A$, and so $\mu_\alpha=0$ for all $\alpha$.  This finishes the proof.
\end{proof}


\pagebreak
%%%%%%%%%%%%%%%%%%%%%%%%%%%%%%%%%%%%%%%%%%%%%%%%%%%%%%%%%%%%%%%%%%%
%%%%%%%%%%%%%%%%%%%%%%%%%%%%%%%%%%%%%%%%%%%%%%%%%%%%%%%%%%%%%%%%%%%
%%%%%%%%%%%%%%%%%%%%%%%%%%%%%%%%%%%%%%%%%%%%%%%%%%%%%%%%%%%%%%%%%%%

\section{Homological dimensions}

\begin{definition}
Let $M$ be an $A$-module.
The \defn{projective dimension} and \defn{injective dimension} of $M$, denoted by $\pdim M$ and $\idim M$ respectively, are the infimum of the length of the projective resolutions and of the injective coresolutions respectively, i.e.
\begin{align*}
\pdim M & := \inf\{d\geq 0\mid 0\to P_d\to \cdots \to P_0\to M\to 0 \text{ is a proj. res.}\}\\
& = d \text{ such that } 0\to P_d\to \cdots \to P_0\to M\to 0 \text{ is the minimal proj. res. of }M;\\
\idim M & := \inf\{d\geq 0\mid 0\to M\to I_0\to \cdots \to I_d\to 0 \text{ is an inj. cores.}\}\\
& = d \text{ such that } 0\to M\to I_0\to \cdots \to I_d\to 0 \text{ is the minimal inj. cores. of }M.
\end{align*}
In case we need to clarify the ring involved, we write $\pdim(M_A), \idim(M_A)$, etc.

The (right) \defn{global dimension} of an algebra $A$ is 
\begin{align*}
\gldim A & := \sup \{ \pdim M \mid M\in \mod A \}
\end{align*}
\end{definition}


\begin{lemma}\label{lem:Ext vanishing pdim}
For $M\in\mod A$, we have
\begin{enumerate}
\item $\pdim M=m \;\Leftrightarrow \; \Ext_A^{>m}(M,-)=0$ and $\Ext_A^m(M,-)\neq 0$.
\item $\idim M=m \;\Leftrightarrow \; \Ext_A^{>m}(-,M)=0$ and $\Ext_A^m(-,M)\neq 0$.
\end{enumerate}
In particular, we have
\begin{align*}
\gldim A := & \sup\{\pdim M\mid M\in \mod A\} \\
 =& \sup\{m\mid \Ext_A^{>m}(M,N)=0 \text{ for } M,N\in\mod A\}\\
 =& \sup\{\idim M\mid M\in \mod A\}\\
 =& \gldim A^\op.
\end{align*}
\end{lemma}
\begin{proof}
(1) and (2) follows from the definition of Ext-groups.  These yields the first three equalities of the last part; the last equality comes from duality $D:\mod A\to \mod A^\op$ which sends projectives to injectives and vice versa.
\end{proof}

\begin{lemma}\label{lem:bdd pdim idim by ses}
For any ses $0\to L\to M\to N\to 0$, we have 
\[
\pdim M\leq \max\{\pdim L,\pdim M\} \text{ and } \idim M\leq \max\{\idim L,\idim N\}.
\]
In particular, we have
\begin{align*}
\gldim A :=& \sup\{\pdim M\mid M\in \mod A\} \\
= &\sup\{\pdim S\mid S\text{ simple $A$-module}\}\\
= &\gldim A^\op\\
= &\sup\{\idim M\mid M\in \mod A\}\\
= &\sup\{\idim S\mid S\text{ simple $A$-module}\}\\
\end{align*}
\end{lemma}
\begin{proof}
Horseshoes lemma, or use long exact sequence.
For the second part, second equality follows from repeatedly applying the first part, and second equality comes from the previous lemma.
\end{proof}

\begin{proposition}
$\gldim A=0$ if and only if $A$ is semisimple.
\end{proposition}
\begin{proof}
Global dimension zero is equivalent to every module - in particular, simple module - being projective (i.e. direct summand of $A_A$).  Hence, every modules are semisimple.
\end{proof}

\begin{proposition}
For an acyclic quiver $Q$, we have $\gldim \Bbbk Q \leq 1$.
Moreover, this is an equality if and only if $Q_1\neq \emptyset$.
\end{proposition}
\begin{proof}
For each simple $\Bbbk Q$-module, we have a short exact sequence
\[
0\to \rad P_x \to P_x \to S_x \to 0,
\]
so it is enough to show that $\rad P_x$ is projective.  In fact, we show that $\rad P_x\cong \bigoplus_{(x\to y)\in Q_1} P_y$, and the inclusion map is given by (left-)multiplication $\alpha\cdot -: P_y \to P_x$ for each $(\alpha:x\to y)\in Q_1$.

By construction, $\top(\rad P_x)=\bigoplus_{(x\to y)\in Q_1} S_y$, and so we have $\bigoplus_{(x\to y)\in Q_1} P_y \twoheadrightarrow \rad P_x$.
Suppose $\sum_p \lambda_p p\in \bigoplus_{(x\to y)\in Q_1} P_y$ such that 
\[
0=\sum_{(\alpha:x\to y)\in Q_1} (\alpha\cdot-)(\sum_p \lambda_p p) = \sum_{(\alpha:x\to y)\in Q_1} \sum_{p:\text{ path with }s(p)=y} \lambda_p (\alpha p).
\]
But all the paths $\alpha p$ appeared are different and thus linear independent in $P_x$; hence, $\lambda_p\equiv 0$ - as required.

Note that if there is an arrow $\alpha \in Q_1$, then the above argument shows that $\pdim S_{s(\alpha)}=1$; thus $\gldim \Bbbk Q=1$.  On the other hand, if there is no arrow in $Q$, then $\Bbbk Q=\Bbbk^{|Q_0|}$ is semisimple.
\end{proof}
\begin{remark}
If $\Bbbk=\overline{\Bbbk}$, then $\Bbbk Q$ are the only $\Bbbk$-algebras, up to Morita equivalence, with global dimension at most one; see any book on quiver representation for proof.
Otherwise, one needs to work with \emph{valued quivers} which is enhancement of quivers with added field extension data (an example is in Homework 1 Exercise 1).
\end{remark}



\pagebreak
%%%%%%%%%%%%%%%%%%%%%%%%%%%%%%%%%%%%%%%%%%%%%%%%%%%%%%%%%%%%%%%%%%%
%%%%%%%%%%%%%%%%%%%%%%%%%%%%%%%%%%%%%%%%%%%%%%%%%%%%%%%%%%%%%%%%%%%
%%%%%%%%%%%%%%%%%%%%%%%%%%%%%%%%%%%%%%%%%%%%%%%%%%%%%%%%%%%%%%%%%%%

\section{Hereditary algebra}

\begin{definition}
An algebra $A$ is (right) \defn{hereditary} if every submodule of a projective (right) $A$-module is projective.
\end{definition}

One can define left hereditary dually (replacing right modules by left; or equivalently, replacing projectives by injectives), but it turns out the two notions are the same - such a property is called `left-right symmetry'.

\begin{theorem}
The following are equivalent for an algebra $A$.
\begin{enumerate}
\item $A$ is right hereditary, i.e. $M\subset P\in\proj A \Rightarrow M\in\proj A$.
\item $A$ is left hereditary, i.e. $\inj A\ni I\twoheadrightarrow M \Rightarrow M\in\inj A$.
\item $\gldim A\leq 1$.
\item $\rad P\in\proj A$ for any $P\in\proj A$.
\item $I/\soc I \in \inj A$ for any $I\in\inj A$.
\end{enumerate}
\end{theorem}
\begin{proof}
(3) $\Rightarrow$ (1): We have a short exact sequence $0\to M\to P\to P/M\to 0$.  Take any $Y\in\mod A$ and consider the long exact sequence induced by applying $\Hom_A(-,Y)$ to this ses.  Then, as $\Ext_A^{>0}(P,Y)=$, we have $\Ext_A^k(M,Y)\cong \Ext_A^{k+1}(P/M,Y)$.
Since $\pdim P/M\leq 1$ by the global dimension assumption, the right-hand space is zero for all $k\geq 0$; hence, $\Ext^{>0}(M,Y)=0$ for all $Y\in\mod A$, i.e. $M$ is projective.

(1) $\Rightarrow$ (3): Let $P$ be the projective cover of an $A$-module $M$.  Then we have a short exact sequence $0\to \Ker(p) \to P\xrightarrow{p} M\to 0$.  It follows from the assumption that $\Ker(p)$ is projective, and thus $\pdim M\leq 1$.

(2) $\Leftrightarrow$ (3):  (Exercise) Dual argument to the case of (1) $\Leftrightarrow$ (3).

(1) $\Rightarrow$ (4): Immediate.

(4) $\Rightarrow$ (1): Induction on $d=\dim_{\Bbbk} P$.  If $d=1$, then there is nothing to show.  Assume $d>1$.  Write $P=P_1\oplus P_2$ with $P_1$ indecomposable and let $p:P\twoheadrightarrow P_1$ be the canonical projection.  Then we have a map $f=pi$ given by the composition of $p$ with the canonical inclusion $i:M\hookrightarrow P$.

If $\Im (pi)=P_1$, i.e. $pi$ is surjective, then $P_1$ being projective implies that $pi$ splits; hence, $M\cong P_1\oplus (M\cap P_2)$.  Since $\dim_{\Bbbk}P_2<d$, by the induction hypothesis we have that the submodule $M\cap P_2$ of the projective module $P_2$ is also projective.  Hence, $M$ is also projective.

If $\Im(pi)\neq P_1$, then $M\subset (\rad P_1)\oplus P_2$.  By the assumption $\rad P_1$ is projective; hence, so is $(\rad P_1)\oplus P_2$.  Since $\dim_\Bbbk (\rad P_1)\oplus P_2 < d$, by the induction hypothesis, we have $M$ projective.

(2) $\Leftrightarrow$ (5): Dual argument to the case of (2) $\Leftrightarrow$ (5).
\end{proof}

\begin{corollary}
Being hereditary is left-right symmetric, i.e. $A$ is hereditary if and only if so is $A^\op$.
\end{corollary}

\begin{corollary}
$\Bbbk Q$ is hereditary for any finite acyclic quiver.
\end{corollary}
\begin{proof}
We have already seen that $\gldim \Bbbk Q\leq 1$.
\end{proof}
\begin{remark}
A better result is that, when $\Bbbk$ is algebraically closed, then hereditary is the same as being Morita equivalent to $\Bbbk Q$.
More generally, being hereditary is the same as being Morita equivalent to Dlab-Ringel's species (roughly, `path algebra' of quiver with added field extensions datum).
\end{remark}





\pagebreak
%%%%%%%%%%%%%%%%%%%%%%%%%%%%%%%%%%%%%%%%%%%%%%%%%%%%%%%%%%%%%%%%%%%
%%%%%%%%%%%%%%%%%%%%%%%%%%%%%%%%%%%%%%%%%%%%%%%%%%%%%%%%%%%%%%%%%%%
%%%%%%%%%%%%%%%%%%%%%%%%%%%%%%%%%%%%%%%%%%%%%%%%%%%%%%%%%%%%%%%%%%%

%\fi
\end{document}
